% ============================================================
%  F. C. Sibbern, Om Erkjendelse og Granskning.
%    Til Indledning i det akademiske Studium.
%  Kjøbenhavn: Paa Fr. Brummers Forlag (Trykt hos C. Græbe), 1822.
%  TRANSCRIPTION (Danish)
%  Source: Google Books scan (Nottingham copy),
%    ~/bibliotek/Sibbern, Frederik/Om_Erkjendelse_og_Granskning.pdf
%  234 PDF pages. Front matter (Fraktur) + Fraktur body (§§ 1--22).
%  PAGE MAP (verified against printed headers + Oversigt):
%    Front matter (roman num.): printed page = PDF page - 6  (Fortale p.III = PDF 9).
%    Body (arabic):             printed page = PDF page - 22 (§.1 opens p.3 = PDF 25).
%    Half-title p.1 = PDF 23, blank verso p.2 = PDF 24, §.1 opens p.3 = PDF 25.
%    (§.2 p.18 = PDF 40, verified; the Oversigt's "§1 S.1" counts the half-title.)
%  See RESUME-NOTES.md for structure, section->page map, and progress.
% ============================================================
\documentclass[12pt,a4paper]{book}
\usepackage[utf8]{inputenc}
\usepackage[T1]{fontenc}
\usepackage{libertinus}
\usepackage{libertinust1math}
\usepackage{textalpha}   % Greek words typed directly (body §§ use Greek: e.g. p.18)
\usepackage[danish]{babel}
\usepackage[protrusion=true,expansion=true]{microtype}
\usepackage{setspace}
\usepackage[margin=1.2in]{geometry}
\usepackage{fancyhdr}
\setlength{\headheight}{28pt}
\addtolength{\topmargin}{-13.5pt}
\usepackage{hyperref}

\hypersetup{
  pdftitle={Om Erkjendelse og Granskning},
  pdfauthor={F. C. Sibbern},
  hidelinks
}

\pagestyle{fancy}
\fancyhf{}
\fancyhead[LE,RO]{\thepage}
\fancyhead[CE]{\textit{\leftmark}}
\fancyhead[CO]{\textit{\rightmark}}

\setstretch{1.3}

\title{\textbf{Om}\\[6pt]
  \Huge\textbf{Erkjendelse}\\[6pt]
  \large og\\[6pt]
  \Large\textbf{Granskning.}\\[12pt]
  \normalsize Til Indledning i det akademiske Studium.\\[16pt]
  \normalsize Af\\[6pt]
  \large Dr.\ Frederik Christian Sibbern,\\[4pt]
  \normalsize Professor i Philosophien.}
\author{}
\date{Kjøbenhavn 1822.\\
  Paa Fr.\ Brummers Forlag.\\
  Trykt hos C.\ Græbe.}

\begin{document}
\maketitle
\thispagestyle{empty}
\cleardoublepage

%% ============================================================
%%  FORTALE  (printed pp. III--XII; PDF 9--18)
%%  Fraktur in the original. Transcribed and image-verified.
%% ============================================================
\chapter*{Fortale.}
\markboth{Fortale}{Fortale}

\begin{center}\rule{0.28\linewidth}{0.4pt}\end{center}

Af Skrivter om Videnskabernes Studium i Almindelighed, om det akademiske
Studium og den Studerendes Kald i Særdeleshed, haves adskillige. Det
nærværende er et nyt Forsøg af denne Art efter en ejendommelig Plan. Kan
det overhovedet ikke andet, end fremme Mennesket i hans Stræben og Virken,
at han tilfulde gjør sig Regnskab for, hvad det er, han gjør og virker, for
at han kan komme til at gjøre det med fuldkommen Klarhed og Besindighed, saa
maa dette vel især være vigtigt for den, hvis hele Stræben netop gaaer ud
paa overalt at eftergranske Tingenes inderste Natur og Væsen. Og kunne der
end, paa den anden Side, vel unægteligen i andre Virkesphærer gives de, som
istedetfor at fremmes, vilde opholdes og standses, ja forvildes i deres
Stræben og Virken, om de i høj Grad vilde tragte efter at trænge ind i deres
eget Selv, for at randsage alt det, der ligger til Grund for deres Virkekraft
og den Naturgave, de have at takke for den, saa kan dette dog ikke gjelde om
den Studerende, som saadan, da han, kaldet til at trænge ind til al
Erkjendelses og alle Tings sidste Grunde, men især at grunde over Aandens og
Aandelivets Væsen, derfra ej kan udelukke den Livets Form, som træder frem i
hans eget Indre.

Man har derfor ogsaa i vor Tid gjort Betragtninger over det akademiske
Studium til Gjenstand for særskilte Foredrag ved Universiteterne.

Men overvejes nærmere, hvad der hører derhen og hvad ikke, saa vise sig
Vanskeligheder, hvilke ogsaa Forfatteren af det nærværende Forsøg ej har
kunnet andet end fornemme. For at oplyse det videnskabelige Studiums Natur og
Væsen maa man gaae ind i Betragtninger, hvilke dog andre videnskabelige
Udviklinger tilegne sig som sig tilhørende, og Spørsmaalet bliver nu,
hvorledes en Grændse drages mellem disse og hiin. Det være tilladt her at
omtale denne Sag med nogle Ord, for at det Forhold, hvori denne Fremstilling
er tænkt at skulle staae til de andre videnskabelige Udviklinger, med hvilke
den paa flere Steder mødes, kan blive oplyst.

For at vise, hvad den, der tragter efter ægte Erkjendelse og Videnskab, bør
at have for Øje, er det det første og vigtigste at betænke, hvad Erkjendelse
og Videnskab ere. Men til Grund for al Erkjenden ligger Tænkning. De
almindelige logiske Udviklinger kunne imidlertid dog ikke høre herhen. Lige
saa lidet kan det, skjønt vistnok det Liv, der skal leve og tage Skikkelse i
den videnskabelige Gransker, er Overvejelsens Gjenstand, høre herhen, at
sætte de forskjellige Livsvirksomheder, der i vor aandelige Leven ere at
adskille, og deres Forhold og Sammenhæng ud fra hinanden, eller at gaae ind i
de psychologiske Undersøgelser. Med den i nyere Tid med Føje som en egen
Disciplin opstillede almindelige Indledning til Philosophien maa en
Indledning til det akademiske Studium ej heller forverles. Imidlertid kan
denne sidste dog ikke andet end at komme i mangfoldig Berøring med alle disse
Undersøgelser. Og saa spørges der da, hvad den dog egentligen skal være og
skal præstere; ja der kan vel opstaae Tvivl om, om den ogsaa kan fordre og
fylde en Plads ved Siden af hine andre videnskabelige Fremstillinger.

Lettest bringe vi os denne Sag paa det klare, naar vi spørge os, hvorledes vi
dog kom til, hvorledes andre før os ere komne til, at gjøre det akademiske
Studium til særegen Gjenstand for særegen Undersøgelse. Det var, fordi det
maatte være den Studerende, der dog som saadan har sit særdeles Kald,
magtpaaliggende at tænke efter over dette sit Kald, over sin Stræben og
Virken. Blant de Former, hvorunder det almene aandelige Liv, hvilket
Psychologien betragter, kan træde frem i en særlig individuel Bestaaen,
hypostaserende sig paa en ejendommelig Maade, er det Liv, der skal leve og
boe i den Studerende, een. Hos andre Individer i andre Virkekredse tager det
Skikkelse paa andre Maader. Hvad der nu hører til fra alle Sider at oplyse
denne Livets Form, og at minde os om, hvad det egentligen er, hvorefter vi
som Studerende stræbe og bør stræbe, og hvad vi, midt i et Liv, hvori saa
meget kan adsprede og forvirre, eller drage hen til Eensidighed, standse og
opholde, stadigen bør at have for Øje, dette hører herhen. Og saaledes viser
sig da et ejendommeligt Middelpunct, hvorom det mangfoldige, der frembyder
sig for Betragtningen, maa samles til Eenhed.

En Fremstilling af denne Art kan betragtes som et Slags Monographie. Af
saadanne for sig fremtrædende videnskabelige Fremstillinger af visse enkelte,
eller til større videnskabelige Systemer henhørende, Gjenstandes hele Natur
og Beskaffenhed, kan man tænke sig to Arter. De sædvanlige naturhistoriske,
og de, der ligne dem, behandle deres Gjenstand ganske saaledes, som den
ellers behandles paa sit Sted i den Lærebygning, hvortil den hører, og ere da
kun at ansee for udhævede, blot maaskee mere end sædvanligen udførte, Afsnit
eller Stykker af det større videnskabelige Hele. Af en anden Beskaffenhed
derimod er en Monographie, som den nærværende. Erkjendelser, der ellers høre
hjemme paa de forskjelligste Steder i andre Videnskaber, skulle her samle sig
om eet Punct, for tilligemed, og i eet med, Betragtningerne over det særegne,
der er ejendommeligt for Gjenstanden, at oplyse denne tilfulde og fra alle
Sider. En saadan Monographie, mere vanskelig og indviklet end en af første
Art, har ogsaa større Adkomst til at træde op som en for sig bestaaende
videnskabelig Betragtning, eller som et mindre Hele i Videnskabernes store
Hele, da der her er et ejendommeligt Middelpunct for en egen videnskabelig
Organisation, hvorved et Indbegreb af sammenhængende Erkjendelser paa en
særegen Maade voxe sammen til Eenhed; og denne Organisation, eller
Fremstillingens indre Grundform, dens Deles indre Forbindelsesmaade og Følge,
eller Udviklingens Gang, bestemmes ved Gjenstandens ejendommelige Grundnatur,
ikke ved en fra en anden Videnskab hentet og ved dennes Gang og Methode
bestemt Norm og Fremstillingsmaade. En Monographie af denne Art forholder sig
til de Videnskaber, til hvis Grundbetragtninger den knytter sig, og til
hvilke meget er at føre tilbage af det, hvad den, skjønt altid paa en ved
Gjenstanden bestemmet ejendommelig Maade og i et eget Lys, fremstiller og
udfører, ligesom et Talent forholder sig til de almindelige
Aandsvirksomheder. Drag bort af Talentet, hvad der hører Forstand og
Judicium, Sands og Phantasie, Fornuft og Følelse til, der vil intet blive
tilbage, som kunde kaldes Talent, og dog er Talentet noget andet, end hine
almenere Aandsvirksomheder. Hvad er Talent da? Den er den ejendommelige
Combination af dem alle i visse bestemte Retninger, den ejendommelige Aandens
Organisation i Henseende til alles forenede Samvirkning, der gaaer ud paa en
vis Art af Productivitet. Ligesom nu et vist Talents grundige psychologiske
Skildring ikke bestaaer i en Skildring af hine almenere Aandsvirksomheder i
Almindelighed, men maa forudsætte dem som bekjendte, medens netop de selv dog
ere det, den overalt maa agte paa og have for Øje, saaledes kan en
Monographie, som den nærværende, have at see hen til mange enkelte Puncter,
som i og for sig tagne høre hen til bestemte Steder i bestemte Videnskaber,
og kan have at fremstille meget, hvilket stykkeviis andre Videnskaber kunne
tilegne sig og drage til sig, og kan dog være, og haandhæve sig som, et
ejendommeligt videnskabeligt Hele for sig.

Indholdet af nærværende Skrivt var det iøvrigt fra først af min Agt ---
Læseren vil maaskee paa flere Steder let bemærke det --- at forelægge
Publicum under en anden Form end den nærværende, under samme Form, hvorunder
jeg adskillige Gange har meddeelt mine Tilhørere det, under Form af
Forelæsninger. Forelæsningen er uden Tvivl hos os den Fremstillingsmaade, der
lettest kunde træde i den Plads, Dialogen har hos Platon; den er, ligesom
denne i hans Tid, ligefrem taget ud af Livet, og almene Begrebers Udvikling
kan neppe paa nogen os naturlig Maade i en sammenhængende Fremstilling bedre
end i Forelæsningens Form forbindes med Talens umiddelbare Henvendelse til
Individualiteter, og faae det livfulde, der maa blive Følgen af en saadan
direct Opmærksomhed paa de tænkende Væsener, der skulle opfordres til med
Selvtænkning at komme det foredragne imøde. Men Frygt for ikke at kunne,
endog kun nogenledes, naae saavidt, at mit Forsøg ikke skulde staae alt for
langt under det Ideal af en saadan Fremstilling, som foresvævede mig, har
bragt mig til at opgive den allerede begyndte Udførelse.

Ogsaa havde jeg fra først af i Sinde, i Anmærkningerne til §§.~13--17, at
omhandle de enkelte Hovedstudier, baade med Hensyn til det philosophiske, saa
og i de øvrige Hovedhenseender. Men hvad jeg derom kunde sige, besluttede jeg
siden at lægge til Side og opbevare til de mundtlige Foredrag, jeg haaber
fremdeles at holde over det akademiske Studium for de Studerende, der endnu
befinde sig i Banens Begyndelse. Jeg torde ej tiltroe mig et saadant
Bekjendtskab med saa mange Videnskaber, at jeg om deres Studium skulde være i
Stand til endog kun i Almindelighed at anstille Betragtninger, der kunde
fortjene at forelægges det Offentlige gjennem Trykken i et for bestandig,
eller dog for en lang Tid, faststaaende Ord.

Endnu maa jeg bemærke, at, naar jeg i nærværende Skrivt er afveget fra min i
de to videnskabelige Arbejder, jeg forhen har leveret, iagttagne Sædvane,
intet Ord at udmærke, da er det skeet, fordi jeg deri har troet at burde give
efter for en Opfordring af en af de blant mine Venner, hvis Omgang jeg især
har meget at takke med Hensyn til min philosopherende Stræben. Men da jeg
ikke er vant til at understrege strax, naar jeg skriver, fordi det næsten
aldrig falder mig ind, at det kunde være fornødent, og jeg denne Gang næsten
havde fuldendt den hele Bog, sinden Beslutningen dertil fremkaldtes hos mig,%
% FLAG p.XII: "sinden" reads thus in Fraktur (long-s); sense wants "inden" (before). Verify.
 saa jeg altsaa først har foretaget Understregelserne ved Skrivtets sidste
Gjennemsigt, men dog det Øjeblik, da der nedskrives, er det rette Øjeblik til
at bestemme, hvad Ord der skal udmærkes, saa tør jeg ikke haabe, at jeg med
bedste Held har udført Beslutningen, og beder derfor Læseren, der dog
alligevel vil finde nok at see igjennem Fingre med, at udstrække sin
Overbærelse ogsaa hertil.

Dette var det, jeg havde at forudskikke det Skrivt, jeg hermed overleverer i
Publicums Hænder.

\begin{flushleft}
Kjøbenhavn d.\ 13de Novb.\ 1821.
\end{flushleft}
\begin{flushright}
Sibbern.
\end{flushright}

\begin{center}\rule{0.28\linewidth}{0.4pt}\end{center}

%% ============================================================
%%  OVERSIGT  (printed pp. XIII--XVI; PDF 19--22)
%%  Synopsis of the whole work, §§ 1--22, each keyed to a page.
%%  Fraktur in the original. Transcribed and image-verified.
%% ============================================================
\chapter*{Oversigt.}
\markboth{Oversigt}{Oversigt}

\begin{center}\rule{0.28\linewidth}{0.4pt}\end{center}

Al Erkjenden er væsentligen en Finden, Modtagen og Skuen. Sandheden, som der
ude,%
% FLAG p.XIII: "som der ude" -- Fraktur reads roughly "deru' der"; uncertain. Verify against scan.
 selv fremstillende sig, refererer sig til en opfattende Individualitet,
medfører for denne en Nødvendighed, der virker Overbeviisning, med hvilken en
personlig Affection er forbunden. §.~1. S.~1.

Men vi erkjende dog kun, idet vi, og forsaavidt vi, selv frembringe. Scimus,
qvia facimus. En levende Personlighed kan ej modtage uden i og under sin egen
Frem-virken. §.~2. S.~18.

Modtagen og Frembringen gaae her i eet. Netop ved selv at uddanne dens Udtryk
eller Fremstillelse i os, modtage vi Sandheden og levendegjøres af den.
§.~3. S.~24.

I den rene Skuen er imidlertid Productiviteten fortabt, og kommer ej til
Syne. Men under en Stræben efter fuldkomment at erkjende gjør Reflectionen,
at den bliver fri og viser sig som saadan. Spørge vi, hvorfor vi dog kun
langsomt gjennem Eftertænkning komme til Sandhed, hvorfor dens Fylde ikke
strax gives, saa er Svaret: fordi et Frihedens Rige skal komme frem, paa det
at et Kjerlighedens Rige kan komme til Virkelighed. §.~4. S.~29.

Den ved Reflectionen fremkaldte Tænken fører nu til en Erkjenden i Kraft af
Grunde. Den højeste Erkjenden vilde vistnok være en over al Gransken og
Grunden ophøjet reen Skuen af den fuldkomne Sandhed, men en saadan umiddelbar
Skuen af det højeste er ikke givet Mennesket. Derimod hviler al hans
Forvisning tilsidst i en umiddelbar Tro. §.~5. S.~36.

Bestræbelsen efter at erkjende i Kraft af Grunde er en Bestræbelse efter at
erkjende Sandheden i dens Sandhed. Denne er det, der fremkalder det egentlige
videnskabelige Studium, som først ret begynder med det akademiske Studium.
§.~6. S.~48.

Erkjendelsen bliver nu en intelligent Erkjenden (modsat den empiriske); en
Erkjendelse af Sandheden i dens Sandhed er nødvendigviis en Erkjenden i Kraft
af det constituerende Første, altsaa en Erkjenden a priori, der fører til en
Construction a priori. §.~7. S.~55.

En saadan Erkjenden gaaer da ikke blot ud paa Grunde for vor Antagen, men paa
Tingenes egne Tilværelsesgrunde. Al Philosophering tragter efter at finde
disse. (I Mathematiken derimod have vi et stort Exempel paa en, endog aldeles
intelligent, Erkjenden, som ej gaaer ud derpaa). §.~8. S.~72.

Tingenes reale Grunde finde vi nærmest i den Forstandighed, som i en
sammenhængende Tingenes Orden fordrer eller medfører dem, tilmed i Tingenes
Begreber. §.~9. S.~88.

Men som sidste Grund viser sig dog et altreelt virkende Første, der
imidlertid, idet det bringer en bestemt Tingenes Orden til Virkelighed,
sætter de enkelte Ting ind i Tilværelsen ifølge en Forstandsmæssighed, tilmed
ifølge Begreb. §.~10. S.~104.

Dette altreelle fremstiller sig for os i Ideen. Erkjendelsen er i sig selv
Ideens egen Sigfremstillen i Bevidstheden, om den end ikke altid viser sig,
eller sees og besiddes af os, som saadan. At Sandhed ej er en blot
Overeensstemmelse imellem en Antagen og dens Gjenstand, men er det, som
levende virker denne Overeensstemmelse, er deraf klart. §.~11. S.~112.

Ideen fremstiller sig for Erkjendelsen og kommer til ydre Leven og Bestaaen i
Videnskaben. §.~12. S.~116.

Det første i al ægte Videnskab er Ideen selv, som Videnskabens Middelpunct og
alt bestemmende Sjæl. I dennes Erkjendelse og alle Erkjendelsers
Tilbageførelse dertil, viser sig det speculative og philosophiske i
Videnskaben. Dog kan Philosophie ogsaa allerede vise sig i det umiddelbare
judiciøse, i en vis videnskabelig Maade at tænke og dømme paa. §.~13. S.~120.

I Speculationens Overgang til det mangfoldige, som derved skal skues i sit
rette Lys, og dettes Bestemmelse, indtræder Dialectiken, deels som en
philosophisk Tilintetgjørelseskonst, deels som en dermed væsentligen og
umiddelbart forbunden videnskabelig Constructions Konst. §.~14. S.~128.

Ægte Erkjendelse er dernæst nødvendigviis og væsentligen systematisk. Dog maa
det umiddelbart frie judiciøse, ligesom det maa berede Systemet Vejen,
saaledes ogsaa ikke fortrænges ved Systemet. §.~15. S.~133.

Endnu er i al Videnskab Detaillets nøjagtige Studium af Vigtighed. §.~16.
S.~146.

Endeligen er i al Videnskab endnu Methoden at bemærke. Denne uddanner sig
bedst af Videnskaben selv og dens historiske Udviklingsgang. Som da Studiet
af en Videnskabs Historie er af væsentlig Vigtighed for dens eget Studium.
§.~17. S.~153.

Ikke blot i Videnskaben, men og gjennem Videnskaben i en personlig
Individualitet, tager Ideen Skikkelse og kommer derved til en fuldkommen ydre
Tilvær eller Bestaaen i den reelle Virkelighed. §.~18. S.~159.

For at komme til denne Bestaaen og Leven i Individet fremkalder den da sin
Modsætning, idet den sætter den individuelle Selvstændighedsfølelse i
Bevægelse, opfordrende til en Er- og Vedkjendelse, som fra Individets Side
skal paafølge. §.~19. S.~160.

Selvtænkningen fremkaldes ved en Spørgen. Mueligheden af at antage, men og
ikke at antage, viser sig da. Benægtelse og Bekræftelse træde imod hinanden.
Men den af nedslaaet Tvivl opstaaende Forvisning, det Ja, der gaaer frem af en
Negations Benægtelse, er af ganske anden Natur, end det, der ej er kommet
gjennem denne Prøve. §.~20. S.~162.

Ideens Kommen til Individet viser sig da som en Forsøgen, og det kommer an
paa, om nu hos Individet det paafølger, der skal paafølge. Heri er nu vel
noget, som staaer i Individets Magt, men det, hvorpaa alt kommer an, staaer
dog ej i Individets Magt, men maa, naar det gives, tages som reen Naade. Den
Modstrid, der her skal bringes til Eenhed, sees især i det theologiske
Studium. Dog i flere Henseender er her en Prøve at gjennemgaae, idet deels
Egoismen let kommer til at trække Næring af Bevidstheden om den Selvtænkning,
hvortil Opfordringen er skeet, deels selve Bestræbelsen efter Indsigt let kan
føre ind i Uvished, Forvirring og Vildfarelse. §.~21. S.~165.

Ej blot i Individer endeligen vil Ideen leve og tage Skikkelse, men i et
Menneskehedens Hele, hvori enhver skal finde sig som Led og practiskt gribe
ind under Udøvelsen af et bestemt Kald. Ved Valget deraf maa Videnskab og Idee
selv lede. Vi maae dernæst tage Videnskaben med ind Livet, og Kaldets
Udøvelse maa tillige være et bestandigen fortsat Studium. §.~22. S.~194.

Men det højeste, det ene fornødne, er, saa høi og betydelig ægte Erkjenden
ogsaa er, dog Erkjendelsen endnu ikke. S.~203.

\begin{center}\rule{0.28\linewidth}{0.4pt}\end{center}

%% ============================================================
%%  BODY -- transcribed in ~10-printed-page batches.
%%  §.1 opens printed p.3 = PDF 25 (half-title p.1 = PDF 23).
%%  Section -> printed-page map (from the Oversigt; §2 verified):
%%    §1 opens p.3  §2 p.18  §3 p.24  §4 p.29  §5 p.36  §6 p.48  §7 p.55
%%    §8 p.72 §9 p.88  §10 p.104 §11 p.112 §12 p.116 §13 p.120
%%    §14 p.128 §15 p.133 §16 p.146 §17 p.153 §18 p.159 §19 p.160
%%    §20 p.162 §21 p.165 §22 p.194 ... body ends ~p.203; then Rettelser (errata).
%%  Emphasis: letterspaced (Sperrung) words -> \emph{}. Greek typed inline.
%% ============================================================

\begin{center}\rule{0.28\linewidth}{0.4pt}\end{center}

\begin{center}\large\S.~1.\end{center}

Al Erkjenden er væsentligen en Finden, en Modtagen og Skuen. (Scimus, qvia
accepimus).

Thi al Erkjenden gaaer ud paa det Sande og i og ved den skulle vi komme til
Sandheden, komme til at have og besidde denne. Men Sandhed er ikke noget, som
er afhængigt af os; den er og bliver, uden alt Hensyn til os, som den i sig
selv er. Vort kan det altsaa ikkun være at søge ganske ligefrem at komme til
Bevidsthed om den, som den selv er og fremstiller sig for os; men en saadan
Kommen til Bevidsthed er en Modtagen og, naar den er bestemt og klar og tillige
umiddelbar, eller vis og indlysende ved sig selv allene, en Skuen. Men dette
umiddelbare, ved sig selv visse, er i al Erkjenden paa ethvert Punct, endog om
vi komme til den paa en middelbar Maade; thi Erkjendelsen beroer da dog altid
paa Indsigten i, at af saadanne foregaaende Erkjendelser den nærværende følger,
hvilken Følgen eller Consequents ligefrem maae indlyse ved sig selv, eller maae
skues, hvormed da tillige selve det, hvad der saaledes følger, viser sig som
det, der, under Forudsætning af de foregaaende Erkjendelsers Sandhed, indlyser
af sig selv, men kun under samme Forudsætning erkjendes det og. Saa at selv den
relative og middelbare Erkjenden, paa det Punct, hvor den ifølge sin
Middelbarhed først indtræder, gaaer frem af en ligefrem Indseen og forsaavidt
selv er en Indseen og Skuen. Den er en sig efterhaanden udviklende Skuen. En
Modtagen er den i alle Tilfælde.

Som saadan er den altid og væsentligen forbunden med en Nødvendighed til at
forestille saaledes, som der forestilles, og ikke anderledes. Den umiddelbare
Fornemmelse af denne Nødvendighed er Overbeviisningen, der kun kan opstaae
derved, at vi umiddelbart gribes af den og føle den. Erkjendelsens Gjenstand
paanøder sig os, og det ganske saaledes, som den nu erkjendes at være.

Det Sande, der erkjendes, viser sig herved som det, der udøver en bestemmende
og virkende Indflydelse hos os, i hvilken det selv udtaler og fremstiller sig,
og derved gjenlever, i os.

\begin{center}Anmærkning og Tillæg.\end{center}

1.~At al Erkjenden er en Modtagen og Skuen, lærer ogsaa et Blik paa
Erkjendelsens forskjellige Kredse snart. Ikke blot, hvad der strax falder i
Øinene, at \emph{Begyndelsen} til al Erkjendelse, og det ligesaavel i de højere
Sphærer, som i den laveste sandselige, er en ligefrem Vaerbliven og Opfatten,
en Bemærken og Sandsen. Ogsaa i sin Fremgang til den højeste Fuldendthed bliver
den ved at være en umiddelbar Tagen og Skuen. I den lavere Sandsnings Kreds er
dette strax ved første Øiekast aldeles øjensynligt; thi om end her vist nok en
vis sammenfattende og combinerende Virksomhed strax maae træde til den
opfattende, saa at Sandsningen ikkun kommer i Stand formedelst en eftergjørende
Construction, saa gaaer dog ogsaa denne ikkun ud paa at fremstille Tingen ganske
som den forefindes og selv fremstiller sig for os. Vi føle os afficerede, og
ganske som Affectionen fører det med sig, bliver Sandsningen, netop saaledes,
som vi bestemmes dertil ved denne Affection, maae vi forestille. Og indtræder
her vel og ikke sjeldent en Forskjel imellem det, der synes at være, og det,
der befindes i sig selv at være, saa befindes dog ogsaa dette kun, og maae netop
erkjendes, som det ifølge en Randsagelse befindes, om ej af Sandsen, dog af den
overvejende Eftertænkning.

Som det her viser sig i Sandsernes og den derved fremkaldte Erfaringens Sphære,
saaledes forholder det sig og med enhver højere Forstandens eller Fornuftens
Erkjenden, ja er just ved den højeste Erkjenden især i Øine faldende. Hvad der
saa erkjendes, uden at Erfaring paaberaabes, ifølge en af reen Tænkning
udspringende Indsigt, være det i Kraft af mathematisk Demonstration eller
apriorisk Reflection eller philosophisk Speculation, saa erkjendes det altid som
noget, der indsees at maatte saa være, og en vor hele Forestillen bestemmende
Forstands- eller Fornuftnødvendighed indtræder derved; som denne tilsiger det at
være, maae det erkjendes at være. Men naar saaledes her al Eftertænkning, al
Sammenstillen, Forbinden og Slutning gaaer ud paa at udforske eller udgrunde
det, der nødvendigviis maae antages, som det nu selv vil vise sig at være, saa
at det kun erkjendes for det Sande, idet det erkjendes for det eneste
antagelige, saa gaaer jo Eftertænkningen ud paa at lade sig det \emph{give} ved
at lade sig det paanøde ved en Tænkningens Nødvendighed. Det erkjendes kun, idet
det indsees, det er: sees og skues (med det højere Aandens Øie), som det selv
er. Triangelen construerer jeg, men at dens Vinkler tilsammen ere lige to rette,
indseer jeg som nødvendigt ifølge Triangelens egen, mig nu eengang givne, Natur;
ja allerede idet jeg construerer den, gjør jeg det paa en vis bestemt Maade og
med visse Betingelser, hvilken Constructionens Maade og hvilke Betingelser jeg
maae tage, som de paanøde sig mig. Alleslags menneskelige Sammentræf og Forhold
kan jeg udtænke mig, (dog at jeg ogsaa allerede heri maae gaae frem efter visse
forefundne Betingelser og Bestemmelser) men hvad der i et givet Forhold er ret
eller uret, godt eller slet, det maae Grundideen om det Gode og Rette lære mig,
saa at jeg ved denne maae lade mig lede i min Dom og i ethvert Tilfælde lade
blive ret og godt, hvad jeg ifølge hiin Idee seer at være ret og godt. Mit kan
det ikkun være at eftergranske det Rettes og Godes Væsen, og nu at erkjende det,
som det selv fremstiller sig for mit aandelige Øje. Saa at, om ogsaa den
Erkjenden, der gaaer ud over den blotte Erfaring, fremkaldes ved Spørgsmaal,
der, som indre Ansporelser til at søge mere end det erfarne, vise, at vi ej lade
os nøje med det i Erfaringen allene givne, saa fører dog den hele derved i
Bevægelse kommende Tænkning ikkun til en ny Modtagen og Opfatten af højere
Orden, under hvilken vi igjen maae søge ligefrem at iagttage og skue, hvad der
fremstiller sig for os.

Men vi blive da nu ogsaa i den Modtagen og Skuen, hvori al Erkjenden bestaaer,
en \emph{Hovedforskjel} vaer. Hvad vi formedelst Sandserne opfatte og komme til
Kundskab om, ja ogsaa hvad den ideelle eller intellectuelle Vaerbliven og Skuen,
som allerede i denne Sphære indtræder ifølge Tænkning, kan lære os, naar de for
Sandserne sig aabenbarende Tings indre Grundformer, deres Forhold og deres
Virksomheder erkjendes ganske saaledes, som de ligefrem forefindes, alt dette
modtages blot, som noget der gives os udenfra i den i Tid og Rum sig
fremstillende Virkelighed. Denne Erkjenden er den blot \emph{empiriske}; den
Nødvendighed, ifølge hvilken den indtræder og ved hvilken den faaer Charakteren
af en Erkjenden, er kun en Nødvendighed af den Art, der gaaer ud fra en vis til
Tid og Rum bundet \emph{Affection}, en aldeles sandselig Affection af lavere
Natur. Hvad der paa denne Maade erkjendes, erkjendes endnu kun som
\emph{tilfældigt}. Det er os givet som det virkelige, og for os er der en
Nødvendighed til at tage det saaledes, men at det selv nødvendigen skulde være,
som det er, sees endnu ej i denne blot empiriske Erkjenden.

Af en ganske anden Art er derimod den mathematiske og philosophiske Erkjenden;
heri er en intellectuel Skuen, der ingenlunde har hiin Tilfældigheds Præg. Hvad
der her erkjendes, tages ej blot som givet, saa at det dermed skulde være nok,
men opfattes og forestilles som noget, der er saa, fordi det maae saa være
ifølge en indre \emph{Forstands- og Fornuftsnødvendighed} af ganske andet
Udspring, end den, hvorved visse sandselige Affectioner paanøde sig os. Vel er
det os givet, men givet indenfra, eller dog givet som noget, der stemmer sammen
med en indre Fordring, og aldeles ikke gives os formedelst sandselige
Affectioner. Det erkjendes som noget, der ej kan være anderledes; det erkjendes
som \emph{væsentligt}, og derfor gyldigt for alle Gange og overalt, hvorfor der
ogsaa kun i denne Sphære kan tales om evige Sandheder. Vi ville kalde denne
Erkjenden den \emph{intelligente}, baade fordi dette Ord synes mere betegnende,
end Ordet intellectuel, saa og fordi dette sidste tillige kan omfatte den
nysnævnte blot empiriske, men dog ideelle Erkjenden, hvorved Tingenes indre
Grundformer og Forhold aldeles opfattes saaledes, som de ligefrem forefindes i
Erfaringen.

2.~Nødvendigheden indtræder, idet Sandheden viser sig i Forhold til (refererer
sig til) \emph{den Individualitet}, der optager den, og i Overbeviisningens
Følelse træder der det Sandes Leven og Virken i Individet noget reent
individuelt imøde,%
% FLAG p.10: "imøde" (Fraktur m/n ambiguous, "inøde" as printed) -- read by sense. Verify.
 hvori en Person lader sig til Syne og træder frem som saadan, idet den ganske
knytter sig til, og giver sig hen til, hiin Virken. I den rene Modtagen og Skuen
saaes endnu kun det Sande som fremstillende sig selv, saa at kun dets egen
(universelle) Fremtræden deri viiste sig, ingen personlig Leven. Modtagelsen var
blot som en Afspejling eller Gjenspejling, liig den, hvormed Vandet gjengiver
Lyset og Farven, ikke en personlig Optagelse, liig den, hvorved Ædelstenen i sin
Farve har en Lysets Kraft og Virkning optaget i sig til Eenhed med sin egen
Natur og Tilvær. Men i Overbeviisningen bliver den til en inderlig Tilegnelse,
hvorved det sig i og for en Individualitet fremstillende Sande nu ogsaa kommer
til kraftigen at boe i denne Individualitet, som i et Væsen, i hvilket det har
taget Sæde, og hvilket nu selv er og bliver noget ved Optagelsen. I denne det
Sandes og Individets Sammenknytning til Eenhed er ogsaa den Nødvendighed, som
vistnok er deri, saaledes fortabt, at den ikke særligen fornemmes eller kommer
til Syne. Her er en \emph{udeelt}, over Nødvendighedens Følelse opløftet, fri
Leven i det Sande. Kun naar det individuelle, rivende sig løs fra det almene,
bevæger sig i en særlig personlig Stræben, kommer Modsætningen og tilligemed
samme Nødvendigheden for Lyset.

3.~I Overbeviisningen træder en personlig \emph{Affection} og Sindsbevægelse
πάθος, til Modtagelsen, med en personlig Tilknytning og Hengivenhed til
Sandhedens Magt, saa at Individet ej blot lader Sandheden faae Rum og Magt hos
sig, men og gjør det med Lyst, følende sig deri at eje noget og sin personlige
Leven derved udvidet og forhøjet. Hengivenheden viser sig nu som
\emph{Kjerlighed}. Denne gjør, at med det Sandes Magt ogsaa Personligheden og
Friheden træder fuldkomment frem; thi i Kjerligheden træder et selvstændigt
Væsen det Sande villig og gjerne imøde. Hvorved Bestræbelsen efter at erkjende
viser sig som en Sandhedens Cultus eller \emph{Dyrkelse}. Granskerens Forhold
til Sandheden og til hans Videnskab bliver nu et \emph{religiøst} Forhold, i den
ikke almindelige, men bogstavelige og vidtomfattende Betydning, hvori Ordet
betegner et ved hellige Baand knyttet Forhold til noget Helligt, i hvis Magt vi
med villig Hengivenhed føle os at være. (Ja ogsaa i Ordets mere almindelige og
egentlige Betydning kan Forholdet vise sig som et religiøst; thi ligesom
Mennesket kan og skal gjøre alt, hvad han gjør, med Gud for Øje og af Kjerlighed
til Gud, saa at dette højeste Motiv kan og skal være det overalt virkende,
saaledes kan ogsaa ikke blot det samme Motiv, som saadant Motiv, omfatte den
videnskabelige Stræben med, men ogsaa kan og skal alting tilsidst erkjendes og
skues saaledes i sit Forhold til det højeste guddommelige, at Ideen om Gud
videnskabeligen bliver den altoplysende, ligesom practiskt den altbesjælende.
Saa at tilsidst alt Forhold til Sandheden ligefrem viser sig som et Forhold til
Gud. Men heri see vi hen til et Ideal, som Menneskets Videnskabelighed endnu er
langt fra at have nærmet sig til). Med Ærefrygt for Sandhedens Stemme og inderlig
Lyst til at bringe den, og kun den allene, til at tale og leve i sin Gransken
maae Videnskabens Dyrker give sig hen til at være dens Organ, dens Tjener og
Forkynder, der kun vil, hvad den vil. Hans Personlighed maae være som et helligt
Tempel for den. Men tillige vorde det Sandes Talen og Leven i hans Indre
fornummet og fastholdt som hans egen inderste og bedste Leven, og Forholdet vise
sig som et sandt Kjerligheds Forhold. Hvad der stundom er bleven sagt, at
Mennesket er viet til sin Videnskab og sin Konst, dette kunne vi mindes her. At
optændes ved Sandheden og nu at optage og troligen bevare og pleje den i vort
Indre, indtil den selv voxer og uddanner sig og tager Skikkelse i os, for at den
selv ret kan komme til Fyldest i os, dette er vort. Kjerlighed til Sandhed være,
som al ægte Kjerlighed, omhyggelig og samvittighedsfuld.

Sandheden viser sig herved som den, der \emph{ved sin egen Magt og Kraft} søger
at tage Skikkelse i Individet, medens dette aldeles giver sig hen til at lade
den virke hos sig, hvad den selv vil. Og ligesom dette gjelder om alt, hvad der
i hvilkensomhelst Videnskab eller Kundskab viser sig som det Sande, paa hvis
Erkjendelse Studiet gaaer ud, saaledes gjelder ogsaa alt, hvad der her er
fremsat, om enhver videnskabelig Stræben, selv om den, der endnu staaer saa
langt fra de højere Ideer, at den kun sjeldent føler sig beaandet deraf. Under
\emph{ethvert} Studium vil den trolige Hengivenhed til den Erkjendelse, som
former sig i os, bringe den Fylde, Kjerlighed giver. Synes i Granskningens
lavere Kredse hvad vi erkjende, ikke at bringe os det Evige nærmere, saa vil dog
det, at det virker fuldkommen Erkjendelse og levende Overbeviisning hos os,
gjøre, at vi føle os som opfyldte af en Guddom. (Som da ogsaa, selv i de laveste
Sphærer, Erkjendelsen i sig selv virkelig kommer fra det højere Udspring, om den
end ikke kommer til Bevidsthed som saadan. See længere hen §.~11 i Anm.~2). Saa
at ethvert Studium, hvori troelig og klar Granskning hersker, har sin
\emph{Muse}. Individet træder deri lige over for noget, som det vil tilegne sig,
trængende dertil, som til noget, uden hvilket det selv ej kan komme til fuldelig
Leven. Men hvad det saaledes vil tilegne sig, er ikke noget, det kan behandle
som Gjenstand for en sandselig Fornødenheds Tilfredsstillelse, fortære eller
tage som et blot Vehikel eller Middel. Det bliver altid staaende lige over for
Individet, som noget, der bestaaer i og for sig i sin Betydenhed og Værd, hvad
enten det kommer til Individet eller ej, ja som noget, der netop kun forsaavidt
har Værd for Individet, som det har Værd og Betydenhed (det er paa dette Sted
har Sandhed) i og for sig. Og saaledes er overalt, hvor Sandhed søges, i hvilken
Sphære det saa er, Forholdet mellem Individet og det Sande, det søger, et
saadant Tilknytningsforhold, der kan, og under en ægte Stræben ogsaa vil, blive
et Kjerlighedsforhold.

Vi see her nu ogsaa, hvad det er at \emph{granske}, \emph{grunde},
\emph{eftertænke}, eller, som man mindre hældigen udtrykker sig, anstille
Eftertænkning, Undersøgelse, Betragtning over noget. At det nemlig ikke kan være
andet, end under denne eller hiin Sammenstilling og Behandling af visse
Forestillinger, at søge at bringe det dertil, at det Sande, vi attraae at vide,
kan komme til af sig selv, ved sin egen Magt og Indflydelse, at fremstille sig
for os, saa at vi nu ligefrem blive det vaer eller, som man siger, komme efter
det.

Vi maae derfor og \emph{fordybe} os i ethvert Studium. Kun derved bliver det
ægte og kraftigt. Skal det Sande, som virker Erkjendelse i os, faae os i sin
Magt, at det selv kan forme og danne sig i os, saa maae vi glemme alt andet
derover og ligesom tabe os hen deri. Endog Tingenes sandselige Former blive
ikkun da iagttagne og opfattede tilfulde i deres Ejendommelighed og
Mangfoldighed, naar Opmærksomheden derved ganske er fængslet, som om vi ganske
vare henne i Beskuelsen eller fortabte i den; thi kun da virker den heri
indtrædende levende nisus formativus, hvorunder det Sande vil tage Skikkelse i
os, heel og udeelt i sin fulde Styrke. En Bog studeres allerbedst og virker
allerfuldest, naar den en Tid lang saaledes fængsler os og ligesom tager vort
hele Sind i Besiddelse, at vi ganske ere fordybede i den. Paa denne Maade, det
er med denne kraftige Opmærksomhed og Hengivenhed, skal ethvert \emph{classisk}
Værk studeres.

Men ogsaa kun \emph{classiske Værker} fortjene og kunne ganske studeres
saaledes. Thi classisk er et videnskabeligt Værk ikke ved den Samling af
Kundskaber, som deri monne indeholdes og som deraf nu Stykke for Stykke kunne
øses og læres, men ved dets Aand i det Hele, dets indre Genius, der taler ud af
det i sin Heelhed, og dets dermed følgende totale Form eller Maaden, hvorpaa alt
deri er fremstillet og sammenstillet og ligesom voret sammen til Eenhed, saa at
vi overalt i Værket møde den samme kraftige og dybe Aand, og at alt det
mangfoldige deri er som af eet Stykke. Et saadant Værk staaer, ved sin Aand og
ved sin Form, gyldigt og betydningsfuldt for alle Tider, hvilket ogsaa Ordet
classisk især synes at tyde hen paa; thi det virker ei blot ved en Mangfoldighed
af Stof, der kan omplantes og omdannes, men det virker som Heelt ved sit
ejendommelige originale Præg, hvorved Aand taler til Aand. Men saa maae da ogsaa
denne Aand virke i sin Heelhed og ligesom gaae ind i Læseren eller Betragteren,
og dette gjør den desto mere, jo mere han fordyber sig i Betragtningen, saa at
han ej behandler det blot som Middel til en Kundskabs Erhvervelse, men som et
Værk, der har fuld Værd ved og i sig selv allene og derfor ogsaa virker til
selve Aandens Dannelse i det Hele.

Vi maae læse saaledes, at vi træde i et \emph{harmonisk Vexelvirkningsforhold}
med det, vi læse, og saaledes sættes i Bevægelse derved, at vi komme til
igjennem Ordene og Sætningerne at føle os tiltalte af det almene aandelige i
Værket, der udgjør Ordenes og Sætningernes indre Udspring og første Bevæger
(primus motor), den Sjæl, der lever og taler i dem. Kun en saadan Læsen er en
sand sjælefuld Læsen. Et Værks Studium bliver saaledes til en indre Omgang med
en herlig Aand, der i flere Henseender kan faae en velgjørende Indflydelse paa
os og fremme os i vor hele Stræben. Det bliver overalt desto mere til ægte
Næring for os, jo mere det er et \emph{hensynsfrit} Studium; thi jo mere vi
granske og grunde ej for disse eller hine besynderlige Øjemeders Skyld, men
allene af reen Kjerlighed til det Sande, for at optage dette aldeles i os, desto
mere vil det virke i os med sin fulde Kraft. Jo mere vor Hengivenhed er udeelt,
fri og af ganske Hjerte, desto lettere og fuldeligere vil den totale
Fordybethed eller Grebethed kunne komme. En Tid lang lade vi da Værket hentage
os og lade det gjøre med os og gjøre af os, hvad det vil. Derpaa begynder
Sandhedens egen Kraft selv at tale i os og at vidne om den i Værket udtalte;
hvorved en Fordybelse i vor egen umiddelbare Betragtning af det i sig selv Sande
træder til Fordybelsen i Værkets Udsagn, og nu bestemmer, hvad Værket skal gjøre
af os.

Overhovedet kan naturligviis den Aand, som lever og tiltaler os i et Værk, ikkun
komme til at hentage og fængsle os, og blive os en Kilde til nye Livets
Virkninger og Formationer i vort Indre, derved, at den kommer fra Sandheden
selv, tilmed fra det, der umiddelbart kan fremstille sig for enhver Gransker,
indtræde i enhver og hos enhver vidne om sig selv. Kun idet vi gjenfinde det
samme hos os selv, saa at vi sympathisere dermed, kan det os overleverede tale
til os igjennem de Former for Overleverelse, hvori det meddeles os. (Endog om
Naturens Værkers Studium gjelder dette; ogsaa Naturen er en Bog, der maa
studeres, som en anden Bog, men hvori ogsaa en os beslægtet Aand taler til os:
den samme Sandhed, der ogsaa i vort Indre vidner om sig selv. Dog dette vil
først fuldkomment indlyse i det Følgende. See §.~7 i Anm.) Til at fatte,
erkjende og nyde et os leveret Værk, komme vi i Ordets egentlige Betydning
ikkun, idet det \emph{reproducerer sig hos os}, og vi nu under denne dets
Reproduction bevæges og oplives af den Aand, der, ligesom den nu gjenlever i os,
saaledes først levede i den Virksomhed, hvorved Værket kom i Stand. Men herved
føres vi ligefrem til den næste Hovedbetragtning, idet vi nærmere tage denne
Reproduction, der viser sig som en saadan, hvorved vi selv ere frembringende, i
Overvejelse.

\begin{center}\large\S.~2.\end{center}

Al Erkjenden kommer til Virkelighed og Bestaaen formedelst en Bestemmen, Gjøren,
Frembringen, og ikkun formedelst en saadan; saa at denne væsentligen er
forbunden med hiin. \emph{Vi erkjende kun, idet vi frembringe.} (Scimus, qvia
facimus).

Der erkjendes, idet det Sande viser sig, som saadant for en Personlighed og
optages af denne. Men en Personlighed maae aabenbare sig i en personlig Leven.
En Leven kan nu ikke tænkes uden et Noget, et Reale, hvori den udtrykker sig i
uafbrudt Virksomhed, uafbrudt formende og bestemmende dette sit Reale eller sit
Ydre; den giver sig selv en bestemt Skikkelse, idet den giver dette en bestemt
Skikkelse. Personlighed derimod viser sig i at finde og sætte sig selv som
nærmeste Middelpunct for sin egen Tilvær, og at have denne sin Tilvær som sin
egen. En personlig Leven maa altsaa udtale sig i en Virksomhed og Productivitet,
i hvilken det personlige Væsen finder og bevæger sig selv, som det producerende.
Skal da nu det Sande optages i Personligheden, maa ogsaa det omfattes med af
denne personlige Productivitet. Individet kan derfor ikkun erkjende, idet det
selv virker og frembringer et til Erkjendelsen svarende Reale eller (sjæleligt)
Ydre, hvori det fastholder Erkjendelsen, idet det bestemmer og uddanner dette
Ydre.

Bevidstheden er Personlighedens første Udtryk. I denne optages alt, hvad der
skal optages i Personligheden. Saaledes kunne vi ogsaa kun erkjende ved at komme
til Bevidsthed om det, der skal erkjendes. Men dette kunne vi ifølge
Bevidsthedens Grundnatur ikkun ved at fremstille det for os eller forestille os
det, (objectivere os det). Men denne Fremstillen eller Forestillen er vor egen
Fastsætten, hvorved vi sætte (eller ponere) det givne og modtagne for os og
ligesom gjøre os et indre Afbillede deraf, i hvilket vi udtrykke det for os og
fastholde det. I denne Forestillen bestaaer her hiin Productivitet.
Forestillingen er vort eget Product og kan, som et Udtryk af vor egen Leven,
hvilket vi beholde som vort, ikke være andet. Men kun i denne bestandige
Forestillen og under dennes bestandige Bestemmelse, Uddannelse og Fuldendelse
komme vi til at erkjende, og have vi, hvad vi erkjende.

Erkjendelsens Besiddelse viser sig derfor og som en \emph{Kunnen}, en
\emph{Formaaen}. Vi vide kun (tilstrækkeligen og tilfulde) det, hvad vi ere i
Stand til at bringe frem under een eller anden Form og derved at give til andre.
Denne vor Productivitet, af hvilken nu det modtagne omfattes, er ogsaa i sig
selv en fri, der vilkaarligen kan fremstille andet end det Sande. I Erkjendelsen
bindes imidlertid dette frie, og bestemmes til en vis eneste Maade at forestille
paa ved den Erkjendelsen ledsagende Nødvendighed, som nu ogsaa i denne Henseende
kommer frem.

\begin{center}Anmærkning og Tillæg.\end{center}

1.~Vi see her \emph{Sprogets og Talens} væsentlige Sammenhæng med Erkjendelsen.
Al Erkjenden, der ikke gaaer ud paa umiddelbare eller reproducerede
Sandseanskuelser, eller paa Constructioner, der ligefrem fuldføres i Tid og Rum
paa samme Maade som disse, (som f.~Ex.\ de mathematiske) skeer formedelst en
bestandig Dømmen, der igjen skeer ved en indre Talen, formedelst Ord (eller
mueligen andre Tegn), og kan ikke komme i Stand paa anden Maade. Men at dette er
saa, hvilket Erfaring ligefrem lærer som Factum, viser ogsaa, hvor umiddelbart
en sættende, bestemmende og frembringende Virksomhed træder til vor Opfatten og
Modtagen. I Ord og Tale stille vi det erkjendte hen for os, binde det derved og
bemægtige os det, gjørende det til noget af vort eget, hvilket vi nu paa vor
Maade, i vor Betegnen og Benævnen, fastholde (jfr.\ Hegels Encyklopædie §.~379
og følgende §§.).

2.~Dog ej blot ved Ordene, overhovedet i al Forestillen gjøre vi det samme.
Forestillingernes Bevarelse, efterat den umiddelbare Skuen er forbi, saavel som
alt, hvad vi med dem kunne foretage, vidner i Almindelighed derom. Ogsaa viser
den til al Erkjenden væsentligen medhenhørende Productivitet sig snart for det
iagttagende Øje i enhver af Erkjendelsens Kredse. Hvad den simple \emph{lavere
Sandsning} angaaer henvise vi til de psychologiske Iagttagelser, der ligefrem
føre til at indsee, at rummelige Former og Forhold saavelsom Tidsbestemmelserne
og de Beskaffenheder, der med Hensyn til disse tillægges Tingene, kun erkjendes
formedelst en eftergjørende Construction af de sig fremstillende Gjenstande;
hvilke Iagttagelser allerede blot som saadanne lede til Slutninger, som
berettige til at antage det samme, hvad de blot synlige, hørlige og andre blot
sandselige Beskaffenheder angaaer. (See Forfatterens Psychologie §.~55).
Hvorledes \emph{Tænkningen} skeer, og den af samme udspringende Erkjenden
fremkommer, formedelst en bestandig Bestemmen, Sætten og Construeren, lære de
logiske Undersøgelser heel igjennem og overalt. Hvad enten concrete Gjenstande
begribes og deres Natur derved erkjendes i Kraft af Begreber, eller Begreber
selv betragtes som de Gjenstande, der skulle erkjendes, saa construeres der. Thi
det første skeer kun, idet en Sammenhæng indsees, men denne sees kun, idet det
mangfoldige, hvis Sammenhæng den er, sættes ud fra hinanden og derpaa igjen
forbindes til Heelhed, formedelst en eftergjørende Construction eller en genetisk
Fremgangsmaade, hvorunder vi lade Tingen opstaae for os, det er: selv i et
Billede danne os den, og nu ogsaa see alt paa sit Sted, idet vi see
Nødvendigheden af at sætte det saaledes ind i det Hele. Sammenhæng og
Gjenstandes reale Betydning i et vist sammenhængende Hele kan kun erkjendes
derved, at den dertil svarende Anskuelse udkastes og fuldføres; at erkjende den
er at producere eller dog kunne producere dens concrete Udtryk. Begreberne selv
dernæst, hvilke derved abstraheres, som de saaledes i Tingene skuede
Sammenhænge og derved bestemte Grundformer og Realbetydninger, eller som de i
Tingene selv sig udtalende Grundtanker, erkjendes kun, naar de hæves frem for
sig, ved at erkjende deres saakaldte Indhold, det er det Indbegreb af
Bestemmelser, der ifølge Begreberne maae findes hos Gjenstanden,%
% FLAG p.22: print reads "Gjenstandenr" (stray r); rendered "Gjenstanden". Verify (poss. "Gjenstandene").
 og som Prædicater kunne udsiges om dem. Hvorved da den foran nævnte Dømmen
indtræder, som en Bestemmen, Fastsætten og Udtalen. Begrebers Deductioner,
hvorved den Grundtanke, der bestemmer Begrebets hele Gyldighed og Indhold
udledes, eller sees i sit Udspring af en højere Erkjenden, skee, ligesom
Slutninger og Beviser, ved Tankegange, hvorved Tanker knyttes til Tanker ved
Bestemmelser, hvorved et Tankehele uddannes. Begrebers Deduction skeer kun ved
en Construction af det videnskabelige Hele, hvori Begrebet nu viser sig som
væsentligt Led; og Beviser fremkomme ligeledes kun ved Combinationer. Ogsaa en
Idee endeligen erkjendes kun, idet den enten skues som Udtryk af højere Idee, i
hvilken Henseende den fra Formens Side viser sig som et Begreb, eller ogsaa idet
den skues som Væsenet, det bestemmende Værende, og tilmed som det reale,
Væsenhed og Tilvær givende, Princip i reelle Gjenstande og Grundformer, hvilke
Gjenstande og Grundformer da, ikke blot som alle andre, kun formedelst
Constructioner erkjendes, men ogsaa kun kunne sees som Ideens Udtryk, idet de
sees, som de, der ved samme blive det de ere, hvilket kun kan sees under en
eftergjørende genetisk Fremstilling.

\begin{center}\large\S.~3.\end{center}

Modtagelse og Frembringelse, Skuen og egen Bestemmen og Gjøren gaae her aldeles
i eet, saa at Erkjendelsen, just idet den træder ind under den ene Form,
væsentligen ogsaa antager den anden Form. (Scientia et potentia in idem
coincidunt).

Productionen og Fremstillingen er i Erkjendelsen ikke en vilkaarlig, men en med
Nødvendighed bestemt, og det, der bestemmer den, er netop vor Modtagen og Skuen
og det i og for denne sig fremstillende Sande. Idet derved Forestillingens hele
Maade og Beskaffenhed bestemmes, udtrykker det Sande sig selv, men igjennem vor
Productivitet, hvilken det styrer og bestemmer, og i vort Product, hvilket nu dog
bliver hiint Sandes Udtryk, og ligesom bliver det Materiale, hvori det Sande
realiserer og individualiserer sig i os, medens dog i det samme ogsaa vor egen
Leven deri realiserer og substantialiserer sig. Vi bestemmes, idet vi skue, men
bestemmes til selv at frembringe noget, nemlig at danne Forestillinger, som
blive vore, dog at vi frembringe dem netop saaledes, at de svare til vor
Bestemtvorden. Og saaledes ere vi bestemmende og bestemte i eet. Det, der gives
os eller fremstiller sig for os, kommer kun til Klarhed hos os, ja overhovedet
kun til at blive vort, i det, hvad vi selv gjøre. Men dog erkjendes her kun
formedelst en bestandig Modtagen og Skuen; thi det, hvoraf vi bestemmes, have vi
kun i denne dets bestemmende Indflydelse, og Forestillingen antager ligeledes,
som nødvendig ifølge hiin Bestemtvorden, den samme Charakteer. \emph{Vi producere
som vi see, at vi maae producere.}

Personligheden, forsaavidt den finder sig som en Bevidsthed og Selvbestemmelse,
eller vi, som kalde os vi, befinde os her imellem \emph{tvende Regioner}. Paa
den ene Side Indflydelsen, hvoraf vi bestemmes, og til hvilken vi henvende os og
skue hen, idet vi modtage og bestemmes; paa den anden Side vor
Forestillingskreds, hvilken vi selv bearbejde og gjøre til det, den bliver, idet
vi deri søge at fastholde det modtagne og gjøre det til vort. Hiin, som
bestemmende, over os, foran os; som den, der skal udtrykkes og komme til
Skikkelse i os, det mere indre og virkende, det, hvoraf vi beherskes. Denne, som
vort Product, nedenfor os; som det, hvori vi udtrykke vor Leven, det mere Ydre,
vort Livs sjælelige Materiale, efter sin Natur Gjenstand for vor Behersken. Her
er altid paa den ene Side noget, som blot erholdes og haves ved det Herredømme,
det udøver i os, eller den bestemmende Indflydelse det har paa vor Virken og
Bestemmen; paa den anden Side noget, som blot er Udtryk, Fremstilling eller
Resultat, der kun haves og beholdes ved en Evne eller Færdighed til at
reproducere det.

Dette sidste kan nu blive tilbage som blot og bart livløst Materiale, ligesom et
Residuum af Virksomheden; som noget, der nu blot haves og antages, fordi det
engang har været antaget, som \emph{staaende Resultat}, uden at det endnu oplives
ved den levende Virken, hvorved det \emph{dannedes}. Men det skal ei mechaniseres
eller chrystallisere sig saaledes, men bestandigen holdes beredt til under Livets
Fremgang at gaae ind igjen under den levende Indflydelse af hiint bestemmende
Første, og under bestandig levende og organisk Gjenfødelse at modtage
Omdannelser og nye Skikkelser. Kun da bevares Forholdet, som det i sig selv er
og derfor ogsaa skal være. Saaledes som Virkningerne fra oven, fra inden, fra
det, der er foran os, tilsige og byde, saaledes skulle vi virke i vor nedre og
ydre Sphære.

\begin{center}Anmærkning og Tillæg.\end{center}

1.~I al Erkjenden formedelst Tænkning (Begreb og Idee) er denne Modsætning og
Dobbelthed klar. Begreber, som saadanne, haves og fastholdes kun enten formedelst
de Constructioner, hvis Normer de ere, eller formedelst de Omdømmer, hvorved
deres saakaldte Indhold bestemmes. Men Normen eller den bestemmende Grundtanke
for en Construction træder vel som styrende og herskende ind og frem i den, men
bliver ei selv construeret, den haves altsaa kun i den levende Virksomhed, der
gaaer frem af den. I Omdømmer udtales et Begreb, idet dets Indhold bestemmes,
men, da dette er et Indbegreb af Bestemmelser, der tillægges, ei Begrebet selv,
men dets Gjenstande, dog ifølge Begrebet, saa haves dette ogsaa her kun som
bestemmende Grundtanke for en Gjøren og Bestemmen, og omfattes ei af denne, men
svæver over den, eller udgjør den indre første bevægende Hersker i samme, dens
indre primus motor. Hvilket endnu mere gjelder om al Idee. Vel kan nu Begrebet
igjen betragtes som Gjenstand, der selv bestemmes, men dette kan da kun skee,
ifølge et andet Begreb, som saa udgjør den herskende Grundtanke i
Bestemmelserne. Saadanne herskende Grundtanker eller Begreber, der kun haves i
det, der i Kraft af dem sættes og udføres, ligge og til Grund for alle
Slutninger. Ogsaa i den lavere Sandsning træder et bestemmende første i
Modsætning til den derved bestemte Forestilling. I det der stræbes efter ret at
see og høre nøjagtigen, og at erholde en fuldkommen Forestilling om Gjenstanden,
hvori denne kan vise sig i sin fulde Ejendommelighed, vil, just fordi nu
Fornødenheden af den egne Gjøren træder frem, ogsaa det, der skal skues, vise
sig som det, hvortil der bestandigen maa skues hen, altsaa som det, hvorved hiin
Gjøren styres og beherskes, og som kun erholdes derved, at ifølge dets egen
bestemmende Indflydelse et Afbillede deraf fuldkomment uddannes.

2.~Den opfattende og erkjendende Individualitet viser sig her som den, der,
ligesom den i sine Forestillinger danner et Materiale, hvori den udtrykker sin
Erkjenden og Leven, saaledes ogsaa selv bliver dannet til at være det, hvori
Sandheden udtrykker sig; ret som om Individualiteten selv igjen var den sig
fremstillende Sandheds Substrat eller Materiale, hvori den søger at komme til
Fremtræden og til at tage Skikkelse (cfr.\ §.~4. Anm.~3.) \emph{Vi voxe derved}
paa den ene Side \emph{ind i den} Heelhed, hvori en almeen Leven udtrykker og
virkeliggjør sig, idet vi som Led eller Organer omfattes, opfyldes og
levendegjøres af den, medens paa den anden Side ogsaa et Hele voxer og uddannes
i os, i hvilket vi selv virkeliggjøre os som personlige for os selv bestaaende
Væsener. Vi assimileres selv af hiin almenere Leven, i det samme at vi søge at
assimilere os det sig for os fremstillende Sande, ved at assimilere os dets
Udtryk i os.

Jeg henviser herved til%
% FLAG p.28: a printer's manicule is set before "den geniale"; reading "til" assumed. Verify.
 den geniale Franz Baaders Afhandling i Schellings og Markus's Jahrbücher der
Medicin tredie Binds første Hefte, samt til flere af hans
philosophisk-biologiske Smaaskrivter.

\begin{center}\large\S.~4.\end{center}

I den simple, rene, aldeles umiddelbare Skuen er Productiviteten \emph{bunden} og
\emph{fortabt}, eller gaaer deri saaledes op, at den ikke for sig kommer til
Syne, eller kommer til Bevidsthed, og rører sig efter sin særlige Natur. Ligesom
ei heller Indflydelsen med dens bestemmende og nødende Magt fornemmes og kommer
til Bevidsthed, som saadan. Modsætningen mellem det erkjendte Sande og den
Individualitet, i og for hvilken det træder frem, viser sig her ikke som en
Modsætning. Denne Erkjenden er ogsaa, i Henseende til Maaden at erkjende paa, den
fuldkomneste. Thi, at Productionen ganske kommer af sig selv, beviser, at
Sandheden virker med fuldkommen Magt, og at den ogsaa optages tilfulde.

Men under en Stræben efter at erkjende det sig i denne Skuen umiddelbart
fremstillende ret tilfulde og tilligemed at erkjende endnu andet og mere, fordi
det saaledes erkjendte viser hen paa andet og mere, træde formedelst en
Reflection \emph{Erkjendelsens Gjenstand og den erkjendende Personlighed},
hvilke i den rene Skuen vare uadskilte, \emph{ud fra hinanden} og imod hinanden.
Productiviteten bliver løsbunden og viser sig som fri, medens den bestemmende
Indflydelse paa den anden Side viser sig, som saadan, med sin Nødvendighed.
Individet finder sig, som det, der selv maa arbeide paa formedelst Overvejelse og
Tænkning at komme til den attraaede Erkjenden, og just derfor maa Productiviteten
vise sig som en fri, der kan fremstille andet end det forefundne. Men nu vil
derfor ogsaa Tilknytningen omfattes af Selvbestemmelsen og vil skee ifølge en
Villen. Sandheden viser sig nu som en Gjenstand, vi først ved Anstrengelse
efterhaanden kunne komme i Besiddelse af ved en fri Virksomhed, der skal stemme
overeens med Sandhedens Tilsagn, paa hvilke vi nu maae agte nøje, saaledes som
de efterhaanden under vor Gransken lyde til os. Hvorved da Productiviteten,
medens den viser sig som fri, bindes og beherskes paa ny, saa at den kun
fremstiller paa den ene sande Maade, men dog saaledes, at den nu skeer med
Selvbestemmelse og som den, der kan virke anderledes; derfor ogsaa ifølge
Overvejelse, og saaledes at et frit Valg er muligt. Men reen aldeles umiddelbar
Skuen, ved en ny fuldkommen Fortabthed af den til Klarhed og Selvbevidsthed
komne Productivitet, er dog det, hvortil der stræbes, for at Sandheden i sin
fulde Fylde kan indtræde i den erkjendende Personlighed.

\begin{center}Anmærkning og Tillæg.\end{center}

1.~Den friblevne Productivitet kan nu ogsaa \emph{ytre sig i egne}, af det givne
uafhængige, Productioner. Men altid vil den dog ogsaa derunder, da den maa skee
paa bestemte Maader, bestemmes ved visse Normer eller Begreber, hvilke af sig
selv herske i den og ligefrem haves i og ved den Magt, hvormed de bestemme
Selvvirksomheden. Og skulle disse Productioner være mere end et Spil med
Forestillinger, saa maae denne friere Productivitet, uagtet dens Frembringelser
nu ei skulle blive blotte Afbilleder af det forefundne, men dens egne
Skabninger, dog paa samme Maade, som under Bestræbelsen efter at erkjende, paa
ny bindes og beherskes ved en højere Erkjenden, hvorunder det Sande under nye
Former viser sig for os og tager Skikkelse i os. Saaledes er det i Mathematiken
og i Konsten.

Overhovedet træder Livet ligesaa fuldt ud i Handling og Selvvirksomhed, som i en
Stræben efter at erkjende. I hiin hersker da den frie Fremvirken, men tillige en
vis Erkjenden, der ligger til Grund som styrende og bevægende Første. Men ogsaa
heri kunne de to Modsætninger paa samme Maade være fortabte eller gaae op i eet,
som i den blotte Skuen, saa at hverken Selvvirksomheden særligen fornemmes som en
fri, ei heller den Grundtanke, der leder den, kommer til Bevidsthed, som saadan,
idet Handlingen skeer ganske af sig selv: en Selvvirksomhed, som i Henseende til
Virkemaaden, er den fuldkomneste da deri det Sande og Værende, den Væsenhed, som
ogsaa under vor Handlen skal tage Skikkelse i os, viser sig, som det, der virker
med fuldkommen, udeelt og ustandset Magt.

2.~Saavidt den frie Productivitet kommer til Bevidsthed som saadan, kommer og
Erkjendelsens bestemmende Første til Bevidsthed som saadant; og omvendt. I en
fri Arbejden paa at erkjende, ligesom i en med selvbevidst Selvbestemmen
foretagen Handlen, kan der derfor, medens \emph{noget} kommer særligen \emph{til
Syne}, være \emph{noget andet} dybere eller mere indre, som med dets bestemmende
Indflydelse er fortabt eller gaaer op i Erkjendelsen eller Selvvirksomheden, i
hvilken Henseende da den, i anden Henseende friblevne, Productivitet ikke viser
sig som fri. Men en Reflection kan da siden, under en Stræben efter dybere
Erkjenden, drage frem, hvad der saaledes hidtil vel var tilstæde, men var
fortabt i det Hele. --- Saaledes kan f.~Ex.\ Aarsagelighedsprincipet een Gang
ligge som bestemmende Grundtanke til Grund for en Stræben efter at udgrunde
visse Phænomener, uden at der tænkes paa dette Princip, en anden Gang
Reflectionen henvendes derpaa; og den Standpunct, paa hvilken den erkjendende
Stræben staaer i hiint Tilfælde, er en ganske anden, end i dette.

3.~Men hvorfor viser sig Tingenes sande fulde Væsen ikke strax saaledes for det
efter Erkjendelse stræbende Individ, at al Erkjenden bliver en reen umiddelbar
Skuen? Det er naturligt, at dette Spørsmaal her opstaaer. Men til dets
fuldkomne Besvarelse vilde udfordres en almindelig Eftertænkning over Grunden
til, at der overhovedet er en Tid og en Timelighed, og at Livet overhovedet
træder frem i en Vorden, hvorunder det kun efterhaanden kommer til en bestemt
Bestaaen, noget, til hvis fuldelige Oplysning en Grunden over al menneskelig
Tilværs inderste Betydning og Væsen vilde være fornøden. --- Kun de første
Grundtræk i en saadan Betragtning kunne faae Plads her.

I den endelige Verden skal det Evige og Guddommelige komme til at træde frem,
herske og raade, for deri at aabenbare sin uendelige Magt og Fylde: dette er
Betydningen af alt Endeligt, at det skal være \emph{det Andet}, end det Evige,
hvori dette skal komme til at træde frem og aabenbare sig. Men kun i en
Bevidsthed og aandelig Leven aabenbarer sig det Evige og Guddommelige \emph{som
saadant}, idet det Væsen, hvori det lever, kommer til at vide af det, og tilmed
ei blot at bære det i sig og at bevæges deraf, men ogsaa selv at kunne skue og
have det som det, det i sin Sandhed er. Men det Rige, der saaledes skal komme
til Mennesket, skal være et \emph{Kjerlighedens Rige}; thi kun i et saadant
kommer den Eviges sande Fylde til at vise sig, og at vise sig i sin rene
Aandelighed, enigende et Samfund af Væsener, der selv ere Aander, til et Hele,
der nu i Sandhed kan kaldes Guds Rige, som det Rige, hvori Gud selv aabenbarer
sig i sin Guddom, idet han aabenbarer sig i evig Kjerlighed. Men et Aandens og
Kjerlighedens Rige forudsætter \emph{nødvendigviis}, og fremkalder derfor
nødvendigviis, som det andet, hvori det skal være det herskende, \emph{en
Intelligents og Frihed}. Dersom fra de Væseners Side, til hvilke dette Rige skal
komme, ikke en fri og velbevidst Hengivenhed kom til at finde Sted, nemlig en
Hengivenhed af eget Hu, egen Lyst og Drivt, saa vilde det ei være et
Kjerlighedens Rige, som vilde komme frem. Men en personlig Leven, der skal være
en fri, som nu af reen Kjerlighed og tilmed med klar Bevidsthed skal give sig
hen til at levendegjøres og opfyldes af Guddommen og kun at røre sig i Gud, maa
selv gjøre sig til en saadan fri Leven, idet den kommer til Klarhed over sin egen
Tilværs, og tillige over hele Tilværelsens, sande Betydning og Væsen. Men saa
maa da Modsætningen ret træde fuldkomment frem, for at Enigelsen kan blive den
sande Enigelse og Eenhed i Kjerlighed og Frihed. Men skal da nu saaledes
Bevidstheden komme til Selvbevidsthed og Intelligents, selv sigende sig, hvad
den er, skal Personligheden komme til Frihed under en Selvbestemmen, hvorunder
den selv tager sig som saadan og gjør sig fri, saa maa den komme til en
Productivitet, i og under hvilken den objectiverer sig for sig selv. Den maa
komme til at udtale sig i en Sphære (hiin foromtalte lavere Region), hvilken er
dens egen, som den selv danner sig, og i hvilken den bliver sig sig selv og sin
Frihed bevidst i en fri Productivitet. Og saaledes stiger da Livet ned i en Tid
og Timelighed.

Vel kan nu det personlige Væsen ligesaa lidet gjøre sig selv fri, (uagtet vi kom
til at bruge dette Udtryk), som den selv kan give sig Sandheden. Det finder sig
som fri, hvad Muligheden angaaer, men til i Gjerningen (actu) at have og nyde
Friheden kommer det kun ved Sandhedens og Kjerlighedens frigjørende Magt. Men nu
kan det dog ej heller frigjøres, uden \emph{i og formedelst dets egen
Selvbestemmen}, og det en Selvbestemmen, hvilken det selv vil og føler som sin
egen. Det, der bringer og giver Individet al Realitet, Væsenhed og Fylde, kan
dog kun, da det er et Kjerlighedens Rige, der skal danne sig, give det Fylden ved
at referere sig til \emph{noget}, som nu hos Individet selv maa \emph{paafølge},
som dettes egen Sigbestemmen, nemlig dets egen Sighengiven og deraf følgende
Handlen. Beslutningen (det, hvori den hele Virken slutter eller ender) maa altid
være Individets egen. Derfor have da og --- for at vi endnu skulle bemærke
dette --- de indre Virkninger, hvori det Liv, der skal komme til Realitet,
opererer eller trænger sig frem i Bevidstheden, Charakteren af Tilskyndelse og
Opfordring, nemlig Opfordring til at søge, hvad Individet kun kan faae, naar det
attraaer og begjærer og søger det, men hvad det visnok ogsaa kun kan faae
derved, at det reent ud gives samme. Dog derom længere hen igjen; see
§.~19--21. For at Friheden fuldkomment kan komme til Bevidsthed og en fuldkommen
Selvbevidsthed kan indtræde, dog saaledes, at denne tillige klarligen skuer,
hvad den skal stræbe at gjøre sig selv til, saa at nu Aanden og Sandheden kan
haves i Aand og Sandhed, derfor oplader denne sig først efterhaanden for
Individet under dets egen Tragten, Gransken og Arbejden.

\begin{center}\large\S.~5.\end{center}

For saavidt Modsætningen imellem Erkjendelsens Gjenstand og den erkjendende
Personlighed træder saaledes frem, at Productiviteten eller Forestillingen viser
sig som en fri, der nu dog bringes til, formedelst indtrædende Selvbestemmelse,
at fremstille paa en vis eneste Maade, bliver Erkjendelsen til en Erkjenden i
Kraft af Grunde.

En saadan \emph{Erkjenden i Kraft af Grunde} skal al Erkjenden være.

Thi al Erkjenden skal komme til fuldkommen Klarhed i og for den erkjendende
Personlighed og Bevidsthed. Men skal den det, saa maa og den dertil fornødne og
deri liggende Productivitet komme til klar Bevidsthed; men dette kan den ikkun,
idet den skeer eller foretages af Individet med Selvbevidsthed og
Selvbestemmelse. Derved viser den sig da som en fri; bliver den da nu dog, som
den, formedelst hvilken en Erkjendelse skal komme i Stand, bunden til en eneste
bestemt Fremstillen, saa maa den bindes dertil ifølge et vist tilsvarende
bestemmende Noget, der, naar det kommer til Bevidsthed som saadant, viser sig som
en Grund.

Erkjendelsen fremkommer nu som Resultat af en Undersøgelse og Betænkning.
Reflection og Overvejelse fremkalde et \emph{Studium}, som træder ind imellem en
første umiddelbar Modtagen og en endelig eller sluttelig Erkjenden. Sindet, som
først aabnede sig for blot at modtage og gav sig reent hen i Beskuelsen og
Modtagelsen, sættes nu i Bevægelse saaledes, at det selv ved sin egen
selvbevidste Arbejden skal bringe den Production i Stand, hvori det skal besidde
Sandheden, og søger nu, formedelst en Gransken og Grunden, et Princip eller en
første Bevæger for denne sin Production. Men den derved indtrædende Erkjenden er
dog, naar den nu kommer i Stand, og saaledes, som den nu bliver, en ny Modtagen
og Skuen; thi vi antage dog tilsidst, hvad vi ifølge de bestemmende Grunde see,
at vi maae antage. Ogsaa skal denne Erkjenden i Kraft af Grunde være en
Gjennemgang til en fuldkommen Skuen, hvori Sandhedens fulde Magt umiddelbart og
udeelt virker. Men hos den frie Personlighed, der skal komme til fuldkommen
Frihed og Selvklarhed, viser den sig som en nødvendig Gjennemgang.

\begin{center}Anmærkning og Tillæg.\end{center}

1.~Grunden kan iøvrigt gjærne være en med den rene umiddelbare Opfatten og Skuen
følgende Nødvendighed allene, saa at Erkjendelsen finder Sted, som den, der
begrunder sig selv; men Spørsmaalet om en Grund er dog nu kommet i Bevægelse, og
om end det, der antages, blot bliver antaget, fordi det umiddelbart ved sig selv
er vist, saa viser nu netop denne dets umiddelbare Vished sig som Grund, og
Erkjendelsen faaer nu en anden Charakteer, end den forhen havde, da denne
Reflection endnu ej var indtraadt.

Dette er i Særdeleshed at erindre om de mathematiske Axiomer. Ogsaa
Sandseanskuelsen træder vistnok op som umiddelbart begrundende sig selv og kan
ogsaa blive staaende med dette Præg. Men allerede da anlediges dog vor
Forestillen ej sjeldent til at finde sig som en fri eftergjørende, især om den
sandselige Skikkelse er saa mangfoldig eller dens Deles Forhold saa forviklede,
at den kun ifølge Betænkning og Overvejelse kan komme til Klarhed, og saa vil da
ogsaa det, der leder til at fixere Forestillingen, komme til, i Modsætning til
denne selv, at vise sig som dens Grund. Dette skeer endnu mere, naar det
bemærkes, at Tingene ofte ikke kunne eller skulle forestilles at være saaledes,
som Indtrykket umiddelbart vilde føre det med sig. Men naar saadanne saakaldte
Sandsebedrag først have viist sig som muelige, kan der nu opstaae Tvivl om
enhver umiddelbar sig selv dannende Sandseforestillings Rigtighed, saa at, om
den end ligefrem antages, skeer dette nu ej mere blot ifølge dens egen
Umiddelbarhed, men tillige ifølge en Erkjendelse af, at intet saakaldet
Sandsebedrag finder Sted, og den Overvejelse, som lærer dette, træder nu ind som
Grund. Ja under Bestræbelsens og Erkjendelsens Fremgang viser tilsidst (for det
philosophiske Blik) den hele Sandseanskuen sig som en saadan, hvis Sandhed og
Gyldighed overhovedet kan omtvivles og tilmed først maa hjemles og godtgjøres.
Og saaledes vil da den Forestillen, som fra først af dannede sig selv ganske af
sig selv og syntes tilforladelig nok i sig, vise sig som den, der overhovedet og
i eet og alt endnu behøver Begrundelse.

Det i sig selv visse, hvorfra denne Begrundelse skal gaae ud eller hvorved al
anden Erkjenden skal hjemles, søges nu paa en ganske anden Side, og Erkjendelsen
stiger til en højeste Idee, i Kraft af hvilken nu alt andet skal skues i sit
rette og fulde Lys, i sin Betydning og Sandhed. Men om end denne højeste Idee
forsaavidt er umiddelbar og hvilende i sig selv, saa komme vi dog kun til den
ifølge en Tænkning, der naturligviis heelt igjennem maa være ledet og bestemmet
ved Grunde. Som da selv det, at alt andet kan opklares ved hiin højeste Idee,
kan være den Grund, hvorpaa dens Erkjendelse hviler. Og saaledes finde vi lige
saa vel paa det højeste Punct en Erkjenden i Kraft af Grunde, som vi fandt den
laveste trængende til Begrundelse, og det fordi al Erkjenden kun kommer til
Virkelighed formedelst en Bestemtvorden til en vis Forestillen, en
Bestemtvorden, som under Bestræbelsens Fremgang maa komme til Bevidsthed som en
saadan. Der spørges overalt, hvorfor vi lade os saaledes bestemme, med hvad Ret
denne bestemmende Magt udøves over os.

2.~Den fuldkomne Erkjenden vilde være den, der, uafhængig af al Betænken og
Grunden med al dennes Uro og stræbende Bevægelse, vilde hvile ganske i sig selv
allene og umiddelbar være klar og vis i sig, uden at et Spørsmaal om dens Grunde
kunde opstaae eller en Tanke paa, at noget andet kunde være det sande, end netop
det skuede. Men saadan Vished og Tryghed, en i en saadan Skuen til Ro og Fylde
kommen Erkjenden kan kun hvile i Visheden om, ja i den rene umiddelbare
Fornemmelse eller Følelse af, at det, der saaledes skues, ligefrem kommer til os
fra Sandhedens Kilde, gives og siges os af den, fra hvem al Sandhed, ligesom
overhovedet al god og fuldkommen Gave, kommer. Dette er den Skuen, man har villet
betegne med det Udtryk at skue alle Ting i Gud. Det er den Erkjenden, ved hvilken
vi skulle erkjende saaledes, som vi selv ere erkjendte; vide, hvad vi vide, fordi
Aanden taler det saaledes i os og giver os det. Da staaer Sandheden ej mere lige
over for os, som den, vi søge; men den boer i os, idet det samme, der giver alle
Ting Væsenhed, nu ogsaa giver os Erkjendelsen, reen, umiddelbar, absolut eller i
og for sig, uden at den gjennemgaaer en Betænken og Eftertænkning og kommer ud
som et Resultat af en saadan. En Bevidsthed og Forvisning om, at alt, hvad vi
erkjende, er os givet ved Gud, at vore Ord ej ere vore, men saaledes, som vi
finde dem i os, virkes i os ved den guddommelige Aand, kan allene være det, hvori
denne Erkjendelse kan hvile. Christi Erkjenden kunne vi kun tænke os paa denne
Maade, ingensinde som en Erkjenden ifølge Reflection og Slutning.

Som højeste Maal, som et Ideal, staaer en saadan Erkjendelse for os. Dog ikke
engang som det Maal, til hvilket den blotte Gransken og Grunden, naar den engang
vilde naae sin Fuldendelse, skulde kunne komme, kan denne højeste Skuen
betragtes. En fuldendt Tænkning kan føre til en fuldendt Erkjendelse ifølge
Tænkning, til en fuldendt aldeles vis, klar og i sig selv hvilende Viden ifølge
de højeste Grunde og i Kraft af de højeste Grunde. Men til hiin højeste Skuens
rene Umiddelbarhed kan det kun komme derved, at det højeste Sande, Sandhedens
evige Kilde, Aanden, paa en ganske anden (en anderledes reel) Maade bliver
levende i os og levendegjør os, end den Maade, paa hvilken Sandheden, som højeste
Idee, besjæler og levendegjør en Tænken og en Grunden.

For det fremadstræbende Menneske er imidlertid \emph{Tænkningens Vej} den, ad
hvilken han skal fremad, for at trænge mere og mere frem til Sandhedens fulde
Erkjendelse. Det er en mere jordisk Vej, paa hvilken der opnaaes Stykke for
Stykke. Den hører til den Sphære, hvori i en Endelighed, under sin egen Stræben,
en Personlighed skal uddanne sig og bestandigen mere og mere skal hæves frem til
Klarhed og Selvstændighed, i et Liv, hvori den mere og mere objectiverer sig for
sig selv og sætter sig fast i en egen Bestaaen, men dog kun for at modnes til, at
Aanden nu i sin Fylde kan træde ind i den som Organ nu fuldeligen uddannede
Individualitet; hvorved da, naar dette Punct naaes, Fylden vil indtræde
istedetfor Stykværket, og nu Tænkningen, der vistnok ikke skal høre op, Tingenes
Udgrundelse, der ikke skal falde bort, selv blive oplyst af det højere Lys, gaae
ind i den højere Fred, og medens en uendelig Fylde staaer for den til Beskuelse
og Udgrundelse, dog selv gaae den trygge Gang, der ej mere bestandigen har sig
selv at spørge om Grundene, hvorfor den antager, hvad den skuer, hvorfor den
erkjender, hvad den under sin Grunden finder, men i denne Henseende har udstødt
„ethvert Vidne om menneskelig Nødtørftighed.“

Saa længe, indtil det kommer dertil, er al Erkjendelse behæftet med
\emph{Uvished}. Det umiddelbart givne, som ligefrem skues, kan ej tilfredsstille
os, fordi det ej er det sande Væsen, paa hvis Erkjendelse alt kommer an, og fordi
vi om dette sande Væsen have en Grundbevidsthed, der driver os til at søge noget
andet end hiint ligefrem skuede. Det selv viser os nu hen til noget andet og
højere, men dette andet og højere komme vi kun til netop ved saaledes at vises og
ledes hen dertil, altsaa paa en middelbar Maade, og denne Middelbarhed bliver
det, saa længe vor Erkjenden er en Erkjenden formedelst Tænkning, altid beheftet
med for os.

Men hvad der dog her allerede kan hæve op over denne Middelbarhed, er den sig
dannende faste \emph{Overbeviisning} om dets Realitet. Denne udspringer under vor
Gransken og Grunden, idet vi gribes af det Erkjendtes Sandhed; føle os hentagne
deraf. Under vor Sammenstillen og Henføren, Betænken og Overvejen, hæver den sig
mere og mere frem og bliver i subjectiv, individuel og personlig Henseende, (hvad
Begrebet og Indsigten i objectiv Henseende ere) den hele Tænkens og Bestemmens
bevægende Sjæl, hvorfra den gaaer ud og hvortil den tilsidst gaaer tilbage. Og
saaledes indtager da en, al Mangelfuldhed i Indsigten og i Tænkningens
Udviklinger og Fastsættelser overvindende og opløsende \emph{Tro}, nemlig en
\emph{Tro paa Grunderkjendelsens Realitet} for os hiin fuldkomne Erkjendelses
Plads. Som Tro er den, om end vakt og næret ved en Grunden og bestandigen fødende
en nærmere Betænken og Begrunden af sig, dog ikke nogen Erkjenden ifølge og i
Kraft af Grundes Udvikling allene, men en mere umiddelbar Fremtræden af
Sandhedens Magt, der vidnende om sig selv griber og uimodstaaeligen nøder til
Erkjendelse, dog ikke heller fordi der umiddelbart skues, men fordi alt saaledes
viser hen og trænger os hen til denne Erkjendelse, at den viser sig som det
Middelpunct, hvorpaa alt kommer an og hvori alting hviler, og uden hvilket vi
aldeles ikke kunne bestaae eller komme ud. Ligesom vort Liv overhovedet, saa maa
og vor Leven i Erkjendelsen søge sin sidste Støtte og Hvile i en uendelig Tillid
og Tiltro, og kan som et endeligt Væsens Leven tilsidst ej hvile i andet. Dog
dette, at ogsaa her, ligesom overalt, den guddommelige Naade tilsidst er det,
hvorpaa alt allene kommer an, og hvori alt beroer, dette komme vi i det følgende
til nærmere at foreholde os. Men fordi vi kun modtage af Naade, derfor er det
ogsaa, vi kun kunne have med Tro. Alt er tilsidst givet og maa tages, som det
gives, i Tro til den Givende. Men netop det, hvad der gjør, at Erkjendelsen ikke
bliver en total Skuen, aabner nu ogsaa denne uendelige Tillid det Rum, hvori den
kan vise sig som uendelig Tillid, saa at just det, at vi ej have Fylden,
medfører, at Tilknytningen, Hengivenheden og Kierligheden desto mere kommer til
Syne.

Og det er det da, hvortil der sigtes med al Livet i os, og tilmed ogsaa med al
den Leven der arbejder paa at naae en fuldkommen Erkjendelse: at vi, jo mere vi
anstrenge os for at have noget ved os selv og besidde noget i os selv, desto mere
skulle komme til at føle, at vi intet kunne og intet have, uden hvad der ligefrem
og reentud gives os, og at vi, jo mere vi stige i Selvstændighed og Intelligents
og Raadighed, hvad det mangfoldige og den relative Bestemmelse af samme angaaer,
desto mere skulle erkjende, at vi i Henseende til det Ene, som gjøres nødigt og
hvorpaa alt kommer an, ganske maae give os Gud i Vold, der allene virker i os
baade at erkjende og at udføre.

3.~Men det er dog kun Grunderkjendelsen, (det, hvorpaa alt kommer an og som for
saavidt er det Ene, som gjøres nødigt), om hvilken dette gjelder, at den tilsidst
hviler i en Tro og Tillid. Alt andet, som ved denne Grunderkjendelse oplyses og
ifølge samme skal bestemmes, maa skues som det, der er givet med den, alt i sin
Sammenhæng dermed og sin Følge deraf. Saa at, om end paa den ene Side en Tro og
Tillid bliver det, hvori al Overbeviisning hviler og ikke kan andet end hvile,
saa skal dog Erkjendelsen derfor ikke høre op paa den anden Side at være en
Erkjendelse i Kraft af Grunde og af en ved Grundes Erkjendelse opstaaet
\emph{Indsigt}. Og ligesom alt, hvad der ifølge Grunderkjendelsen erkjendes, maa
erkjendes saaledes, fordi vi klarligen see, at Grunderkjendelsen nødvendigviis
fører det med sig og kunne forfølge den hele Sammenhæng fra det enkelte
betragtede Punct indtil det højeste, saaledes maa og Grunderkjendelsen selv netop
klarligen sees som den, der bærer den hele Sammenhæng, og til hvilken denne
Sammenhæng nødvendigviis fører; hvorunder da Gangen vil blive den, at det, der
maaskee fra først af erkjendtes, ifølge en vis Sandhedsfølelse, ja en Relation
til det subjective Menneskelige, mere og mere kommer til at indlyse for
Tænkningen i og for sig og paa en objectiv Maade, indtil tilsidst et Punct vil
opnaaes, hvor det højeste, hvortil alt andet viser hen, ikke blot indlyser, fordi
det klarligen sees, at alt andet viser hen derpaa, men fordi det synes
umiddelbart \emph{evident i sig selv allene}. Men denne Betragtning, i hvilken vi
mødes med den dybsindige Tanke, som ligger til Grund for det berømte saakaldede
ontologiske Beviis for Guds Tilværelse, fører os over til vor næste Hovedsætning.
--- Gud vidner om sig selv i Bevidsthederne, og ingen Erkjendelse af, at Gud er,
vilde være muelig, dersom Gud ikke saaledes vidnede om sig selv i vort Indre,
d.\ e.\ i vor sig efterhaanden dannende aandelige Leven, Skuen og Betragten; ja
denne Grundbevidsthed om Gud (og om sit eget inderste Samfund med Gud) gjør
Mennesket først til Menneske. Heri er da umiddelbar Tro. Men skulde nu ikke den
samme Erkjendelse ogsaa kunne blive ligesaa umiddelbar evident for
Speculationen, saa at ogsaa denne Evidents, som Evidents kun hvilede i sig selv?
Hvormed dog, da ogsaa denne Evidents saa maatte gives os, idet den netop ved sig
selv indlyste, Troen ikke vilde høre op at være det, hvori al Erkjendelse
tilsidst vilde hvile. Selve Indsigten og Indlysningen eller Evidentsen
\emph{virkes} dog ogsaa i os, men hvor der er Tale om et paa en reel Maade sig i
Bevidstheden fremarbejdende Liv, og det, der fremkommer, sees som dette Livs indre
Operationer, der er der Naade og Styrelse, og tilmed maa der der vise sig Tro,
Tilknytning og Tillid. Det rene intelligente er selv noget, som ogsaa kun kommer
ved reelle Virkninger fra oven og maa derfor tages som Gave, tilmed med Tillid
til den Givende. I den Erkjenden, vi tænke os som Christi Erkjenden, gaaer en
umiddelbar Indlysen i eet med den Forvisning, Tillid giver; den er Tro og Skuen i
eet.

\begin{center}\large\S.~6.\end{center}

I det Granskningen gaaer ud paa en Erkjenden i Kraft af Grunde, tragter den efter
ej blot at erkjende det Sande, men og at erkjende det som det Sande og ifølge
dets Sandhed. En saadan Erkjendelse af Sandheden i \emph{dens Sandhed} skal al
Erkjenden være. Kun da er den ægte og fuldkommen. Det er Bestræbelsen efter at
erkjende saaledes, der sætter den hele Gransken efter Grunde i Bevægelse; den er
det, som egentligen udtaler sig i denne.

Under den første umiddelbare Opfatten og Skuen modtage vi, hvad der umiddelbart
fremstiller sig for os og lade det dermed være nok. Om dets Gyldighed er der endnu
ej Tale; være det og det Sande, som deri har viist sig for os, saa er det dog
endnu ej traadt frem som det sande; det er det endnu ej for os; (hvorfor og
Platon i Theætet søger en ganske anden Definition paa Erkjendelsen, end den, at
den skulde være den rigtige Forestilling). Men skal Sandheden fuldkomment komme
til os og tage Skikkelse i os, saa maa den selv give sig os ved sig selv, derved,
at den indlyser for os ved sin Sandhed; kun da skue og have vi den, som den selv
er. Derved indtræder først den sande Erkjenden, der ej blot er en Optagen og
Skuen, men en Optagen med Forvisning. Men dette skeer allerede, om end ej
tilfulde og i højeste Maade, forsaavidt Erkjendelsen bliver en Erkjenden i Kraft
af tilstrækkelige Grunde, og saa langt som denne strækker sig; thi saavidt og saa
langt antage vi, hvad vi antage, fordi det ifølge disse Grunde indlyser for os
som det, der under disse Grundes Forudsætning er det eneste antagelige, altsaa
som det, der allene kan erkjendes for det sande.

Men skal Erkjendelsen blive den fuldkomne Erkjenden af Sandheden i dens Sandhed,
hvori denne lyser ved sit eget Lys, saa maa paa ethvert Punct Grundene, af
hvilke den gaaer frem og hvoraf dens Gyldighed forsaavidt er betinget, forfølges
indtil de \emph{højeste og sidste Grunde}; saa at, hvad der erkjendes, erkjendes
i Kraft af disse, der iøvrigt, betragtede fra en anden Side, ogsaa kunne kaldes
de \emph{første Grunde}.

Disse, der da ej igjen kunne føres tilbage til højere Grunde, maae nu indlyse
umiddelbart ved deres egen Sandhed.

\begin{center}Anmærkning og Tillæg.\end{center}

1.~Den frie Selvbestemmelse og Productivitet med den egne Betænken og
Overvejelse, som hos den individuelle Personlighed fremkaldtes og traadte
Modtagelsen prøvende imøde, sattes altsaa kun i Bevægelse, for at Sandheden
fuldkomnere kunde blive erkjendt og optaget. Et højere Trin i Erkjendelse skal
den erkjendende Personlighed derved stige til, Friheden hæves frem for desto
fuldeligere at høre Sandheden til. For at denne kan lyse for os i sit eget Lys og
derved fuldkomment vinde Skikkelse i os, derfor gaae vi ind i den bestandige
Prøvelse, derfor kræve vi overalt Rede og Rigtighed, som man udtrykker sig, kræve
i Videnskaben Grunde og Beviser, endog der, hvor, f.~Ex.\ ved visse af de
simpleste mathematiske Sætninger, Visheden ellers er urokkelig nok og ingen Tvivl
opkommer eller kan opkomme. Det er ikke nok, at det Sande gives os, det skal gives
os ved den fuldkomne Indlysen af dets Sandhed allene. Fra ingen anden Haand ville
vi modtage det.

See vi da altsaa en højere \emph{Tendents til Selvstændighed} at komme i
Bevægelse og Personligheden at gjøre sig gjeldende, som den, der nu saavidt
mueligt kun vil støtte sig paa sig selv og sin egen Randsagen i al sin
Erkjendelse, saa see vi dog, at det kun skeer, for at det Sandes Magt ret kan
blive fuldkommen hos os. Ville vi end selv prøve alt, ja faaer det vel endog det
Udseende, som om vi krævede det Sande selv til Regnskab, som om vi selv vilde
raade for, hvad der skal være os antageligt og kun vilde lade Sandheden gjelde,
naar vi selv fik stadfæstet den, saa er det dog ikke os eller vort eget, vi deri
søge, men kun det, at Sandheden selv i dens Sandhed fuldeligen maa komme til os.
Søge vi Stadfæstelse, saa er det, at det Sande ogsaa for vore Øjne maa stadfæste
sig selv. Kræve vi til Regnskab, saa er det, vi kræve til Regnskab, kun Former og
Forestillinger, som byde sig frem for os, men som vi ej strax ville antage paa Tro
og Love, og vi gjøre det, for at lade os give ved Sandheden selv det, vi skulle
antage. Det er her, hvor der endnu kun er Tale om rationes cognoscendi, os selv
vi kræve til Regnskab. Rive vi os løs fra enhver anden Autoritet, hvoraf vi forhen
maaskee lededes, saa er det for desto mere at \emph{underordne os den højeste
Autoritet}, det er Sandhedens egen. Ligesom det ellers hedder, at kun den skal
finde Livet, som veed retteligen at tabe Livet for Herrens Skyld, saaledes kunne
vi her sige, at vi her netop kun søge os selv eller søge at vinde os selv, for
desto fuldkomnere at kunne give os selv hen og tabe os selv.

2.~Vi see her ogsaa det Punct, paa hvilket det egentlige videnskabelige
\emph{Studium}, der først ret begynder med det \emph{akademiske Studium}, hæver
sig frem, og hvor det indlyser, hvad det har at betyde. Netop i den Alder, hvori
legemligen og organiskt en højere Drivt til Selvstændighed eller Selvbestaaen
rører sig, og hvori Overgangen fra Barnets Liv til den modnede og fuldendt
Menneskeexistents skal skee, vaagner og virker ogsaa i videnskabelig Henseende
Drivten til aandelig Selvbestaaen i friske frodige Bevægelser. For den Studerende
er det det Trin i Livet, da han skal begynde at tragte efter at erkjende
Sandheden i dens Sandhed. Derfor skal han nu gaae til de Kilder, hvoraf hans
Lærere selv øsede, og vide alt i Kraft af de Grunde og Overvejelser, ved hvilke
Granskerne for ham, af hvilke han lærte og naturligviis bliver ved at lære, kom
til deres Erkjendelser. Det er den videnskabelige Fuldmyndigheds Trin, som nu af
ham skal opnaaes. Han skal hæve sig til sin Tidsalders Højde i sin Videnskab, for
derfra, om Anlægget dertil er ham givet, at føre den videre, være det nu ad en
theoretisk eller en practisk Vej. Derfor skal han nu ei blot lære, modtagende,
hvad der overleveres ham som Granskningers Resultat, men \emph{Selvtænkning} skal
indtræde, egen Gransken og Grunden: han skal studere. I Skolen var den blotte
\emph{Læren}, om end vistnok en tænksom Læren, Hovedsagen. Tilfredsstille sine
Lærere, og kun forsaavidt ogsaa sig selv, skal og vil Disciplen. Nu gjelder det
derimod fornemmeligen at tilfredsstille Sandhedsfølelsens umiddelbart fornumne
Fordringer. Bestræbelsen skal nu være overalt at trænge hen til de \emph{sidste
og højeste Grunde}; mod disse skal Studiets hele Retning være stilet hen; og kan
ogsaa Granskningen endnu ej med Bestemthed og Klarhed fra et givet Punct af og
i en vis Sphære naae saavidt, at den seer det nærværende i bestemt Forbindelse
med alle Tings sidste Grunde, at den kan forfølge Gjenstandenes Sammenhæng og
deres Kjæde saa vidt som mueligt, saa skulle de dog altid haves for Øje, som
det, om end fjærne, Maal, hvortil der skal stræbes, og en Følelse for
Vigtigheden og Væsentligheden af engang at naae til dem, skal bevares og være
levende. Dette er det, der charakteriserer den akademiske Borger, og den ved det
akademiske Studium dannede Videnskabsmand eller Lærde, som saadan. Hvorfor og det
philosophiske Studium, som en Grunden over alle Tings sidste Grunde, og en ved en
saadan Grunden fremkaldet og bevæget videnskabelig Stræben, med det ejendommelige
Præg, en saadan Grunden maa give den, et Studium, der først paa dette Livets Trin
begyndes, nu ogsaa bliver her af saa afgjørende Væsentlighed. --- Imidlertid
udfordres der vistnok ikke, for at et Studium skal kunne kaldes ægte
videnskabeligt, at Tendentsen skal være at erkjende det givne i sin Sammenhæng
med alle Tings højeste Grunde, eller at Formaalet skal være egentligen
philosophiskt. Om ogsaa Granskeren holder sig inden visse Grændser og allene
eftersporer Sammenhæng og Forhold inden en vis Sphære, vil dog Granskningen
kunne være ægte videnskabelig, naar han kun gaaer ud paa indenfor denne Sphære
at gaae saavidt, han kan, baade ved at trænge her saa vidt ind til Sandheden som
mueligt og lade den lyse saa vidt mueligt i sit eget Lys, saa og, deels for at
vinde den, deels for at prøve det fundne og bestyrke det, om det synes at holde
Prøve, at udstrække sin Eftergransken til Randsagelsen af alle de Data, der ligge
i den engang fixerede Sphære. Ogsaa maa Sandsen for alle de fjærnere
Betragtninger, til hvilke der kan være at gaae tilbage, holdes aaben og vaagen
tilligemed Beredvilligheden til at gaae ind i dem, forsaavidt kun det nærværende
Arbejde derved ej vil afbrydes, men derved vil kunne fremmes. Et bestandigt
almindeligt Hensyn til dem maa vise, at i det laveste, nærmeste og specielleste
dog et Bidrag til det højeste, fjærneste, almeneste tilsigtes. --- Til en bestemt
\emph{Begrændsning i Bestræbelsen} kunne iøvrigt ej blot personlige Hensyn og
Grunde føre, forsaavidt ethvert Individ maa mere eller mindre begrændse sig for
at udrette noget med Kraft og Virkning, men ogsaa objective Grunde kunne fordre
den, forsaavidt enhver Gjenstand bør undersøges ej blot i dens Forbindelse med
det Hele og dens derved bestemte Betydning, men og i sin særlige Fremtræden for
sig, hvorledes den i og for sig viser sig; hvorved den da først betragtes som Led
af et vist mindre Hele, hvilket undersøges for sig, førend dette Hele og den selv
tilligemed det betragtes i sine større Forhold. Den philosophiske Aand viser sig
herunder i den Retning, Studiet har til det yderste og sidste, der indenfor denne
Sphære frembyder sig, tilmed i Henarbejdelsen til det Punct, hvor nu en højere
Betragtning kan tage fat.

3.~Hvorledes igjen den her betragtede Bestræbelse efter at erkjende Sandheden i
dens egen Sandhed og lade sig den give ved denne allene, fører til den i det
ontologiske Beviis for Guds Tilværelse i Bevægelse værende Tanke, vil maaskee
letteligen bemærkes; men at gaae ind i alle de Betragtninger, som derved
fremkaldes, dertil er Stedet ikke her.

\begin{center}\large\S.~7.\end{center}

Forsaavidt en Erkjendelse af Sandheden i dens Sandhed opnaaes, opstaaer en
\emph{intelligent Erkjenden} eller en Erkjenden med Indsigt og i Kraft af en
Indseen. Denne, hvorved det erkjendte indlyser ifølge sin Samstemming med det,
hvad Forstand eller Fornuft fordre eller dog tilsige eller tilstæde, og nu antages
ifølge en Forstands- eller Fornuftnødvendighed til at erkjende det, træder nu i
Modsætning til den Erkjenden af det givne og forefundne, hvorved dette antages
ifølge dets blotte Givethed eller Forefundenhed allene, en Erkjenden, der som
saadan har en empirisk Charakteer.

Allerede en Grund kan kun erkjendes som Grund og en Følge som Følge, idet hiin
sees som den, der nødvendigviis fører til sin Følge, og denne sees som den, til
hvilken dens Grund nødvendigviis fører, saa at Følgen viser sig som den, der under
Forudsætning af Grunden ikke kan tænkes anderledes, at altsaa Nødvendigheden til
at antage den viser sig som en Nødvendighed for Tænkningen. Hvad der antages i
Kraft af Grunde, antages altid ifølge noget, Forstanden kan prøve og har prøvet.
Men ligesom saaledes den Nexus, der er imellem Grund og Følge, dennes Grundethed
i hiin, blot kan skues ved en intelligent Skuen, saaledes kan ogsaa, naar det
Sande ej blot skal modtages, men skal modtages som det Sande og erkjendes, fordi
det viser sig i sin Sandhed, den blotte Forefundenhed ej være det, hvorpaa
Erkjendelsen kan støtte sig, men den maa støtte sig paa dets Indlysen. Det Sande
sees først da i sin Sandhed, naar det sees som det, vi ej kunne unddrage os for
at tillægge Gyldighed, naar altsaa vor Tænkning saaledes bindes derved, at vi ej
kunne tænke det anderledes; men denne intelligente Nødvendighed indtræder kun,
idet en Samstemming med en Forstandens egen Fordring eller Forventning viser sig.

En ægte og fuldkommen Erkjenden er altsaa en \emph{intelligent Erkjenden}. Ikkun
ved en saadan erkjendes Sandheden i dens Sandhed. Den empiriske Opfatten og
Anskuen maa derfor, naar den skal vorde en ægte Erkjenden, omfattes saaledes af
den intelligente Stræben, at alt, hvad der forefindes som det virkelige, nu ogsaa
skues i sin, i det mindste relative Forstands- eller Fornuftmæssighed. Kun da
fremkommer en egentlig Videnskab. Hvorimod den Kundskab, hvortil vi komme ved den
blotte empiriske (kun i en vidtløftigere Betydning saa kaldede) Erkjenden, om den
end fører til en, vistnok fornøden og vigtig, ordnet og fuldstændig Oversigt over
de empiriskt givne Gjenstande, og en tydelig og nøjagtig Fremstilling af deres
Udseender eller Skikkelser og forefundne Beskaffenheder, dog allene frembyder et
Indbegreb af Phænomener, hvori Tanken først skal bringe Lys og Klarhed, idet den
forklarer dem, d.\ e.\ viser deres Fremkommen og Tilvær som en saadan, Forstanden
selv, være det og kun under visse særdeles, nok saa snævre Betingelser, maatte
fordre eller dog vente. Hvorved de da antage Væsentlighedens Præg og løftes op
over den Tilfældighed, hvormed de viiste sig, saa længe de endnu kun vare optagne
som blot forefundne.

I den intelligente Erkjenden virkes Forvisningen ved en ganske anden og højere
Nødvendighed, end den empiriske Nødvendighed til at antage, hvilken bestemte
Affectioner medføre. Skjønt det opfattende Individ vistnok ogsaa ifølge disse maa
forestille, som det forestiller, saa at dets Productivitet forsaavidt finder sig
aldeles bundet og bestemt, saa kan det dog, saa længe der ikke indsees, at det
forefundne nødvendigviis er, som det er, tænke sig det, d.\ e.\ forestille sig
det, anderledes, og Productiviteten er altsaa forsaavidt ikke fuldkommen bunden.
Den empiriske Opfatten modtager blot, men Betragtningen, som hæver sig op derover,
kan nu og lade det staae ved sit Værd. Under Granskningen efter det Sande i dets
Sandhed fremkaldes derimod Tanken, der ej blot vil modtage simpelhen, men indsee.
Men naar den nu, idet den bevæger sig efter sin egen Grundnatur og sine
Grundbegreber, dog kommer til det, hvad der i Gjenstanden selv fremstiller sig,
bliver den, just fordi den bevægede sig i sin Frihed og dog kun kom til dette
samme ene, bundet med en højere Kraft; saa at den fuldkomne Bundenhed og
Beherskethed indtræder formedelst den intelligente Erkjenden. Netop fordi den frie
Prøvelse og Overvejelse er fremkaldet og gjør sig gjeldende, kommer nu det, der
ifølge samme erkjendes, til at herske med den fuldkomneste Magt. Sandheden kan ej
lyse i sin Sandhed og indlyse ved denne sin Sandhed selv, uden at gribe og
hentage med en højere Nødvendighed.

Saavidt den intelligente Erkjenden gaaer og træder klarligen frem som saadan,
udtaler den sig nu i en \emph{Construction a priori}, under hvilken det erkjendte,
idet det afledes af den Grundtanke, til hvilken den intelligente Skuen er trængt
frem, fremstilles i dets Forstandsnødvendighed. Den Tanke, som bragde Lys ind i
det erkjendte Hele, stilles nu og hen som Erkjendelsens første Bevæger, ved
hvilken da alt bestemmes, som den medfører, at det maa være.

\begin{center}Anmærkning og Tillæg.\end{center}

Vi forfølge her Erkjendelsen, som den stiger opad igjennem en Række af Trin til
det højeste. Først naar vi have os til dette, kan alt foregaaende sees i sin
sande Betydning og i sit fulde Lys. Men vi ville dog strax her paa Vejen gaae ind
i de Spørsmaal, vi efterhaanden møde og overveje, hvad de føre os hen til.

At en Intelligents oplader sig i Erkjendelsen, og at denne overalt stræber frem
til \emph{Intelligenthed}, sees allerede snart, naar Bestræbelsen efter
fuldkommen Erkjendelse kommer i Bevægelse og bliver sig selv klar. Men hvorfor er
det saaledes? \emph{hvorfor maa al ægte Erkjenden være intelligent?}%
% FLAG p.59: emphasis span (Sperrung) — question set letterspaced; extent approximate. Verify.

Fremdeles: \emph{hvorledes kan dog en udenfra given Gjenstands Natur og
Beskaffenhed erkjendes ved en Erkjendelse, hvis levende Grundnormer, hvis
bestemmende første Bevæger og bestemmende Maae, endog paa ethvert Punct, er af
indre Udspring},%
% FLAG p.59: emphasis span (Sperrung) extent approximate; also "Maae" printed thus (poss. "Maade"). Verify.
 den i os talende Forstands og Fornufts eget?

1.~At dette sidste --- vi ville tage det først i Betragtning --- er saa, er
klart. Ethvert Exempel paa en Erkjendelse, som gaaer frem af en Betænkning, eller
hvorved Sammenhæng, Forhold, Grund og Følge komme i Betragtning, kan tjene til at
vise det. Allerede naar i den simpleste lavere Sandsning Forestillingen fastsættes
paa en vis bestemt Maade, være det endog ganske overeensstemmende med det
umiddelbare momentane Udseende, ja naar der endog kun er Tale om at fixere dette,
og dette skeer med Bevidsthed, saa at en Betænkning opstaaer, eller dog et
Spørsmaal om en Betænkens Fornødenhed eller Ikke-Fornødenhed, vil strax, om end ej
en Forskjel imellem et Skin og det virkelige antages, men Mueligheden af et
Sandsebedrag nægtes, noget intelligent vise sig, forsaavidt noget derved altid
bestemmes ifølge en Slutning, hvis Gyldighed og Consequents, ifølge Grunde, hvis
begrundende Magt, ikkun kan sees med Forstandens Øje. Det er overhovedet kun en
Erkjendelse formedelst Tænkning, der bestemmer Sandseanskuelsen dens Realitet,
hvilket baade gjelder, naar der indenfor Sandsernes Sphære spørges, om og hvorvidt
noget er Skin eller Virkelighed, saa og, naar, foranlediget ved Betragtninger af
en ganske anden Art, det Spørsmaal opstaaer, om den hele Sandseanskuen overhovedet
har nogen Realitet; vi modtage forsaavidt den hele Sandseerkjenden selv, hvad den
Betydning angaaer, den faaer for os, fra den intelligente Erkjendelses Haand.

Den intelligente og den blot empiriske Erkjenden gribe da allerede under den
lavere Sandseanskuelses Dannelse i hinanden, og hiin strækker sig heelt hen
igjennem den Erkjendelse, der med Hensyn til sit Hovedpræg, nemlig med Hensyn til
det, hvad der antages, endnu kun beroer paa en simpel Opfatten eller Modtagen af
det i Erfaringen givne i dets Givethed. Det samme sees da og, naar nu formedelst
en intellectuel Anskuelse de Grundformer erkjendes, der i Gjenstandene opdages.
Uagtet ogsaa disse maae tages, som de forefindes, saa indtræder ogsaa her den
intelligente Erkjenden; ej blot, at vi, naar Sammenhængen ikke er aabenbar, ved
Grunde maae ledes til at bestemme Forestillingen om den paa en vis bestemt Maade;
vi kunne ogsaa overhovedet ikkun see en Grundform som Grundform ved at indsee,
hvorledes ifølge samme den ydre Skikkelse eller Tilsyneladelse, hvis Grundform den
er, virkelig kommer i Stand. Endnu mere træder det intelligente frem, naar et
Hvorfor besvares, naar altsaa Phænomenernes Forklaring findes; thi et Phænomen
forklares ikkun, naar der gaaes tilbage til visse Grundbestemmelser og
Grundforhold, og Phænomenerne nu sees som nødvendige Følger af disse, Følger,
hvilke Forstanden selv afleder og erkjender som Følger deraf, fordi der indsees,
at af saadanne Grundbestemmelser kun saadanne Phænomener kunne tænkes som Følger
og ingen andre; som naar Tyngdens bekjendte Love, der maaskee paa en reen
empirisk Maade kunde været fundne ved de bekjendte Experimenter, indsees at være
en nødvendig Følge af en saadan uafbrudt tiltrækkende Krafts Natur og
Beskaffenhed, som Tyngden er, der ej kan tænkes at virke efter andre Love.

Gaaes herved nu maaskee end ofte ud fra første Bestemmelser, der blot ere reent
empiriskt opfattede og antagne ifølge deres Forefundenhed, uden endnu at være
førte tilbage til saadanne første Bestemmelser, der indlyse for Forstanden selv,
ja er end den Samstemming med det, hvad Forstanden selv medfører, der viser sig,
idet noget gjøres begribeligt, ej sjeldent højst relativ og betinget, (som naar
endog det fornuftstridige i Menneskers Handlinger gjøres begribeligt ved at sees
i sin Consequents af visse Forudsætninger, af hvilke Forstanden selv maa erkjende
dem at være en naturlig Følge, det er en Følge, hvis Følgen eller Affølge er
forstandsmæssig), saa er dog Forklaringen og Begribningen og den deraf
fremgaaende Erkjenden, saa vidt den som saadan gaaer, altid aldeles intelligent;
indenfor den givne Sphære af Betingelser staaer nu det antagne ej længere som
blot forefundet og tilfældigt. Befindende os midt i Empirien søge vi at komme til
Lys og Klarhed ved at gaae ud fra det nærmeste og opsøge dets Sammenhæng med det
nærmeste, og vi kunne nu ogsaa skue de nærmest sig til hinanden knyttende Led og
Dele i denne deres Sammenhæng, tilmed skue dem i deres Betydning og
Væsentlighed for hinanden og med Hensyn til visse første Led, uden endnu at
have forfulgt denne Indseen og Begriben til dens yderste Grændser. Saaledes kan
en vis Dyreform, eller et vist Mineral, man ud af Rækken vil tage til Exempel,
sees som integrerende Led, tilmed som væsentligen medhenhørende Led, i den hele
Række, ja det hele System af Dyreformer eller af Arter af Mineralier kan blive
udkastet paa en apriorisk Maade, idet disse Formers hele Udviklingsgang, som om
den var Naturens Ideegang under disse Formers Production, fremstilles, om end de
Bestemmelser, hvorfra der derved gaaes ud endnu blot tages som de, der ere
abstraherede af Erfaringen selv ved et stort Overblik over Naturen. Men
naturligviis tragter Videnskaben efter overalt at udvide de Grændser, indenfor
hvilke den intelligente Erkjenden en Tid lang er bleven staaende, og saaledes at
overvinde denne dens Relativitet. Saavidt Gjenstandens sande, ægte, med Indsigt
forbundne Erkjendelse gaaer, saavidt gaaer ogsaa en intelligent Erkjenden, og
hvor denne slipper, slipper og hiin.

Heri gaaer nu Forstanden, idet den gjør sine Conseqventser, Constructioner og
Deductioner, overalt ud fra sine \emph{egne} Grundnormer og Principer. En ikke
af Erfaringen modtaget eller abstraheret, men paa en selvstændig Maade af
Tænkningen udspringende Logik, Mathematik, Metaphysik (metaphysiske
Grundbegreber og Grundsætninger) lægges herved overalt til Grund. Og den
bestemmende Indflydelse heraf strækker sig ind det mindste Detail, hvis særlige
Beskaffenhed, idet den forklares og begribes, just ogsaa afledes og erkjendes
som den, der under de givne Betingelser netop tilsiges saaledes og ikke
anderledes af Forstanden, og kun fattes og begribes, idet den netop i denne sin
særlige Beskaffenhed viser sig som den, hvortil hiin Logik, Mathematik,
Metaphysik og den af disses Grundbestemmelser ledede Construction, eller
Udførelse og Uddannelse af en vis Anskuelse, nødvendigviis leder.

Og dog træder denne Erkjendelse, uagtet dens bestemmende Grundbestemmelser ere af
indre apriorisk Udspring, op med Krav paa objectiv Gyldighed, og kan ikke andet.
Hvad der træder frem som det, Forstanden nødvendigviis fører med sig, træder
nødvendigviis frem med denne Fordring paa Realitet. Det er os aldeles utænkeligt,
at det ikke skulde have den. Kunde en virkelig Strid indtræde imellem Erfaringen
og Forstandens nødvendige Fordringer, saa vilde der intet Forhold kunde være
imellem det tænkende Væsen og Virkeligheden. Dog en saadan Strid vilde vi ikke
nogensinde kunde finde eller statuere. Kommer der noget ud som Resultat af en
Tænkning eller Beregning, hvilket ikke stemmer med Erfaringen, saa søge vi
Grunden til denne Uovereensstemmelse altid deri, at enten vore Forudsætninger
have været urigtige eller at i vore Slutninger er indløbet en Fejl. Aldrig kan
det falde os ind, at Naturens Conseqventser maaskee kunde være andre end
Tænkningens, at Naturens Logik, Mathematik eller metaphysiske Principer kunde
være andre, end den i os virkende og talende Forstands og Fornufts.

Men hvorledes kan dog en saadan af Forstanden allene udspringende Erkjendelse
overalt gyldigen lægges til Grund i Erfaringen og derpaa en heel empirisk
Erkjenden bygges?

Aabenbar kun fordi den Intelligents, der i det betragtende Menneske viser sig
som saadan, gaaer igjennem den hele Virkelighed og ligger til Grund i denne selv;
fordi netop \emph{Intelligentsen selv} er et Udtryk af Naturens \emph{sidste
Love}. Hvorledes skulde og i Mennesket, som et Led med af denne Virkelighed,
noget andet kunne være første og inderste Bevæger, end hvad der er det i den hele
Virkelighed overhovedet? Saa at her ikke er at tænke paa en dobbelt Forstand,
som om vor Forstand mødtes med, og samstemmede med, en anden Forstand, Naturens.
Det er den selv samme Intelligents, der i sin udeelte Eenhed aabenbarer sig i
Naturen og i vor Tænken og Betragten, fremstillende sig, sig selv lig, her og
hisset. Hvad der knytter Forbindelsen og Samstemmingen imellem os eller vor
Tænkning og Virkeligheden, er her, som overalt, noget tredie, der, højere end
begge de forbundne, ligemeget ligger til Grund i dem begge. Derom vidner da
ogsaa, saavidt det ufuldkomne og stykkevise kan vidne om det fuldkomne og totale,
idet det viser hen derpaa som paa Kilden for sin Muelighed, derom, sige vi,
vidner det, at der under den ved Selvtænkning ledede Stræben efter at trænge ind
i Naturen, jo mere den gaaer frem, desto mere virkelig gaaer et Lys op for os
over Naturen. Ogsaa bliver, hvad Forstanden taler, medfører, fordrer, os
ligesaafuldt givet, som det vi komme til paa en aldeles empirisk Maade, om det
end gives os fra en ganske anden Side, paa en ganske anden Maade og med et andet
Præg, end det empiriske; saa at vi idet intelligente ikke mindre have med noget
objectivt og ikke mere have med os selv eller med vort eget at gjøre, end i det
empiriske. Er her paa ethvert Punct en intelligent Afleden og Bestemmen, saa er
Erkjendelsen dog paa ethvert Punct en Modtagen og Skuen, da vi forstandsmæssigen
aflede og bestemme, som vi see, at vi maae det.

2.~Men i Mennesket træder Intelligentsen frem som saadan, eller i \emph{sin
Intelligenthed}. Derfor er det selv samme her ikke det selv samme mere. Derfor
indtræder her hiint \emph{a priori}; hvilket nu viser sig for os i sin reale
Betydning og Væsentlighed. Skal Intelligentsen træde frem i sin Intelligenthed,
skal den Forstandighed, der er i Naturen, aabenbare sig i Mennesket, som den er,
og er den i Naturen det, der \emph{ligger til Grund}, saa maa den og i Mennesket
vise sig som det til Grund liggende, tilmed vise sig reen i og for sig i sin
oprindelige Selvstændighed, som det \emph{constitutive og bestemmende}, og tilmed
som det \emph{aprioriske} i Modsætning til den sig i Virkeligheden fremstillende
Mangfoldighed, og som Princip for en selvstændig Construction. Og saaledes see
vi, hvorledes en udenfra given Gjenstands Natur og Beskaffenhed ikke blot kan,
men ogsaa, naar den skal erkjendes tilfulde, maa erkjendes ved en Erkjenden,
hvis bestemmende Princip og bestemmende Maae er af indre Udspring, Forstandens og
Fornuftens eget.

Det bliver nu ogsaa klart, hvorfor, uagtet det er den selvsamme Intelligents, som
træder frem i os og i Naturen, dog, og det just fordi den er en Intelligents og i
os træder frem i sin Intelligenthed, noget indre eller indenfra sig fremstillende
(subjectivt) kommer frem i Modsætning til det udenfra sig fremstillende
(objective), baade forsaavidt i Almindelighed oprindelige af Tænkningen selv
fremgaaende Fordringer vise sig, med hvilke det givne maa findes at samstemme,
naar det skal vorde begrebet, saa og forsaavidt ogsaa i det Enkelte ethvert Punct
kun erkjendes tilfulde, naar det sees som det, hvortil ogsaa en apriorisk
Afledning eller Deduction kunde føre. I den aprioriske Construction vise sig
begge, det indre og det ydre, i deres største Modsætning, hvorledes de træde lige
overfor hinanden, men og, naar nu det givne just ogsaa viser sig som denne
Constructions nødvendige Resultat, i deres fuldeste Samstemming til Eenhed, en
Samstemming, uden hvilken det ej kommer til en ægte Erkjenden. --- Forsaavidt
denne aprioriske Construction gaaer ud paa at stille en given Virkelighed i sit
fulde, sande Lys, viser den sig som en blot eftergjørende Construction, hvorved
vi, til Vidne og Forvisning om, at vi fuldkomment ere trængt ind i Erkjendelsen
af det givne i dets Sandhed, nu i Kraft af denne Erkjendelse udføre det for os i
et Afbillede, hvilket da aldeles maa samstemme med det, hvad en nøjagtig
Iagttagelse lader os finde og skue i den givne Virkelighed. Men forsaavidt dog
den opnaaede Grunderkjendelse udgjør det bevægende, alt bestemmende Første i denne
Construction, træder denne frem som en selvstændig med den aprioriske Charakteer.

3.~Men ligesom en nøjagtig empirisk Iagttagen maa afgive Prøven paa denne
Constructions Rigtighed og Gyldighed, saaledes kan ogsaa kun en saadan Iagttagen,
et nøjagtigt og omhyggeligt Studium af den givne Virkelighed selv, d.\ e.\ af det
forstandige eller intelligente i den selv, lede til samme. Tanken vil
gjennemtrænge det forefundne, men dertil udfordres, at den Forstandighed, som
deri selv aabenbarer sig, opfattes d.\ e.\ forstaaes, som den er, og det ganske
og aldeles som den er. Det faaer her derfor altid det Udseende, som om vor
Forstand mødtes med en i Tilværelsen sig aabenbarende Forstand. Kun forsaavidt vi
finde Mening og Forstand i Natur og Virkelighed, kunne vi fatte og forstaae den,
d.\ e.\ erkjende den tilfulde; men denne finde vi nu deri og derfor kunne vi
erkjende den saaledes. Naturens Studium er som Studiet af en Bog, i hvis Mening
og i hvis hele Aand vi igjennem dens Bogstav søge at trænge ind, men ogsaa kun
kunne trænge ind, forsaavidt vi finde Mening og Aand deri. Men hvorledes trænge
vi ind deri? Kun idet det, vi deri finde udtalt, indlyser for os selv, og det
saaledes, at vi see, hvorfor det har udtalt sig, som vi finde det udtalt, og at
vi nu selv samstemme med deri. Skal et ypperligt, dyb Sandhed fremstillende Værk
forstaaes af os tilfulde, maa det Lys, der lyste for dets Forfatter, gaae op for
os, og vi komme i Besiddelse af det bevægende første Princip, hvorfra alle hans
Bestemmelser med Nødvendighed gik ud, saa at vi nu med ham tage Sandheden fra
første Haand.

Men netop saaledes har det sig med Naturens Studium; hvad der i den selv
bestemmer Tingene deres Betydning, maa af os fattes saaledes, at der deraf gaaer
et Lys op for os over Naturen, og ogsaa hvad den udtaler, maa vi, trængende
gjennem dens Bogstav ind til dens Aand, søge at tage fra første Haand. Hvoraf da
noksom sees, at Hensigten ej kan være ved en fra Erfaringen sig løsrivende
aldeles paa egen Haand gaaende apriorisk Construeren at gjøre den empiriske
Iagttagen og Eftergranskning overflødig; men netop igjennem denne, som igjennem
et nøjagtigt Studium af Ordene og Sætningerne i Naturens Bog, at trænge ind til
dens Mening og Aand.

Forholdet imellem \emph{Empirisme og Rationalisme}, en Modsætning, der især i
visse Videnskaber, f.~Ex i Medicinen, har vakt saa megen Splid og Kamp, maa nu
vel heraf være klar. At nemlig --- ligesom vi, hvad en anden Adskillelse angaaer,
kunne sige: kun Fornuften gjør tilgavns forstandig, men ogsaa kun Forstanden gjør
tilgavns fornuftig, --- at saaledes ogsaa vistnok kun et omhyggeligt og nøjagtigt
Studium af det empiriskt givne kan føre til en fuldkommen rational Erkjenden, men
at dog ogsaa kun en ægte rational Stræben fører til en solid Empirie. Er det
Empirie eller Erfaring at vide, at dette eller hiint anvendte Middel i
Almindelighed har disse eller hine Virkninger, uden at indsee, hvorledes dog og
hvorfor just af et saadant Middel en saadan Virkning kommer til at blive Følgen,
saa er Empirie eller Erfaring ingen sand Erkjenden. Den er da at sætte i Classe
med den Kundskab, der blot er et Indbegreb af visse Resultater og lærte
Sætninger, uden at man er i Besiddelse af de Grunde og den Tænkning, der førte til
disse Resultater, uden at man seer, hvad der ligger i de lærte Sætninger og
egentligen er udsagt i dem.

4.~Formedelst den frie Tænkning skeer en højere og inderligere Tilegnelse af den
til at erkjendes sig frembydende Erfaringssphære. Den opfattende Personlighed,
der derved ligesom assimilerer sig den og optager den til inderlig Eenhed med sin
egen (fornuftige) Leven, bliver derved ligesom hjemme deri og raadig derover, saa
at den bevæger sig og gaaer omkring deri som i sin egen Ejendom, hvori intet mere
er den fremmed. Men i og ved den samme Sigtilegnen bliver denne Personlighed mere
og mere taget i Besiddelse af den erkjendte Sandhed, saa at den personlige Leven
gaaer op i denne Sandheds egen Sigfremstillen, og mere og mere røres og bevæger
sig i Sandheden og efter Sandheden.

\begin{center}\large\S.~8.\end{center}

Bestræbelsen efter at erkjende Sandheden i dens Sandhed fører, idet det viser
sig, at kun den Erkjenden, hvorved det givne skues i sin Samstemming med
Forstandens Fordringer, er en ægte Erkjenden, og idet Spørsmaalet nu overalt
bliver, om det forefundne ogsaa maa antages nødvendigviis at være, som det er,
til at overføre Granskningen efter Grunde til Gjenstandene selv, og at spørge,
hvorfor disse selv ere, som de findes at være. Søgtes forhen Grunde, som skulde
bestemme vor Antagen, Erkjendelsesgrunde, (rationes cognoscendi), saa opstaaer nu
en Gransken af ganske anden Art, idet Tingenes \emph{Tilværelsesgrunde} (rationes
essendi) opsøges.

Først dermed indtræder den intelligente Erkjenden. Med den fuldkomneste Indsigt i
Nødvendigheden af at antage, hvad der antages, kan dog selve dette antagne endnu
staae for os som noget blot forefundet og tilfældigt, om hvis Væsentlighed der
endnu først bliver Spørsmaal. Ikke dette antagne selv, kun dets Antagelse, hvis
Tilværelsesgrunde nu Erkjendelsesgrundene vise sig at være, omfattes da af den
intelligente Erkjenden. Denne er altid en Erkjenden af en \emph{real
Nødvendighed}. Men skal denne reale Nødvendighed ikke blot indlyse, men og
indlyse i sin Sandhed og tilfulde, saa maa den ej blot indsees, men og indsees i
Kraft af de reale Grundes Erkjendelse.

Kun da, naar vi skue Tingene i deres egen reale Nødvendighed, kunne vi ej
undslaae os for at tillægge dem Væsenhed og Gyldighed, og da see vi ogsaa,
hvorvidt og i hvilke Relationer den er at tillægge dem, nemlig saavidt denne
Nødvendighed viser sig, og i de Relationer, hvori den viser sig. Kun da skue vi
altsaa Tingene i deres Sandhed. Men i en reel Virkelighed, hvor Tingene staae som
Led og Dele af det virkelige Hele, maa den reale Nødvendighed, eller
Nødvendigheden i og for Tingene selv, vise sig som den, ifølge hvilken disse Ting
virkelig komme eller ere tilstæde, og haandhæve deres Plads. Og saaledes maa da
deres Nødvendighed vise sig som en af visse Tilværelsesgrunde udspringende
Nødvendighed, hvilke Tilværelsesgrunde altsaa maae eftergranskes. Skal
Erkjendelsen fuldkomment svare til det erkjendte, skal den Productivitet, hvorved
formedelst Forestillingens Fuldendelse Erkjendelsen selv først kommer i Stand,
fremstille det erkjendte aldeles, som det er, saa maa, hvad der er det første og
bestemmende for Gjenstanden, ogsaa staae som saadant i Erkjendelsen og
Forestillingen. Ere Tingene det, de ere, ifølge et dem saaledes bestemmende
Noget, der som deres reale Grund constituerer dem, saa vilde Erkjendelsen være
ufuldkommen, naar de nu ikke ogsaa skuedes som de, der just ifølge dette Noget
ere, som de ere, eller naar nu ikke ogsaa, tilligemed dem selv, dette Noget
erkjendtes som det, hvorved de saaledes constitueres.

\begin{center}Anmærkning og Tillæg.\end{center}

1.~Erkjendelsen af Tingenes rationes essendi lader os nu ogsaa først, idet vi see
disse at bestemme Tingenes Tilværelsesmaader, skue dem i deres mere eller mindre
\emph{relative Gyldighed} eller i deres større eller ringere Væsenhed. De hjemle
ligesom deres Tilvær og deres Afkomst til at tillægges en vis Betydenhed, idet
deres reale Grunde vise sig; men nu sees det og, at denne deres Betydenhed og
Væsenhed overhovedet kun er en betinget. Videre end det, hvorpaa deres egen
Realitet støtter sig, gaaer, kan denne selv ikke gaae. Og ligesom det først er
den tænkende Betragtning, for hvilken Spørsmaalet om Realitet overhovedet opstaaer
og det er et intelligent Grundbegreb, som derved er i Bevægelse, saaledes er det
og Intelligentsen, der bestemmer Forskjellen baade imellem Skin og Virkelighed
saa og imellem meer og mindre gyldig Realitet, og nu derefter anviser alt givet
sin Plads i denne Henseende. Der er intet, som jo har sin Realitet, men denne
sees først, naar det sees, hvorledes, tilmed hvorvidt, Tingen ej kan andet end
være til, naar den altsaa føres tilbage til sin reale Grund. Endog Skinnet har
sin Realitet, nemlig som Skin, som det tilsyneladende, og har en større eller
ringere Realitet, jo mere eller mindre et saadant Skin viser sig som nødvendigt,
jo mere udvidede de Relationer ere, der maae overstiges for at slippe det; men
det sees først som et saadant Skin og nu ogsaa som et forsaavidt nødvendigt Skin,
naar der indsees, hvorfor og hvorledes det nødvendigviis opstaaer ifølge det, der
som dets Grund fører det med sig.

2.~Man seer let, at det er noget ganske andet at undersøge, hvad vi skulle antage
at finde Sted, og tilmed at fremstille de Erkjendelsesgrunde, som dertil maae
bestemme os, end det er at eftergranske de reale Grunde, hvorfor det er saaledes,
som vi finde det at være. Den ene Undersøgelse og Indsigt ligger ofte i en ganske
anden Sphære end den anden. Ved den første fastsættes først det Phænomen eller
den Beskaffenhed og de Forhold, hvis Hvorfor ved den anden skal besvares. Men
ligesom det er vigtigt først at vide, hvad der maa antages som givet, inden dets
rationes essendi opsøges, saaledes kaste disses Erkjendelse ofte igjen et saadant
Lys paa hiin første Betragtnings Gjenstand, at \emph{Antagelsens} egen Vished og
Tilforladelighed bestyrkes og fuldendes ved \emph{Indsigt i de reale Grunde}. See
vi, at et vist Led eller et vist Punct i et givet Hele, ifølge den i dette Hele
erkjendte eller antagne Sammenhæng, ikke kan andet end være saa og saa, og
stemmer det, denne Forstandsmæssighed fordrer, sammen med det, hvad Iagttagelser
og Slutninger fra Iagttagelse eller andre Erkjendelsesgrunde lede os til at antage
om dette Led eller enkelte Punct, saa kan Tilliden til de sidstes Rigtighed ikke
andet end vinde derved, ja stige til uomstødelig Vished.

3.~Paa tilstrækkelige Erkjendelsesgrundes Fremstilling gaae Beviser ud, der føres
ved \emph{Slutninger}. Slutningen er overhovedet Formen for denne Art af
Erkjenden. Den fører ifølge sin Form i Almindelighed kun til at see
Nødvendigheden af at antage dette eller hiint. Dette falder ikke alene i Øinene i
saa mange Tilfælde, hvori Følgerne eller Virkningerne af en Gjenstands Grundform
eller Grundbeskaffenhed og Grundyttringer blive det, hvorfra vi slutte os til den
selv, og hvoraf vi altsaa udlede Beviset for den derom dannede Forestillings
Sandhed; som naar det af de bekjendte Grunde godtgjøres, at Jorden er sphæroidisk
og at den drejer sig om sin Axe og gaaer omkring Solen, hvorved det Spørsmaal,
hvorfor den har en saadan Figur og Bevægelse, endnu aldeles ikke besvares, ja end
ej engang berøres, da det ligger i en ganske anden Undersøgelses Sphære. Men
ogsaa, hvor der ikke saaledes fra visse Phænomeners Beskaffenhed sluttes tilbage
til en vis factisk Grundtilstand eller Grundform, eller et vist Grundfactum, men
hvor en, fra det i sig selv første til det deraf følgende lige fremad skridende,
Tankegang fører til en Række af Erkjendelser ifølge strenge Slutninger, og til en
heel videnskabelig Lærebygning, endog paa en aldeles intelligent Maade, kan
Erkjendelsens Art endnu være den, at den blot, ifølge gyldige Erkjendelsesgrunde,
viser, hvorledes Gjenstandene ere beskafne, hvilke deres Forhold ere, uden at
Tilværelsesgrundene dertil udgrundes og drages frem.

\emph{Mathematiken} er, i det mindste efter den sædvanlige Fremstillingsmaade, et
stort og talende Exempel herpaa. Den hele Beviisførelse giver her i Almindelighed
kun Erkjendelsesgrunde, og andre søges eller fordres ej. Nødvendigheden, med
hvilken og ifølge hvilken der erkjendes, er vistnok intelligent, og den betræffer,
hvilket ikke maa oversees, selve Gjenstandene, selv de mathematiske Former eller
Anskuelser; der indsees virkelig, at disse Gjenstande selv nødvendigviis have de
erkjendte Beskaffenheder og Forhold, saa at Nødvendigheden ikke er en blot ydre,
os bestemmende, men en indre Nødvendighed i og for Gjenstandene (Trianglerne,
Qvadraterne o.\ s.\ fr.) selv. Men bestandigen er dog det, hvis intelligente
Nødvendighed indsees, ikkun, at Beskaffenheder og Forhold ere, som de ere; men
hvad Grundene angaaer, da vises kun, hvorfor vi maae \emph{erkjende}, at det er
saa, ikke, hvorfor det \emph{i sig selv} er saa. I de Hjelpelinier, der saa ofte
construeres, overhovedet i de Biconstructioner, der saa ofte gjøres, for Bevisets
Skyld, og i andre Omveje, som bruges, kunne dog de betragtede Beskaffenheders og
Forholds virkelige Tilværelsesgrunde ej være at søge. Kun paa Formers og Forholds
Bestemmelse, ej paa deres indre Væsen og Sammenhæng med et højere og højeste
Værende, gaaer overhovedet Mathematiken ud. Den mathematiske \emph{Demonstration}
er endnu ingen philosophisk \emph{Deduction}. De Beviser, ifølge hvilke der
erkjendes, med den Maade, hvorpaa der erkjendes i Kraft af dem, ere, med Hensyn
til det erkjendte selv i sig, intet i Sandhed væsentligt. Beviset giver en vissere
Erkjenden og fører til en mere sammenhængende, men giver ingen højere og
væsentligere Erkjenden. Forholdenes og Beskaffenhedernes Kilde opvises ej, de
udledes ej af deres Udspring eller Herkomst; (deres Hjemmel er ikke en sand
Retsgrund, men snarere et Slags causa in facto posita). Der erkjendes derfor ogsaa
stykkeviis, ej i en altomfattende Sammenhæng, hvori alt skuedes i gjensidig
væsentlig Forbindelse. Denne Videnskab har ligesaa lidet ægte fuldkommen
Videnskabs Form, som den selv er en ægte og fuldkommen Viden fra Grunden af. At
den mathematiske Methode ej bør optages i Philosophien, dertil er (for at vi i
forbigaaende skulle bemærke det) den stærkeste Grund den, at denne Methode
ingenlunde er fuldkommen nok dertil, ingenlunde i Sandhed grundig (til Grunden
gaaende) og ægte videnskabelig nok. Til at optage den philosophiske Gehalt skal
en ganske anderledes systematisk og organisk Form, med en ganske anderledes nøje
sammenknyttende Methode, end i det mindste den euklidiske og enhver anden, der
gaaer frem i hans Aand, saa meget end dens Strenghed og Nøjagtighed ellers er at
prise. Saa lidet det, uagtet al Erkjendelse af saa mange Chrystallisationsformers
Skjønhed og Betydningsfuldhed, vilde være passende at tvinge et Træ ind i en
saadan Form, saa lidet passer det sig, om man end fuldkomment erkjender det i sig
fuldendte og elegante i en ægte mathematisk Methode, at ville tvinge den
philosophiske Stræben, der tragter efter en friere, livfuldere og mere organisk
Form, til at gaae ind i den mathematiske, naar den skal udtale sig i bestemt
Fremstilling. I den organiske Fremtrædelsesmaade, hvortil den stræber, maa det
sidste ligesaavel bære det første, som det første begrunder det sidste, idet hele
Erkjendelsen hviler i det, der som herskende Sjæl gaaer igjennem det Hele, hvoraf
følger, at ingen enkelt bestemt Sætning kan være fuldkommen evident i sig, uden
alt Hensyn til de følgende og til det hele videnskabelige Indbegreb af Sætninger;
hvorimod i Mathematiken, der mere gaaer frem i lige Linier, end i en Cyklus, hver
tidligere Sætning er vis nok uden at trænge til den efterfølgende, og de første
Sætningers Evidents ikke hviler i det hele Systems Udførelse, noget, som derimod i
Philosophien er Tilfældet, hvilket i det følgende --- thi her anticipere vi det
--- vil blive klart. Det kan endog være Spørsmaal om det, Mathematiken efter dens
sædvanlige Fremgangsmaade giver os, kan kaldes et videnskabeligt System i
egentlig Forstand; i det mindste sees de mange, saa ofte uafhængigt af, og uden
Hensyn til, hinanden evidente Sætninger, ikke i nogen saadan væsentlig Sammenhæng
med hinanden, at de føres tilbage til en Eenhed og at de vise sig som dem, der
tilsammen danne et i sig afrundet Hele.

[Mathematiken i det Hele bliver overhovedet fra Philosophiens Standpunct, ligesom
Empirien, at betragte som et stort Phænomen (ligesom en vis særlig Forstandens
Verden), hvis Muelighed og sidste Grunde blive en Opgave for Philosophien. En
Lære om blotte Former og Forhold kan da ej heller endnu være en fuldkommen
Videnskab, og kan derfor ikke heller unddrages den philosophiske Stræben, som om
den var tilstrækkelig fuldendt i sig og sig selv nok. Den ægte videnskabelige
Mathematiker vil ogsaa stræbe at bringe en philosophisk Aand til at bevæge og
røre sig i denne Erkjendelsens Sphære.

Men spørge vi, hvorledes den philosophiske Stræben kan vinde Rum i denne
Videnskab, og hvorledes i Særdeleshed en fuldkommen Erkjendelse af reale Grunde
og i Kraft af reale Grundes Erkjendelse ogsaa heri kan opnaaes, da er dette
vistnok ikke saa let at sige: ingen Videnskab er for Tiden fjærnere fra det, der
kunde kaldes dens Philosophie, end Mathematiken. Tvende Hovedpuncter kan man
imidlertid i denne Henseende vise hen paa. Vi have allerede --- dette er det
\emph{første} --- bemærket, at, skjønt Beviserne kun give Erkjendelsesgrunde, saa
blive dog de beviiste Beskaffenheder og Forhold erkjendte som saadanne, der hos
deres Gjenstande ej kunne være anderledes, saa at deres reale Nødvendighed
indsees. Og denne indsees, fordi der indsees, at i ethvert mueligt Tilfælde vil
under slige Forudsætninger det samme kunne bevises, da Gjenstandens (f.~Ex.\
Triangelens) Grundform altid vil tilstæde det; men idet denne Gjenstandens
Grundform viser sig som den samme Beviismaades bestandige reale Muelighedsgrund,
sees det, at den egentligen er de beviiste Beskaffenheders og Forholds
Tilværelsesgrund. Det kommer altsaa kun an paa, at denne egentlige
Tilværelsesgrund ogsaa bliver seet som saadan, og at det, der skal bevises, sees
som \emph{dennes} Følge, og det som dennes ligefremme ikke desuden paa visse
Hjelpeconstructioner beroende Følge, eller rettere at det sees som medhenhørende
til denne Grundform selv og til dens væsentlige Udtryk. En saadan Erkjendelses- og
Beviismaade gik det Udkast til en Geometrie efter en genetisk Methode ud paa,
hvilket Ørsted for en Deel Aar siden begyndte at udarbejde og meddele til
Adskillige, og hvilket han, som vi ønske og haabe, vel engang finder Lejlighed og
Lyst til at lade komme for Lyset. Den Indvending, et saadant Forsøg vel snarest og
først vil møde, nemlig, at Beviserne paa denne Maade ikke kunde faae den
euklidiske Stringents, men at, naar nu Figurer skulle betragtes i deres Opstaaen
ved Liniers Bevægelser, vel endog det crux mathematicorum, der i den højere
Mathematik giver den efter Stringents overalt stræbende Mathematiker, saa meget at
skaffe, Begreberne om det Uendeliglille og Uendeligstore, nu kunde faae Indpas
ogsaa i den lavere Mathematik: denne Indvending vil i Philosophens Øjne betyde saa
lidet, at det befrygtede for ham netop vil være det glædelige derved; thi, at just
der, hvor det saakaldte irrationale i Mathematiken indfinder sig, netop det
rationale kommer til Syne, er en, nu og allerede for længe siden, ikke mere ny
Bemærkning. (See herved den baade dybt og skarpt seende Hegels Logik første Bind,
Anmærkningen pag.\ 206).

Men om nu end saaledes en mathematisk Constructions Grundnorm, der tillige er
Grundformen i den derved frembragte mathematiske Anskuelse, overalt haves saaledes
frem, at den overalt sees som Grunden til de Beskaffenheder og Forhold, der
erkjendes at finde Sted i den construerede Anskuelse, og fremdeles deraf en vis
nødvendig Gang i Sætningernes Følge efter hinanden udvikler sig, saa bliver dog
endnu Spørsmaal om selve de forskjellige Constructioners (selve Trianglers og
Qvadraters og andre Formers) reale Tilværelsesgrunde og deres reale Betydning i
Erkjendelsens Hele. Realbegreber (altsaa egentlige og sande Begreber) om dem
giver Mathematiken endnu ikke. Skulde imidlertid dog ikke saadanne ogsaa være at
finde? men hvor og hvorledes skulle vi søge dem?

For at see endelige Former i deres reale Betydning og Væsentlighed, maae vi ---
hvilket er \emph{det andet}, vi vilde vise hen paa --- see dem som integrerende
Led i den reelle Virkelighed; kun som Led og Dele af denne, eller som Former, der
hver paa deres Sted nødvendigviis fremkomme i denne Virkelighed, sees de i en
saadan Sammenhæng med det Hele, at nu denne ogsaa vil kunne forfølges saavidt,
til de sees i deres Sammenhæng med alle Tings sidste Grunde. Saaledes ville de
for det første allerede blive betragtede, naar de sees som Constructioner, til
hvilke den reelle anskuende og productive Activitet, hvorved de frembringes, i sin
nødvendige Udviklingsgang nødvendigviis maa komme. Den Orden og Følge, hvori
Sætningerne skulle foredrages, og det Lys, hvori de derved skulle stilles, vil da
især blive Gjenstand for Overvejelse; men da vil ogsaa allerede deri noget
philosophiskt vise sig i Mathematiken, der nu og vil kunne blive Gjenstand for
Debat, som og for dialectisk Betragtning, noget, Mathematiken ifølge sin
apodictiske men ikke philosophiske Erkjendelsesmaade ellers ej kjender til. Men
vigtigere, end dette blot subjective Hensyn, vil det være, at opsøge dem i
\emph{Naturen}, og naar man der finder mathematiske Former, være det som
faststaaende Skikkelser, f.~Ex.\ hos Chrystallerne, eller som Bevægelsesformer,
f.~Ex.\ i Planeternes Gang, at eftergranske, hvorfor i Naturen netop \emph{paa
saadanne Puncter saadanne Former} fremkomme. De mathematiske Former findes her
\emph{som givne Gjenstande i Virkeligheden}, altsaa i en vis nødvendig
Forbindelse med den hele givne Natur, og nu kan Bestræbelsen efter, overalt i
Naturen at trænge ind til de sidste Grunde, ikke andet end omfatte dem ogsaa.
Netop fordi de i den rene Mathematik ikke bleve betragtede som givne, opstod
intet Spørsmaal om deres rationes essendi. I en saadan Philosopheren over de
mathematiske Former vilde det nu vel ogsaa vise sig, at en mathematisk
Constructions reale Betydning kan være forskjellig i forskjellig Henseende;
hvoraf da vel ogsaa en Uvished i Henseende til dens rette Bestemmelse, eller i al
Fald en Forskjellighed i Henseende til Gangen i, og Maaden for, dens Bestemmelse
vilde kunne blive Følgen og den philosophiske Forskjellighed i Anskuelsesmaade og
deraf følgende Strid vel igjen komme frem i denne, for slige Anfægtelser ellers
befriede Erkjendelsessphære. Men dette vilde dog være til Vidne om, at
Erkjendelsen var bleven organisk.]

3.~Skjønt de Grunde, i Kraft af hvilke der erkjendes i Mathematiken, blot ere%
% FLAG p.85: this point is printed "3." again (a repeat; the prior point was also "3.", p.76). Verbatim.
 Grunde for Antagelsen, saa, da dog det, til hvis Erkjendelse disse Grunde føre,
viser sig som noget, der i de construerede Former eller Anskuelser, tilmed i
Gjenstandene selv, er, som det er, ifølge en objectiv intelligent Nødvendighed,
saa have vi i Mathematiken et Exempel paa en reen intelligent og det en absolut
intelligent Erkjenden. Intet tages her som blot forefundet, men alt findes og
erkjendes som noget, Forstanden fordrer. Forstanden finder sig her aldeles i sin
egen Sphære, hvori den er fuldkommen hjemme, derfor er Visheden her en fuldkommen
Evidents.

Ogsaa Philosophien tragter efter en saadan, men dens langt højere Natur og
Formaal gjør den her langt vanskeligere at opnaae. Medens Mathematiken gaaer ud
paa at udgranske Beskaffenheder og Forhold hos Gjenstande, der ere Producter af
dens egne Constructioner: Constructioner, om hvis væsentlige Nødvendighed der ej
engang er Tale, da deres mathematiske Muelighed her er nok til deres Realitet,
medens Mathematiken derfor ogsaa kun har at udfinde, hvad der af visse saaledes
muelige Bestemmelser er en væsentlig Følge, saa har derimod Philosophien ikke
allene at udfinde og bestemme, hvad der i og for sig er det sande i Henseende til
Gjenstande (physiske eller moralske), der ere virkelig givne, og hvilke maae
tages med de Bestemmelser, hvormed de findes, men den gaaer endog ud paa, ja skal
gaae ud fra, det, der overhovedet er det absolut sande, der ligger til Grund for
alle de Betingelser og Bestemmelser, hvormed hvilkesomhelst Gjenstande kunne vise
sig. Da nu her alle Ting vise sig i uendelig mange Forbindelser og Forhold, og
disse Relationer netop bestemme Tingenes egne Tilværelsesmaader eller
Fremtrædelsesformer, idet alle Ting overhovedet kun have en relativ Væren og
Væsentlighed, saa er her en uendelig Række at gjennemlebe, et uendeligt Indbegreb
af Gjenstande at skue i deres Sammenhæng, inden den absolute Vished kan opnaaes i
Henseende til nogen af dem. Vel lyse nu mange Puncter snart frem som saa evidente
i sig, at om deres Vished ingen Tvivl synes at kunne være, men dog kan netop det,
hvoraf den fuldkomne Vished skulde udspringe, blive det, der ryster deres Vished,
fordi de, naar nu Regnskab fordres for deres Gyldighed, ej synes at kunne
godtgjøres med en Evidents, imod hvilken jo dog en Skepsis kan rejse sig. Her
gjelder det da tilsidst, om \emph{et højeste absolut fast Punct} kan findes,
hvortil alt saaledes kan knyttes, at ved samme dets mere eller mindre relative
Gyldighed med Fasthed kan bestemmes. Dette højeste Punct vil nu snart vise sig
ikke blot at være det højeste, men tillige det Middelpunct, hvorom alt samles til
Eenhed og Gyldighed, og ved hvis Lys ethvert Punct skal belyses, saa at det skues
som det algyldige, der gaaer igjennem alt. (See længere hen §.~10). Men nu spørges
da igjen, hvori Visheden om dette rene Absolutes Realitet og Sandhed hviler. Og
saa faaer Sagen paa den ene Side det Udseende, som om denne Vished kun kunde
støtte sig derpaa, at alt \emph{andet} netop saaledes \emph{viste hen til hiint},
som det første og alt betingende, at, saavist saa meget andet, som nu efter sine
Hovedpuncter kunde gjennemgaaes, var sandt og unægteligt, maatte ogsaa hiint være
sandt og unægteligt. Men paa den anden Side er det klart, at alt andet kun kan
saaledes vise hen derpaa under dets egen Forudsætning, og at det maa have viist
sig i sin Realitet, for at vi kunne see alt andet vise hen derpaa. Men saa maa det
da erkjendes som det, hvis Realitet ikke blot maa hvile i det selv allene, men
hvis Realitet ogsaa maa vise \emph{sig ved sig selv allene}.

Og saa see vi da, hvor naturlig den allerede forhen i Betragtning tagne Tanke er,
som ligger til Grund for det ontologiske Beviis for Guds Tilværelse, en Tanke,
til hvilken en Stræben efter at erkjende Sandheden i dens Sandhed og ved dens egen
Sandhed ikke kunde andet end føre, og til hvilken alle de Overtænkninger, i
hvilke vi her befinde os, næsten paa ethvert Punct vise hen. Ligesom Troen paa
Gud nødvendigviis faaer, som den da fra først af allerede har, Charakteren af en
fast i sig selv hvilende Tro, saaledes stræber Intelligentsen ogsaa efter at
bringe denne Erkjendelse til en ligesaa umiddelbar, i sig selv hvilende Evidents.
Dog, til Udførelsen af disse Betragtninger er Stedet ikke her.

\begin{center}\large\S.~9.\end{center}

Granskningen efter Tingenes Tilværelsesgrunde gaaer nærmest ud paa at udfinde
deres Sammenhæng og Forhold, og deres derved bestemte reale Betydninger, hvoraf
da deres Tilvær og Beskaffenhed begribes og forklares. Men en saadan Erkjendelse
er en Erkjendelse af \emph{et Begreb} og en Erkjendelse \emph{i Kraft af et
Begreb}. Tingenes egentlige sande Begreber (sensu eminenti), nemlig deres
\emph{Realbegreber}, vise sig nu som dem, hvori Tingenes sande Grunde nærmest ere
at søge.

At spørge efter en Tings Grund er nemlig at spørge efter noget \emph{andet} end
den selv, af hvilket den følger, med hvilket altsaa dens Tilvær er givet som
nødvendig. En Gransken efter Grunden til en Gjenstand gaaer altsaa, med mindre
den findes at have sin Grund aldeles i sig selv, eller selv at være det omspurgte
Andet, ud over Tingen. En anden Gjenstand søges da, hvoraf den givne kan ansees
som en Følge. Men da er (ikke at tale om, at denne anden Gjenstand, naar den er
af samme Art og Sphære, som den første, vil vise sig som lige saa tilfældig,
ligesaa langt fra at være nok i sig selv, som den første) det, der nærmest fører
Gjenstanden med sig, tilmed dens nærmeste Grund, egentligen den Forbindelse, den
Nexus imellem hiin og denne, som gjør, at hiin anden kan og maa føre denne med
sig. Hvilket ligefuldt har sin Gyldighed, hvad enten den ene som Virkning følger
af den anden som Aarsag, eller flere Gjenstande i Forening vise sig som dem, der
væsentligen følge tilsammen med hinanden. Altid vil og maa der, naar
nogensomhelst anden Gjenstand anføres som Grund til en omspurgt, spørges, hvorfor
hiin har denne til Følge, hvilket da er et Spørsmaal efter en vis
Sammenhængsmaade imellem dem, ifølge hvilken den ene fører den anden med sig.
Sammenhængen viser sig da, som det nærmest begrundende; selve Forstandsmæssigheden,
forsaavidt den skues som den Forstandsmæssighed og Forstandsnødvendighed, der i
Tingen selv med den ene kræver den anden, er her nærmest den egentlige reale
Grund; kun ved at indsee denne, begribes og forklares Tingen.

Men Sammenhæng erkjendes kun tilfulde, naar den sees som den, formedelst hvilken
et Indbegreb af Dele eller andre Enkeltheder forbindes til fuldelig Eenhed og
Heelhed. En saadan Sammenhæng er Udtrykket af et Begreb, og en Erkjendelse af en
Tings Væsentlighed og Betydning ifølge en saadan Sammenhæng er baade en
Erkjendelse af et Begreb og en Erkjendelse i Kraft af et Begreb. En Tings Begreb,
som den Sammenhængsmaaden bestemmende Grundtanke, viser sig altsaa som den Grund,
ifølge hvilken de til Tingen henhørende enkelte Led ere, som de ere. Fordi Tingen
er dette eller hiint, det er svarer til dette eller hiint Begreb, har den de
Beskaffenheder og Dele, den har, og staae disse i de Forhold, hvori de staae. Men
Tingen er igjen den Ting, den er, og svarer til sit Begreb, fordi den som Led og
Deel væsentligen hører med til det større reelle Hele, hvori den findes og hvis
Begreb fører den med sig. Dens eget Begreb viser sig her som en af det større
Heles Begreb nødvendigviis følgende enkelt Tanke; hvorved da det Begreb, der før
stod som Grund med Hensyn til sine enkelte Led, nu selv begrundes eller sees i sin
væsentlige Følge af en Forstandsmæssighed.

Erkjendelsen af en, alle givne Gjenstande omfattende, stor Sammenhæng, ifølge
hvilken de alle vilde sees som væsentligen fornødne Led i et altomfattende stort
Hele, saa at det vilde være klart, hvorledes en Forstandsnødvendighed medførte
eller fordrede dem alle med alle deres Beskaffenheder, er Granskningens yderste
Formaal i denne dens Retning.

\begin{center}Anmærkning og Tillæg.\end{center}

1.~Hvor Phænomenerne vise sig som Facta, --- men i Naturen overalt kunne alle
enten betragtes som Facta eller som Resultater af Facta --- gaaer Granskningen ud
paa de \emph{Love}, hvorefter der virkes, og at udgrunde Tingenes Grunde bliver
forsaavidt at udgrunde deres Love, dog naturligviis som saadanne, der indsees som
nødvendige ifølge en intelligent Nødvendighed; men disse Love ere da, som
Virksomhedernes Grundformer, en Art af Begreber.

2.~Vi see her, hvorledes den platoniske Socrates (i Phædon pag.\ 97, c.\ og 99,
e.) kunde føle sig dreven til at gribe efter Anaxagoras's νοῦς, og, for at komme
til en sand Erkjendelse af Tingenes Grunde, tage sin Tilflugt εἰς τοὺς λόγους.%
% FLAG p.91: Greek copied by sense (Phaedo 99e "εἰς τοὺς λόγους"); verify article/accents against scan. Also νοῦς.
 Vistnok er denne Erkjendelse endnu ikke den højeste; den viser selv hen til en
højere, men denne højere Erkjendelse, hvortil den leder hen, til ej fortrænge
den, men bestyrke den, ja den vil først vise den i sin fuldeste Gyldighed.
Imidlertid kunne og ville vi dog strax her, hvor den indtræder, befæste os ganske
i den Erkjendelse, at selve Forstandsmæssigheden er det, hvori den søgte Realitet
hviler, at Tingenes Begreber ere deres egentlige nærmeste Tilværelsesgrunde.

Den almindelige Grund for denne Erkjendelse er den, at overalt det egentlige Fordi
nærmest ej kan være en Gjenstand, men den Nexus, som medfører, at denne Gjenstand
har den første til Følge. Anføres f.~Ex.\ Varmen som Grund til, at et Legeme er
udvidet, Kulden som Grund til Vandets faste Tilstand, naar det er blevet til Iis,
saa erkjendes dog disse Gjenstande endnu ikke som Grunde til disse Phænomener, men
denne Erkjendelse af, at den anførte Gjenstand virkelig er Grund til det omspurgte
Phænomen, opstaaer først, naar der indsees, hvorledes Varmen udvider, Kulden
fastner, og hvorfor den ene ikke kan andet end udvide, den anden at fastne. Men
skal dette indsees, saa maa Varmens og Kuldens hele Natur undersøges med det hele
Indbegreb af Virkninger, den ifølge denne sin Natur indsees ikke at kunne andet
end have. Men Spørsmaalet bliver da, hvad Varme, hvad Kulde egentligen ere. Og
dette Spørsmaal er et Spørsmaal efter Varmens og Kuldens Begreb, hvoraf da Lovene
for deres Virkninger indsees. Dertil vil da nu udfordres at undersøge, hvad
Betydning Varme og Kulde have i Naturens Hele, tilmed og under hvilke Betingelser
de komme frem. Men naar da nu indsees, at Varme f.~Ex.\ maa komme frem under
disse eller hine Betingelser, saa kan dette dog kun blive virkelig indseet, naar
der indsees, hvorledes det af Varmens hele Begreb, af den reale Betydning, den
har i det Hele, følger at den under saadanne Betingelser maa indtræde, hvorved da
Varmens egen indre Natur viser sig som det egentlige aarsagelige (αἴτιον
Phædon.\ 99, b),%
% FLAG p.93: Greek "αἴτιον" (Phaedo 99b, "the cause"); print shows an aspirated form. Verify accents.
 der under de givne Betingelser kommer frem, fordi dets Natur fører med sig, at
det da maa komme frem. Men Varmens Natur er dens ved dens reale Betydning bestemte
ejendommelige Beskaffenhed eller Maade at være paa, tilmed er Erkjendelsen af dets
Natur eet med Erkjendelsen af dets Begreb. Og saaledes forholder det sig ikke
blot, naar Grunden til en vis Art af Phænomener i Almindelighed er Granskningens
Formaal, men ogsaa naar i et bestemt enkelt Tilfælde Grunden til et vist Factum
søges; ogsaa da kan den kun findes, ved at gaae ind i Undersøgelsen af en vis
Sammenhæng, men da vil Indsigten i Grunden beroe paa Indsigten i, at en Gjenstand,
som den omspurgte, ifølge en saadan Gjenstands Natur overhovedet, altsaa ifølge
dens Begreb, maatte træde i Virkeligheden under saadanne Omstændigheder. Man
bemærke vel, at en Grund ikke erkjendes og indsees som Grund til et Factum, blot
fordi Tingen antages og erkjendes som den, paa hvilken altid og efter en bestemt
Regel et saadant Factum følger, men kun, naar det tillige indsees, at dette Factum
maa følge paa den, fordi den nødvendigviis fører et saadant Factum med sig; og
dette kan ikke indsees uden at indsee den Sammenhæng eller Nexus, der fra den ene
fører til den anden; saa at denne Nexus er det, hvorpaa alt beroer. Og dette
gjelder ej blot med Hensyn til vor Erkjenden, men lige saa fuldt med Hensyn til
Tingen selv; thi hvad er det dog ogsaa i denne selv, som gjør, at efter eller med
den ene Ting altid og nødvendigen den anden følger, andet end netop den reelle
Sammenhæng mellem dem, som fører med sig at med den ene den anden maa indtræde?
Indsees det først som aldeles nødvendigt, at af en vis Gjenstand en anden følger,
--- men uden denne Indsigt er Grundethedens Erkjendelse ikke fuldkommen --- saa
maa dette ogsaa erkjendes som noget, der i Gjenstanden selv er saaledes, fordi det
i den ej kunde være anderledes, og Forstandsnødvendigheden viser sig da som den,
der ogsaa i Gjenstanden selv gjør, at Grunden har den Følge, den har, saa at
Forstandsmæssigheden viser sig som det egentligen begrundende. Naar, for at gaae
tilbage til et forhen (§.~7. Anm.~1.) forekommet Exempel Tyngdens Love forklares
og indsees ifølge en mathematisk Speculation, hvad kan saa nævnes andet, som
deres reale Grund i Naturen, end netop det, at en Kraft som Tyngden ej kan tænkes
at virke, og derfor ej heller virkeligen kan virke, paa anden Maade.

Hvad der her er oplyst om den Art af Grund, der kaldes Aarsag, er endnu mere
indlysende, naar flere Ting, der paa eengang findes tilsammen uden Tidsfølge, vise
sig som dem, der føre, den ene den anden, med sig. At Grunden til en vis Dyreforms
Tilværelse maa søges i den hele Dyrerækkes Udviklingsgang, tilmed i den Sammenhæng
i den hele Række, der kræver en saadan Form med, som et integrerende Led i den
hele Række, er let at see; ligesaa, at Hjertet har den Form, det har, fordi den
hele Organisation, altsaa Organismens Grundform eller dens Begreb, fordrer et
saadant Organ med til de øvrige; og saa fremdeles.

3.~Men naar Tingenes Grunde saaledes søges i deres Begreber, vil da ikke
Granskningen efter Grunde tilsidst ende ved noget ganske andet end den søgte
efter, komme an paa et Sted, der syntes allermindst at skulle være det, hvor den
skulde komme an? Spørsmaalet efter Grunden til noget, kræver noget \emph{andet}
end Tingen selv, af hvilket andet denne skal være en Følge, eller med hvilket
andet denne tillige skal være givet, og det er aabenbar intet Svar paa et saadant
Spørsmaal, naar selve det, om hvis Grund der spørges, under en anden Skikkelse
gives tilbage som Grunden. Men skeer dette nu ikke ganske aabenbart, naar man med
Platon, (Phædon anførte Sted), søgende Tingenes Grunde i deres Begreber, siger at
Grunden til, at f.~Ex.\ to Ting ere to, ikke er Tillægningen, hvorved til den ene
den anden er føjet, eller Delingen, hvorved een Ting er skilt ad i to Dele, men er
intet andet, end Toheden; eller naar man f.~Ex.\ vilde sige, at Aarsagen til, at
et Menneske døer, er ikke Giften eller Sygdommen eller deslige, men er Døden,
eller at Grunden til, at det fryser, ej er Nordenvinden, som siges at have bragt
en saadan Kulde med sig, men er intet andet end Frosten? Er deslige ikke den
aabenbare Tautologie?

Og dog er det klart, at, for at blive ved Exemplerne, ikke Giften allene uden
videre gjør, at Mennesket døer, ikke Nordenvinden allene uden videre gjør Vandet
frossent. Ifølge Giftens Indbringelse i Legemet kommer formedelst det Forhold, der
er imellem den og den organiske Leveproces en saadan Tilstand frem, at nu Livet
selv under sin dertil svarende Virken foraarsager Døden, d.\ e.\ at under et
Sammentræf af disse Betingelser, deels de ydre, Giften, deels de indre, den ifølge
sin Natur virkende Leven, den Tilstand indtræder, som nu netop er Døden selv, der
indtræder, fordi under disse Betingelser netop den, og intet andet, maa indtræde,
fordi en saadan Tilstands, et saadant Factums Natur er den at et saadant, og
ligesom føjer denne nye Tanke, og kræver dens Realisation, til de foregaaende.

Og saaledes er, hvad for det første \emph{Facta og Tilstande} angaaer, deres reale
Grunde baade at søge i de Ting selv, hos hvilke de findes, eller disse Tings egen
Natur, d.\ e.\ i deres Begreber, saa og udenfor dem i det, der plejer at nævnes
som deres Grunde og Aarsager. Alt indtræder kun og har kun sin saa og saa beskafne
Bestaaen under visse bestemte Relationer, hvilke altsaa ere Betingelser for samme,
og naar de alle ere tilstæde, maa det selv nødvendigviis ogsaa komme tilstæde og
saaledes antage disse Betingelser ogsaa Charakteren af Grunde. Men naar det nu
indtræder og kommer tilstæde, saa kommer det frem med sin særegne Beskaffenhed,
som noget nyt, der kommer til det øvrige og hvis Plads ej kan indtages af nogen af
disse; det kommer saaledes frem, fordi det selv er, som det er, og med denne sin
særegne Natur og Beskaffenhed hører hen paa denne Plads; det indtræder altsaa
ifølge den særegne, et saadant Noget paa et saadant Sted krævende Tanke, eller
ifølge dets eget Begreb. Det Punct, hvori Tanken om saadanne Betingelser træffer
sammen med Tanken om et saadant Factum, saaledes at det indsees, at dette sidste
under hines Forudsætning nødvendigviis maa tænkes paa denne bestemte Maade, er det
Punct, hvor den egentlige Grund ligger, men dette er netop Tingens eget Begreb. I
det specielle Tilfælde kommer det an paa, at see, hvorledes en vis almindelig
Tanke, et vist Begreb, under disse givne Omstændigheder, nødvendigviis kom til
Realitet; hvorledes f.~Ex.\ under disse Betingelser Frosten eller Døden, ifølge
Begrebet om Frost eller Død, kom til Virkelighed; først da sees Frosten eller
Døden som Aarsag til, at Vandet frøs eller Mennesket døede.

Og spørge vi dernæst for det andet, hvorledes da selve disse Tanker, som fordre
saadanne Factas Indtræden, \emph{selve disse Begreber}, hvori vi søge disse Factas
Grunde, komme frem, saa see vi, at de ifølge en nødvendig Tankegang afledes eller
deduceres af visse foregaaende Begreber, som nødvendigviis ogsaa føre til disse i
et heelt Tankesystem. En vis Hovedtanke fordrer, naar den skal realiseres dette og
hiint bestemte, hvilket ifølge den maa være, men i dette ere nu mange
Enkeltheder, hvilke hver for sig maae være saa og saaledes beskafne, til hvilke
der altsaa svare lige saa mange Begreber, der tilsige, ej blot hvorfor de skulle
nødvendigviis tænkes med, men og hvorledes de desaarsag maae være. Idet da Tingen,
Factumet eller den factiske Tilstand, forklares eller sees i sin Grundethed, fordi
der indsees, at under disse Forhold en nødvendig Tankegang nødvendigviis kræver
netop Tingens Tilvær paa denne Maade, sees nu ogsaa tillige den hertil bestemmende
Tanke, Tingens Begreb, i sin nødvendige Fremgang af en saadan Tankegang.
Imidlertid kommer den dog frem som en særlig Tanke, et særligt Begreb, for sig,
der i Tankegangens Hele vistnok fordres saaledes ved de øvriges paatrængende
Nødvendighed, men dog er en ny Tanke for sig og et lige saa fuldt i sig selv
bestaaende Led, som alle de øvrige, med hvilke den tilsammen danner dette
Tankehele. Saa at Begrebet selv lige meget hviler i sig og i den hele Sammenhæng,
hvori det kommer frem og hvoraf det afledes. Tanken viser sig som det bestemmende
for Tingen og derfor maa dennes Tilværelses og dens Beskaffenheds Grund søges i
dens Begreb; spørges derimod om selve Begrebets Grund, saa er denne at søge i dets
Deduction, i den Forstandsnødvendighed, der i en bestemt nødvendig Tankegang fører
til et saadant Begreb; saa at ej just Grundbegrebet, hvorfra der gaaes ud, men den
Tankesammenhæng og Tankenødvendighed i det Hele, der fra Grundbegrebet fører til
det omspurgte Begreb, er at nævne som dettes Grund.

4.~Spørsmaalet om Grunde træder herunder endnu overalt kun frem som \emph{et
Hvorfor}. Dette fører nærmest fra Phænomen til Phænomen, eller fra en given
Gjenstand til en anden paa samme Maade given. Men fordi ikke denne anden uden
videre kan anføres som Svar, men kun nok; det er ej altid let at bestemme, hvad
der i et Factum er det egentlige Factum. Men er det udfundet, saa bliver det dog
nu først Gjenstand for en Grunden, for nu at udfinde dets ratio essendi. Det
Spørsmaal, hvorfor dog Planeterne gaae om Solen o.\ s.\ v., eller Granskningen
over en Forklaring af det hele Planetsystem, betragtet som et stort forefundet
Phænomen, fører allerede over i en ganske anden Sphære. Den kan nu føre til at see
dette Grundfactum i sin væsentlige Forbindelse med andre, og saaledes Begrebet om
en større Sammenhæng oplade sig, hvorved et \emph{dybere liggende Grundfactum}
vilde komme til Syne, som Opgave for en ny Grunden; men ogsaa kan maaskee
Forklaringen føre tilbage til \emph{intelligente Bestemmelser}, der indlyse ved
deres egen Evidents, og af hvilke hiint Factum sees som nødvendig Følge, hvorved
dog altid en vis Sammenhængsmaade kommer til Syne, som hele Forklaringens
Grundlæg. Overhovedet fører denne Forklaringens Gang alle Spørsmaal om Grunde
tilbage til visse Grundspørsmaal, paa hvis Besvarelse alt tilsidst vil beroe, men
med hvis Besvarelse der ogsaa tilsidst vil blive bragt Lys ind i det Hele. Deels
ville efterhaanden flere sideordnede \emph{isolerede Facta} mere og mere saaledes
sees som \emph{eens eller lige}, at et og samme Grundfactum gjenerkjendes i dem
alle; som naar man i Forbrændingen, i Metallers Forrustning og saa videre,
gjenfinder den samme Iltningsproces (Oxydationsproces). Deels ville flere
sideordnede \emph{i hinanden gribende Facta} eller paa en reel Maade
sammenknyttede Tilstande eller Grundformer saaledes kunne skues som Følger af en
dem alle omfattende \emph{reel Sammenhæng}, at kun denne tilsidst bliver det
tilbagestaaende Problem; som naar alle Sindets Virksomheder og Yttringer vise sig
i en saa væsentlig Sammenhæng med hinanden, at tilsidst kun Tilværelsen af en
saadan sig bevidst og sig selv bestemmende Personlighed midt i en Endelighed
bliver det Problem, der bliver tilbage; eller naar alle enkelte Dyr saaledes
synes at være væsentlige Led i en heel sig udviklende Række, at nu den udgrundende
Stræben vender sig til hele denne i en saadan Række udtrykte Formation overhovedet
med dens Grundbestemmelser, for at spørge, hvorfor den overhovedet existerer. Men
saaledes vil da tilsidst Nødvendigheden af at antage endnu noget \emph{ganske
andet}, end \emph{den hele Kjæde} med alle dens Led, noget, hvori den hele ligesom
hænger, eller hvoraf den overhovedet bæres og holdes, paatrænge sig os mere og
mere; hvilket fører over til den næstfølgende Betragtning.

\begin{center}\large\S.~10.\end{center}

Den Forstandsnødvendighed eller Forstandsmæssighed, der fordrer Tingenes Tilvær
og bestemmer deres Beskaffenheder, de Tingenes Begreber, paa hvis grundige
Erkjendelse Tingenes egen Erkjendelse beroer, og hvilke i Tingene selv vise sig
som det, ifølge hvilket de ere til og ere, som de ere, altsaa som et Grundlæg for
deres Bestaaen og hele Ejendommelighed, kunne vel sættes som de \emph{nærmere}
rationes essendi til Tingene, men i samme i og for sig allene kan dog endnu ikke
deres \emph{første og højeste Grund} være at søge. Om der ogsaa i denne
intelligente Begrundelse vilde blive gaaet ud fra Bestemmelser, --- og fra noget
maa der gaaes ud, --- hvilke (liig Mathematikens første) vare indlysende ved sig
selv, og maatte antages ifølge en umiddelbar intelligent Nødvendighed, saa
begrundedes derved dog ikkun en Antagelse, ikke de tilsvarende Tings Tilvær i en
Virkelighed; og den hele reale Tilvær vilde forsaavidt endnu staae som tilfældig.
Intelligents og Begreb kan vel \emph{medføre}, at saadanne og saadanne Gjenstande
ere der, den kan vel \emph{fordre} dem, men den sætter dem dog ikke ind i
Tilværelsen eller gjør dem virkelige. Den hele Totalsammenhæng, hvoraf Tingenes
Væsentlighed og Nødvendigheden af deres Tilvær følger og indlyser, vil endnu,
saalænge ikkun den allene viser sig som Erkjendelsens sidste Grundlæg eller første
Udspring, staae som et stort Problem, hvorom der ligesaavel maa spørges, hvorfor
den overhovedet og i sin Heelhed er saaledes, som den er, som der spurgtes,
hvorledes dog det virkelig Hele, hvori den findes virkeliggjort, kommer til at
være virkelig til. Vi maae altsaa endnu søge Tilværelsens sidste Grunde i noget
andet end den hele Forstandighed, i et dem begge, den givne Virkelighed paa den
ene Side, Forstandigheden paa den anden Side, forbindende Tredie, der overalt
bringer Begrebet ind i det umiddelbart virkelige (ind i Objectiviteten), og gjør,
at hiint er udtrykt i dette. Hvilket Tredie da maa være lige saa omfattende, som
den hele Sammenhæng, ligesaa fuldt, som denne, vise sig som Eet og alt.

Dette kan da ej blot være en Forstandighed, men maa tillige være en Kraftighed:
et \emph{virkende Første}, der viser sig som et Princip, \emph{hvorved} Tingene
ere, som de ere, saa at det overalt i sin Virken er tilstæde i dem og de kun ere
dets ydre Udtryk og Virkninger. Et Princip, der ikke blot \emph{bestemmer} (paa
en ideel Maade), men realiserer eller \emph{virkeliggjør} (paa en reel Maade). Et
saadant begrunder og giver Væsenhed ved selv virkende at være tilstæde i Tingene
og ligesom at udfylde dem med Væsenhed, ret ligesom Sjælen overalt er tilstæde i
enhver Deel af Organismen. Det er det sande \emph{Værende} og Væsen, det sande
Reale eller Substantielle, og kun, naar det skues i dem, eller de skues som
medhenhørende til dets fuldelige Fremtraden, skues de selv saaledes til Grunden,
at de derved, som ved et indre Lys, oplyses. Det skues kun, som det skal skues,
naar det sees som \emph{en Leven}, der levendegjør eller dog virkeliggjør alle de
Ting, der ved samme begrundes, og alle disse Ting, der nu vise sig som et blot
Substrat eller Materiale for sammes Leven og Virken, sees kun, som de skulle sees,
naar de sees, som levendegjorte eller virkeliggjorte derved, saa at disse i deres
hele Indbegreb kun sees som dets Udtryk, og det selv sees som det, der boer i dem
og straaler os imøde eller taler til os ud af dem.

Istedet for hiint Hvorfor viser sig her et \emph{Hvorved}, og det saaledes, at
det, hvorved Tingene ere, viser sig som det, der som det indre Væsen er
\emph{tilstæde} i disse Ting. Men hiint \emph{Hvorfor} falder dermed ikke bort.
Der kan og maa endnu bestandigen spørges: hvorfor virker dog hiint første Værende
og Virkende netop paa \emph{denne} bestemte Maade, hvorfor træder det just ud i
\emph{disse} Virkninger? Et Spørsmaal, der, naar man gaaer ud fra dets
Betragtning, som Grundbetragtning allene og ikke seer hen til en allerede given
Virkelighed, vil lyde: i \emph{hvilke} Virkninger og Udtryksformer vil hiint sande
Værende træde frem? Overhovedet kan det ikke tænkes som Tilværelsesgrund til nogen
enkelt Gjenstand uden ifølge en vis altomfattende Sammenhæng, en vis
Planmæssighed, en bestemt \emph{Tingenes Orden}, hvorefter just disse eller hine
bestemte Ting fremvirkes, idet den fører dem med sig. Men denne hele
Sammenhængsmaade eller Orden i Tingene maa dog nu selv være grundet i hiint
virkende Første, hiint sande Reale. Men spørges: hvorfor just denne og ingen
anden? eller spørges bedre: hvilken Tingenes Orden maa da hiint første sætte ind i
Tilværelsen? saa kan intet andet Svar gives, end det samme, der maa gives, naar
Spørsmaalet er: hvorfor er overhovedet endnu noget andet til, end hiint sande
væsentlige allene? hvorfor sætter det noget andet ind i Tilværelsen? Og herpaa
kan Svaret kun være: for deri fuldkomment \emph{at aabenbare sig selv}; altsaa for
at dette andet kan være (eller vorde) dets fuldkomne Udtryk, saa at det selv
fuldkomment kan opfylde det og være deri. Det, hvortil vi maae see hen for at
finde alle Tings \emph{første Udspring} og bestemmende Grund, er ogsaa det, hvorpaa
vi maae vise hen, naar der spørges: til \emph{hvad Ende} alt er til. Det,
\emph{hvorved} det er, ogsaa det, \emph{hvorfor} eller for hvis Skyld det er. Men
saa maa da ogsaa hiin Sammenhængsmaade eller Tingenes Orden være saaledes, som
dette fordrer den, for fuldkommen at kunne aabenbare sin Væsenhed. Vi see altsaa,
hvorledes det indre alt begrundende Værende medfører og fordrer, idet det tillige
virker og virkeliggjør. Medens det selv er eet og alt, som det, der overalt i alt
er tilstæde som det allene i Sandhed værende, fordrer det dog tillige dette alt,
som noget \emph{andet} end sig selv, hvori det nu skal være som det altopfyldende,
og staaer saaledes tillige selv i Forhold til dette andet, som om det blot var Led
i Kjæden, om end første bevægende Led. Hvilket især indlyser, naar nu den hele
Sammenhæng skal indsees og erkjendes som ved samme er udtrykt i Realiteten, og
ifølge hvilken hver enkelt Ting bestemmes og virkes frem; thi da maa der gaaes ud
fra visse første Bestemmelser, som dem, hvilke hiint altreale værende først
fordrer, til hvilke altsaa Erkjendelsen først føres, og da deri strax en
Sammenhæng maa vise sig, staaer det altreale værende selv som det allerførste i
denne Sammenhængs Række eller Kjæde, altsaa som begyndende første, som Led. Dets
Betydning med Hensyn til det andet, dets Forhold dertil, maa nu bestemmes og
tilmed et \emph{Begreb} dannes derom. Men uagtet det saaledes staaer som første
bestemmende Led, saa er det dog ej blot det første, men tillige det altbevægende,
der besjælende og bestemmende gaaer igjennem det Hele, og overalt viser sig deri
paa ethvert Punct, saa at det Hele, som nu fremkommer, i sin Heelhed skal blive
dets totale Udtryk. --- Dette sande Værende, Absolute og Evige, forsaavidt det
skues i eet med den Form, under hvilken det træder frem, eller det Begreb,
igjennem hvilket det kommer til sin reelle Fremtræden, er det, vi kalde
\emph{Idee}. I Ideen finder altsaa al Erkjenden sin sidste Støtte og Hvile, idet
den deri finder sit absolut første Udgangspunct.

\begin{center}Anmærkning og Tillæg.\end{center}

1.~Vi have kun faa Træk at tilføje for at føre Tankegangen op til det allerhøjeste
Punct, eller til det Punct, hvor den højeste Idee kommer. Er Tingenes første og
eneste sande ratio essendi at søge i et altreelt ene Værende og Princip, men er
dette et saadant, der tillige medfører og sætter et uendeligt sammenhængende Hele
og en uendelig Forstandighed eller Planmæssighed i dette Hele, saa viser det sig
ej blot som en Væren og Leven, men som et altreelt værende og levende Væsen, som
sætter en Uendelighed, for at aabenbare sin Uendelighed deri, men som dog i sig
selv er uafhængigt af det, hvad det saaledes sætter, bestemmende det, men ikke
bestemmet derved.

Men hvorfor sætter eller fremvirker det da dette Hele, naar det i sig bestaaer
uden samme? \emph{Hvorfor er det ej blot, men aabenbarer sig tillige?} Vi svare:
for at det kan træde frem i uendelig \emph{Fylde}. Altsaa dog for at det kan end
mere fuldeligen \emph{være}. Der maa være noget hos samme selv, som kræver dets
Aabenbarelse, saa at det ej kan og vil beroe i sig selv allene. Vi kunne ej tænke
os det at træde ud og aabenbare sig, uden at tænke os det, som om det \emph{maatte}
aabenbare sig og ej kunde andet, men dog kan der ikke existere noget nødende eller
nogen bydende Nødvendighed for samme; maa det aabenbare sig, saa er det, fordi det
\emph{maa ville} aabenbare sig. Men hvorledes skal et Maae tænkes med fuldkommen
Frihed, hvorledes skal det, der aldeles ikke \emph{behøver} noget andet udenfor
sig til sin Bestaaen, dog tænkes, som det, der \emph{kræver} noget andet udenfor
sig, som om det dog alligevel ej var sig selv nok? Det kan kun tænkes som en
Nødvendighed og et Krav af \emph{Kjerlighed}. Thi Kjerligheden har det i sig, at
den ej kan være uden det andet, uden sin Kjerligheds Gjenstand, og at den dog ej
kræver denne af nogen anden Grund, end allene for at elske, hvilket her ikke er
for at modtage, men blot for at give. Det, der foruden det højeste Væsen, hvilket
nu, som elskende Væsen, nødvendigviis maa tænkes som højeste \emph{Personlighed},
endnu maa være til, er \emph{dets Rige}, Guds Rige, som et Hellighedens og
Kjerlighedens, tilmed og Skjønhedens og Fyldens Rige. Men derfor maa og en
\emph{Frihed} være og et \emph{Frihedens Rige}, og for at dette kan være, maa
Friheden selv gjøre sig dertil, hvorved vi da komme til den forhen (§.~4, Anm.~3)
allerede fremkaldte Betragtning. (Endnu bemærkes kun, hvorledes Ideen om Guds Søn,
som den evige Selvaabenbarelses og Kjerligheds første og tillige altomfattende
Gjenstand, her frembyder sig ganske ligefrem; men denne Overvejelse ligger for
fjærnt fra de nærværende Undersøgelser, til at vi kunne gaae ind i den her.)

Anm.~2. Det højeste Formaal i denne Granskningens Retning, og for al Stræben efter
Erkjendelse overhovedet, bliver nu at skue alt saaledes i Forhold til Gud, at
denne højeste Idee, ligesom den \emph{practiskt} skal være den \emph{altbesjælende},
nu ogsaa \emph{videnskabeligen} bliver den \emph{altoplysende}, som vi allerede
forhen i Anm.~3 til §.~1 i Forbigaaende have bemærket. Men en philosophisk
Erkjenden (en Erkjendelse i Kraft af Idee) kan paa mange Puncter allerede
fuldkomment indtræde, som saadan, uden at endnu Gjenstanden skues i det højeste
Lys. Det, der gjør en Erkjenden til en philosophisk, er overhovedet ikke det, hvad
der erkjendes, allene, men er Maaden, hvorpaa det erkjendes; at det nemlig
erkjendes i Kraft af en Grunderkjenden af noget, der skues som det fremvirkende
Princip, men derved og som det al Realitet givende ene Reale eller værende og
væsentlige, og med Hensyn til hiint,%
% FLAG p.111: print reads "iog med Hensyn"; rendered "og med Hensyn" (broken ligature?). Verify.
 hvad det nu saa er, der skal erkjendes i Kraft deraf.

\begin{center}\large\S.~11.\end{center}

Hvad der bestemmer alle Tingene deres, mere eller mindre relative, Væsenhed, og
med Hensyn til hvilket de allene ere og betyde noget, maa naturligviis ogsaa være
det, der giver selve Erkjendelsen sin Væsenhed og gjør den til den, den er; saa at
selve \emph{Erkjendelsens sidste Tilværelsesgrund} ogsaa er at søge deri. Ogsaa
Erkjendelsen er da en Form, hvorunder det Evige aabenbarer sig eller realgjør sig
selv, og den fuldkomne Sandhedens Erkjendelse er i sit Væsen selve \emph{Ideens
Leven og Sigfremstillen i Bevidstheden}. Men skal den være en fuldkommen
Erkjenden, saa maa den og vise sig for Bevidstheden og haves og besiddes af den
erkjendende Personlighed som det, den saaledes er, saa at den erkjendes for at
være Ideens egen Leven i os, til hvilken vi derfor med villig Hu skulle hengive
os, at være dens rene Kar eller Organer.

\begin{center}Anmærkning og Tillæg.\end{center}

1.~Det Punct, paa hvilket denne, til de højeste Former henhørende, Form for det
Eviges Aabenbarelse indtræder og skues i sin væsentlige Sammenhæng med, og sin
Betydenhed i, det Hele, (hvilket Hele overhovedet kun er til for at tjene det
Evige, at dette kan komme til et Rige og Herredømme deri), er indlysende af det
foregaaende, hvor vi have seet, at et Frihedens Rige maa komme frem, for at et
Kjerlighedens kan komme, men at hiint igjen nødvendigviis fører til en sig som
saadan fremstillende, det er til Bevidsthed kommende, Intelligents.

2.~Hvad den \emph{Sandhed} er, hvis Dyrker Granskeren skal være, bliver nu
aldeles klart: at her nemlig endnu er noget ganske andet end en blot
Overeensstemmelse imellem vor Forestillen og dens, højere eller lavere,
Gjenstande. Sandheden eller det Sande er det, der \emph{levende virker denne
Overeensstemmelse}, idet det selv virkende bestemmer Forestillingens hele Maade
og Dannelse. Netop hvad der giver Tingene, Formerne og Forholdene selv, Realitet
eller Væsenhed og Sandhed, viser sig nu for os (og har allerede forhen forkyndt
sig for os, Anmærkning til §.~7,) som det selvsamme, der ogsaa, gjenlevende i os
og vor bevidste Erkjenden giver denne sin Realitet. Sandheden selv, der nu viser
sig ogsaa at være, ligesom f.~Ex.\ Dyden, Udtrykket for en Idee, skal være det
levendegjørende Princip i vor Erkjenden, ligesom Dyden skal være det i vor
Handlen, og ligesom denne sidste er det under een Form, saaledes er Sandheden det
under en anden Form i os levende, tilmed tilstædeværende, Uendelige og Evige,
hvis Organer vi skulle vorde ved fuldkomment at bringe dets indre Tilsagn til
Bevidsthed, Klarhed og et i indre Udtalelser og Productioner kraftigt Herredømme
i os. Saa at vi ogsaa i Relation til Naturen og Tilværelsen selv skulle tage, og
det med klar og velbevidst Erkjendelse tage, alt, hvad vi i vor Erkjenden have og
besidde fra første Haand. Skal da nu et reproduceret Billede af det Hele gjenleve
i os og vi netop derved voxe fuldkomment ind i, og sammen med, dette Hele, saa
skulle vi dog heri have mere end et blot Afbillede; vi skulle deri tillige have
den mægtige Indflydelse og Tilstædeværelse af det Evige selv, som ogsaa i selve
hiint virkelige Hele er det sande reelle og værende.

Det forstaaer sig iøvrigt af sig selv, hvad vi dog ej ville undlade udtrykkeligen
at bemærke, at det her udtalte ej maa udtydes saaledes, som om Meningen kunde
være, at f.~Ex.\ Tyngdeprincipet, eller den under Tyngdens Form sig yttrende
Virkning i Naturen, nu som saadan skulde antages at være det, der fremkaldte vort
Begreb derom eller vor Erkjendelse deraf. Hvorimod vel det samme Ene, der
medfører eller fordrende og virkende sætter en Tyngde i Naturens System, nu ogsaa
sætter Begrebet derom i vort Forstandssystem.

3.~Men hvad saaledes det (væsentlige og ægte) Sande, hvad selve Ideens Magt
virker i os, er endnu ikke derfor strax erkjendt i Kraft af det Udspring, hvorfra
det i sig selv dog virkelig kommer. Fordi en Erkjendelse i sig selv gaaer frem af
en Ideens Leven i os, derfor erkjendes der endnu ikke strax i Kraft af Ideens
Erkjendelse, derfor sees Gjenstandene endnu ej i Ideens, til Bevidsthedens
Klarhed komne, Lys. Saaledes er f.~Ex.\ Mathematiken vistnok, som Videnskab, og
som den denne Videnskab til Bevidsthed kaldende og bringende Tanken, Udtrykket af
en Idee og en Ideens Leven i os. Men mathematisk Erkjenden er dog ej Erkjenden i
Kraft af Idee.

Endog den blotte lavere Sandsning og aldeles empiriske Opfatten og Erkjenden er
at betragte paa samme Maade, som den, der i sig selv fremkommer ved det
Uendeliges egen Leven og Virken i os, skjønt den er saa langt fra at komme til
Bevidsthed som et Udtryk af en højere Leven i os, at den ikke engang synes at
være virket indenfra, hvilket den dog, som et Udtryk af vor Leven, som en
Livsyttring og et Resultat af en Livsvirksomhed, ikke kan andet end være, som den
da og allerede viser sig saaledes, naar den Productivitet bemærkes, hvorved den
kommer i Stand. Men denne Productivitets egentlige Bevæger, dens bestemmende
Første, kan dog nu ligeledes kun være en Leven, eller et egentligt, tilmed indre,
Princip. Og nu føre ogsaa alle Overvejelser til at see en almeen, i det Ydre som
i det Indre ligefuldt tilstædeværende, og baade i hiint og i dette alt
bestemmende, Virken som sidste og sande Grund til al lavere Sandsning. See
Forfatterens Psychologie §.~58 pag.\ 212.

\begin{center}\large\S.~12.\end{center}

Skal Ideen komme til at leve som saadan i Individet ved at fremstille sig som Idee
i dets Erkjenden, saa maa den træde frem som det styrende og bestemmende Første,
som det dannende og ordnende Princip, i et Indbegreb af Erkjendelser, der komme
til at udgjøre et sammenhængende Erkjendelsens Hele, i hvilket Ideen, idet den
fremkalder og bestemmende uddanner det, fremstiller sig selv. Denne Ideens
Fremstilling ved og ifølge dens egen Sigfremstilling, hvori den selv (som det
mere indre) haves og skues i eet med dens (mere ydre) Fremstillelse, er
Videnskaben, den ægte, fuldkomne Videnskab. De enkelte Erkjendelser ere
forsaavidt de haves og skues som væsentligen medhenhørende Led i et saadant Hele
og nu derfor gaae frem af den Gransken og Indseen, hvorved et saadant Hele
efterhaanden søges bragt i Stand, \emph{videnskabelige}.

I Videnskaben erkjendes altsaa ej blot \emph{det Ene}, hvilket overalt er at
betragte som Væsenet, hvortil alt er at føre tilbage og hvorpaa alt kommer an,
men ogsaa den derved bestemte \emph{Mangfoldighed} i sin uendelige Forskjellighed,
dog saaledes, at alt i den skues baade i sit nærmere eller fjærnere Forhold til
hiint Ene, tilmed i sin større eller ringere Realitet, saa og i sit bestemte
Forhold til alt det øvrige, med hvilket det tilsammen udgjør Ideens fuldelige
Fremstilling, tilmed paa sin bestemte Plads i det sammenhængende Hele. Saa at i
den fuldkomne Videnskab disse trende: \emph{Idee}, \emph{System} og
\emph{Detail} ere uadskillelige.

Til saadan Videnskab skal \emph{al} Erkjendelse uddannes, en saaledes
videnskabelig Erkjendelse skal \emph{enhver} Erkjendelse være.

Thi skal --- dette er det første --- Ideen skues som Idee i sin Sandhed, maa den
nødvendigviis skues som den, der giver Væsenhed, tilmed maa den skues i eet med
det, hvilket den giver denne Væsenhed, eller skues som det bestemmende og
altoplysende i hiint andet, hvilket ved samme har sin Realitet; Ideen vilde
ellers ikke skues i sin fulde Sandhed; thi det hører dens Væsen til at give og
bestemme Væsenhed. Dernæst: skal det hele Indbegreb af enkelte Gjenstande, som
fremstiller sig for os og paa hvis Udgranskelse vor Stræben gaaer ud, erkjendes i
sin Sandhed, eller erkjendes tilfulde, som det er, og er det nu kun noget med
Hensyn til det Evige, har det kun Betydning formedelst sin Relation til dette, saa
maa det ogsaa skues i sine Forhold til dette og i denne sin Relativitet; men skal
saaledes alt andet bestemmes og oplyses med Hensyn til hiint Evige og ved dette
Eviges Idee, saa maa nu ogsaa denne Idee vise sig som det altoplysende og
altbestemmende i denne hele mangfoldige Erkjendelses Kreds, men da maa denne ogsaa
ved hiin Idees Kraft dannes til et Erkjendelseshele, hvori Ideen da vil vise sig
som det altreelle og altopfyldende. Til at skue de enkelte Led eller Dele i deres
Sandhed hører ej blot, at det reale Væsen, der er deres sande Grund, skues i dem,
eller at de skues i Relation til dette, men ogsaa, at den Forstandsnødvendighed,
tilmed den Tankesammenhæng, sees, ifølge hvilken hiint højeste værende just
kommer til foruden alt det andet, ogsaa at føre til, at netop et saadant Enkelt
faaer Realitet, thi kun derved kommer dette Enkelte jo overhovedet til sin
Realitet. Men det enkelte Led kan ikke sees paa sin Plads og i sin Relativitet,
uden idet en heel systematisk, alle dermed sammenhængende Led tillige omfattende,
Erkjenden kommer i Stand.

\begin{center}Anmærkning og Tillæg.\end{center}

Ideen træder i Videnskaben frem i et \emph{Andet}, i hvilket Andet den nu viser
sig at være Sandheden eller at være det Sande og Væsenet. Ja den viser sig først
ret egentligen som Idee, naar den lyser frem som det Sande i det Andet, hvis
Sandhed den er.

Den bliver saaledes vistnok underlagt en \emph{Modsætning} og betragtet som den,
der kun i en saadan Modsætning sees i sin Sandhed og Betydning. Selv naar vi tale
om Ideen i og for sig, i dens Reenhed, tænke vi os den lige overfor det Andet,
over hvilket den skal hæve sig for at sees i sin Reenhed. Men anderledes kan det
aldeles ikke være; thi allerede, idet Ideen fremstiller sig for Bevidstheden, er
her noget foruden den selv, nemlig den Bevidsthed, til hvilken den kommer; saa at
den allerede forsaavidt strax træder frem under Relation, og det endnu mere, fordi
vor Skuen nødvendigviis gaaer i eet med en Productivitet, og denne, hvori vor egen
Leven udtrykker sig, strax træder ind i alle de Relationer, hvilke vor Leven
overhovedet er underlagt. Fordi vi selv høre med til hiint Andet, i hvilket Ideen
træder frem, kunne vi ej have den aldeles uden al Tiltræden af dette Andet. Saa
at det Udtryk: Ideen i dens Reenhed, kun kan have relativ Gyldighed. Men denne
relative Gyldighed har det ogsaa. Idet Ideen lyser frem igjennem det Andet, maae
vi dog skue den selv som den, der er noget andet, end dette Andet, hvorigjennem vi
skue den, og idet vi, saa meget mueligt, see begge i deres Modsætning, see vi dem
hver for sig, tilmed see vi og Ideen i dens Reenhed, ligesom udskilt fra det
Andet, hvoraf den lyser frem.

\begin{center}\large\S.~13.\end{center}

Det første i den ægte og fuldkomne Videnskab, det sande reale og substantielle i
den, ved hvilket først alt dens Mangfoldige har nogen Betydning, er \emph{Ideen}.
Den skal, som det sig selv constituerende og realgjørende levende Princip, gaae
igjennem den og leve i den og derved bære den Hele; derfor som et sandt
Middelpunct møde os overalt i Videnskaben, eller være det, hvorpaa ethvert Punct i
samme viser hen. Som Aanden i den hele Videnskab, som det, hvorved alt deri først
finder sin Sandhed og Realitet, er den ikke i disse eller hine Hovedbegreber eller
Hovedsætninger allene at søge; kun den hele Videnskab i sin uendelige Fuldendthed
vilde være dens adæqvate Udtryk, saavidt overhovedet det, der har bestemt Form,
kan udtrykke eller udtale det sande uendelige og tilmed over al Form i reen
Aandelighed og Frihed levende. Ikke heller er i nogetsomhelst udtalt Hele, som i
et fuldført Værk, Ideen anderledes tilstæde end i den virksomme Leven, hvorved
dette Værk, dette dens Udtryk, bestandigen paa ny kommer i Stand og bringes frem
eller skabes, eller bestandigen reproducerer sig selv. \emph{Ikke i staaende
Resultater}, men i det som gjør, at saadanne Resultater \emph{fremkomme}, have vi
Ideen og det sande rationale i Erkjendelsen.

I at henføre alt bestemt, alt, hvad der viser sig med en bestemt Form og som et
bestemt Noget, til det Uendelige og Evige, som til det Absolute, hvorfra al
Betragtning maa gaae ud og hvortil den maa gaae tilbage, viser sig det
\emph{philosophiske} i Erkjendelsen; og da alt skal saaledes henføres til det
Absolute, saa maa al Erkjenden stræbe efter at vorde en philosophisk. Saavidt nu
herunder overalt en umiddelbar Øinen og Skuen af dette Absolute ligger til Grund,
indtræder heri overalt en \emph{Speculation}. Det speculative viser sig deri, at
under denne Henføren og Bestemmen det samme ene Absolute, som det altomfattende og
altværende, overalt i sin Heelhed og uendelige Betydenhed staaer for Aandens Øje,
som det, der idet det bestemmer og ligesom tilkjender alle Ting deres relative
Væsenhed, selv er over al Form og udgjør det Skikkelsesløse og usynlige
Middelpunct, til hvilke alle Retninger vise hen, men som kun kommer til Synlighed
ved at være det, hvortil alt andet henføres. Speculationen er her en reen Skuen;
thi, hvad der som det aldeles første og umiddelbare kun haves, idet det skues som
det levendegjørende Princip og Væsen i alt andet, kan kun erkjendes, idet det nu
ligefrem sættes og sees som det umiddelbare Værende, som Væsenet i de Former, som
ved samme blive dets Udtryk, ligesom Livet i Organismen maa ligefrem skues som det
første reale værende i den.

Af Ideen gaaer vistnok nødvendigviis en Indseen og Begriben frem. Men Ideens
Erkjendelse er dog ingen Indseen og Begriben. Thi denne sidste forudsætter
allerede, at det betragtede skues i en vis Sammenhæng og bestemmes paa en vis
Maade ifølge en Sammenhæng, saa at det altsaa allerede sees i Relationer og som
Led i en Kjæde. Imidlertid er Begreb væsentligen forbunden med al Idee, og idet
Ideen træder frem i Videnskaben, træder den frem i et System, hvori alt er Begreb
paa ethvert Punct, saa at nu det Absolute selv (som forhen er bemærket) kommer til
at staae som Led, vel som første og ypperste, men dog altid som Led. Men i og
under den saaledes fortløbende Indseen og begrebsrette Afleden og Bestemmen haves
da Ideen ej i dette ypperste Led, men haves som det, der, ligesom det i den
betragtede Mangfoldighed er det ene reale Væsen, der i sin Eenhed træder frem af
den, saaledes for os bliver Sjælen i vor hele Afleden og Bestemmen. Erkjendelsen
er en ideel Skuen, der haves som Grundlægget i en levende uendelig Tænkning,
hvori det Absolute bestandigen, idet det middelbargjør sig i en Mangfoldighed,
viser sig selv i sin rene umiddelbare Absoluthed.

\begin{center}Anmærkning og Tillæg.\end{center}

1.~Ved Ideens Erkjendelse og den Oplysning, denne giver, kommer os det
mangfoldige empiriskt opfattede eller gjennem Reflection udfundne \emph{Detail}
af enkelte, mere eller mindre dybt gaaende, Erkjendelser først \emph{ret tilgode},
fordi vi derved først see, hvad vi egentligen have i disse. Forsaavidt er Ideen
det, hvorpaa alt først kommer an, det, hvorefter der først skal tragtes. Men paa
den anden Side er igjen det mangfoldige Details nøjagtige Studium det, hvorved
først Ideen ret kommer os tilgode. En vis Fordybelse i det philosophiske,
hvorunder da det enkelte og specielle kun betragtes som Middel og Vehikel for
Ideens Betragtning, og visse Hovedsandheder eller Hovedsætninger med deres
alsidige Overvejelse først haves frem, er paa den ene Side lige saa vigtig i det
akademiske Studium, som Materialets og Detaillets nøjagtige Eftergranskning i
dets hele Omfang paa den anden Side.

2.~Ligesom Viisdom bestaaer i at vide, hvad der har ubetinget Værd, hvad der kun
har betinget Værd, men dog har Værd, og hvad der aldeles ikke har nogen eller kun
har en ringe Værd, og nu at have dette for Øje i alle Domme og Handlinger og
bestemme disse i Overeensstemmelse dermed, saaledes, viser Philosophie sig i at
erkjende og vide, hvad der har absolut Væsentlighed, hvad der kun har en relativ
Væren og Væsenhed, og hvad der endelig ingen har eller kun har den som Skin eller
dog i den laveste Forstand, og nu derefter at opløse alle videnskabelige
Spørsmaal, saa at det indsees, \emph{hvad} enhver Ting med Hensyn til det højeste
værende kan siges egentlig \emph{at være}. Skjønt nu dette vistnok ej kan skee
uden at see enhver Ting paa sit Sted i et videnskabeligt Hele, og al philosophisk
Grunden og Bestemmen derfor nødvendigviis har Hensyn til et saadant, saa viser
dog ægte og dyb Philosophie sig ej blot i et videnskabeligt Systems Udførelse
eller Besiddelse, men den kan og lægge sig for Dagen og træde ud med levende Kraft
og Klarhed i den Maade, hvorpaa enkeltviis givne Spørsmaal behandles og enkelte
Tilfælde betragtes og bedømmes. Ja dette umiddelbart \emph{judiciøse}, hvori
Philosophien kommer til Syne paa en mere practisk Maade end i Systemet, om end
endnu blot i Erkjendelsens Gebeet, hører væsentligen med til en ægte Philosophies
Fremtræden. Philosophie viser sig forsaavidt i en vis \emph{videnskabelig Maade at
tænke og dømme paa}; ja selv, naar den træder frem i et vist systematiskt
Indbegreb af bestemte Erkjendelser, haves den ej i disse, men i den levende Tænken
og Skuen, som har fremkaldet dem og bestandigen kan fremkalde dem paa ny i andre
Udtryk og under nye Vendinger, saa at ogsaa heri det samme maa vise sig, som i
hiint umiddelbare frie judiciøse.

3.~Hvad Speculationen og den Skuen, for hvilken Ideen bestandigen staaer i sin
Totalitet, som det eet og alt, hvorfra alt gaaer ud og hvortil alt gaaer tilbage,
er i videnskabelig Henseende, det er i blot personlig og individuel Henseende
det: at have et \emph{Ideal} for Øje, bære et Ideal i sig og være begejstret af og
for samme. Idealet er her den aldeles fuldendte Videnskab, som Ideens fuldelige,
intet mere manglende Fremstilling, hvori alt vilde skues i sine uendelige
Forbindelser og Betydninger. Til et saadant skue vi hen midt under det Opnaaedes
Mangelfuldhed, for at styrke os under vor stykkevise Arbejden og Bestemmen ved
Henblikket til det Fuldkomne, hvorpaa dog det ufuldkomne, vi allerede have, viser
hen, og hvis Realitet vi allerede skue og ligesom forud besidde, idet vi finde os
paa Vejen dertil.

4.~Men hvorledes er dog Begreb forskjelligt \emph{fra Idee}, hvad Maaden angaaer,
paa hvilken begge træde frem og haves? Thi ogsaa et Begreb haves kun paa en ægte
Maade som det bevægende og bestemmende Første og Ideelle i de Omdømmer, i hvilke
det udtales. Det er i disse Omdømmer ej at søge i Subjectet, ej i Prædicatet, men
i den Grundtanke, ifølge hvilken et saadant Subject forbindes med et saadant
Prædicat, ja ej i eet, men i mange Omdømmer udtales et Begreb, ej af eet allene,
men først af mange lyser det frem, som det, hvorfra disse Omdømmer gaae ud, men
hvilket ligesom i Frihed svæver over dem. Hvorledes er da Idee endnu noget andet
end Begreb?

Begrebet gaaer ud paa det særlige Led, hvilket tænkes saaledes, som det tænkes,
fordi Begrebet selv, som den Grundtanke, der bestemmer dets reale Betydenhed og
dets Tilværelsesmaade, bestemmer, hvorledes det skal tænkes; og da dets reale
Betydenhed beroer paa Relationer og bestemmes med Hensyn til Relationer, hvilke
kunne være utallige og variere paa mange Maader, saa bliver derfor nu Begrebet
herskende og bestemmende Grundtanke for en Mængde, ja en Utallighed af
Determinationer, og vi see tillige, hvorfor det saa ofte viser sig, at det
selvsamme kan henføres og fra en anden Side igjen ikke kan henføres under eet og
samme Begreb, hvorfor stringente Definitioner ofte ere saa vanskelige, om ej
uopnaaelige. Men forsaavidt heri dog kun \emph{et Led}, et vist Enkelt, betragtes
i sin reale Betydning og nu bestemmes efter sine Relationer, (dog at just denne
dets reale Betydnings Relativitet medfører, at Bestemmelsen gaaer ud i en
Utallighed af Determinationer, der kunne variere ideligen, og hvori det constante
allene er det, at Dommen i ethvert enkelt Tilfælde falder saaledes ud, som den
falder ud), forsaavidt er her, uagtet her er noget ideelt som bestemmende Første,
dog kun et Begreb. Skues derimod i det givne Enkelte tillige det ej blot
bestemmende men og reelt \emph{virkende Første}, som gjør, at en saadan Tanke
overalt i sine Relationer kommer til reel Virkelighed, eller som sætter Begrebet
ind i Virkeligheden, saa skue vi deri Ideen. Men nu realgjøres Begrebet ej blot i
Tingene, vi betragte, men og i os selv fremstaae disse bestemmende Grundtanker.
Heri er da og Ideen tilstæde. Det, der gjør, at et saadant Begreb kommer til
Leven i os, det der levendegjør det nu i os i sin særlige Sphære sig bevægende
Begreb, dette er Ideen. Derfor er det ogsaa, at kun de Begreber, som gaae ud paa
Gjenstande, der skues som Led i det Hele og betragtes i denne deres Samknytning
med alt andet, vise sig som organiske Middelpuncter for en Utallighed af
Bestemmelser og Omdømmer; de mathematiske Begreber have ej denne Charakteer; i
Mathematiken sees hver Form aldeles isoleret for sig, deri er ej heller Tale om
reale Betydninger. Hvor disse haves for Øje, skues strax hen paa et Hele, men i
Erkjendelsen af den Relativitet, der derved overalt viser sig, rører sig i sig
selv en Fornemmelse af det Absolute. Kommer det til Erkjendelse, hvad det er, der
gjør, at vi bestemme og dømme og indsee, som vi i Kraft af Begrebet gjøre det, saa
føres vi allerede ind i et System og bringes nu og til at øine det, hvorfra alt
System gaaer ud, hvori alt System hænger, og saa see vi allerede Ideen, som det,
der aabenbarer sig paa det enkelte Punct igjennem Begrebet, og den Erkjenden, der
i sig selv allerede var et Udtryk af Ideens Leven i os, bliver nu en Erkjenden i
Kraft af Ideens Erkjendelse.

\begin{center}\large\S.~14.\end{center}

Idet det rene speculative vil udtale sig i et videnskabeligt Hele, indtræder det
\emph{dialectiske}, der er Udtrykket af den sig i Forhold til det Mangfoldige
sættende, og dette bestemmende og efter sig selv formende, Speculation selv.
Dialectiken indtræder i Ideens og Speculationens Overgang til et systematisk
Heles Uddannelse. Den fremtræder derved paa eengang og i eet, for det første paa
en \emph{negativ} Maade som en Opløsning og Tilintetgjørelse af alle de paa lavere
Standpuncter sig fixerende Forestillingsmaader, forsaavidt disse, maaskee endog
naturligen, ja nødvendigviis, have antaget Præget af en Betydning, ja en
Absoluthed, der ej tilkom dem; men dernæst for det andet tillige paa en
\emph{positiv} Maade som den, der just, idet den viser Tingene tilligemed deres
Begreber i deres relative Ugyldighed eller Intethed, \emph{eo ipso} viser dem i
deres relative Betydenhed eller Nogethed, og saaledes fører til at bestemme dem
deres Plads og at udføre eller \emph{construere det videnskabelige Hele}, hvori
alt sees i sine Relationer, og tilmed sees, som det skal sees. Saa at det
dialectiske, som det, der gaaer frem af Ideen, der skues og haves som det rene
substantielle i Erkjendelsen, just fordi det oplyser den Intethed, hvormed alt,
hvad der nævnes som noget bestemt i Virkeligheden, og det ej blot Gjenstandene men
og Begreberne selv, er behæftet, og saaledes ligesom ryster de fixerede
Antagelser i deres Grundvold, i det samme tillige fører til at anvise dem deres
relative Gyldighed og befæste dem i denne.

\begin{center}Anmærkning og Tillæg.\end{center}

1.~Dialectiken viser sig her som en \emph{Konst}; først forsaavidt den bestaaer i
den Tankebevægelse eller Fremgang i Tænkningen, hvorved i Gjerningen Ideen gjøres
gjeldende med Hensyn til Noget, hvilket, om det end først skal bestemmes og
oplyses ved den, dog uafhængigt af den har fremstillet sig i sin Mangfoldighed og
sat sig fast, hvert for sig, med en Betydenhed, ja Absoluthed, hvilken, naar den
nu skal hjemles for den philosophiske Betragtning, der skal vurdere og bestemme
dens Adkomst, ikke kan bestaae sig. Men da intet grundigen kan vises i dets
Ugyldighed, uden at der tillige maa sees, om og hvorvidt det dog i sin Sphære
eller i sine Relationer er noget, der ej uden Grund er indtraadt eller er bleven
nævnet med det øvrige som noget særligt for sig, saa vil denne \emph{philosophiske
Tilintetgjørelseskonst}, der ej, som Skepsis, har et blot negativt Formaal, men
gaaer ud paa at tilintetgjøre alle den lavere Forstands Betragtningsmaader for at
sætte den højeste Fornufterkjendelse som den ene gyldige i Stedet derfor, tillige
nødvendigviis og væsentligen gaae i eet med den philosophiske videnskabelige
Constructions Konst.

Ligesom i den organiske Verden den idelige Vexel er en nødvendig Følge af Livets
uafbrudte Leven i sit Materiale (see Forfatterens Psychologie §.~5), men nu ogsaa
viser hen paa, og vidner om, Livets indre af ingen Form bundne Kraftighed,
saaledes er \emph{det relative} og det derfor ideligen anderledes bestemmelige,
aldrig i bestemte Udtryk og Sætninger udtømmelige, i al philosophisk Erkjendelse
en nødvendig Følge af det Absolutes Absoluthed og deraf, at det Andet, hvori det
Absolute træder frem, aldrig er det Absolute selv; men denne Relativitet uden Ende
viser da og hen paa dette Absolute, og af de dialectiske Tankebevægelser (i
sædvanlig negativ Forstand) maa derfor det Absolutes Erkjendelse altid gaae
bestyrket frem. Kan endog det selvsamme baade bekræftes og benægtes om det
selvsamme, (kan, for at tage simple Exempler, Materien baade siges at være noget
for sig og ikke at være noget for sig, kan det baade siges, at Begrebet om flere
Dyder har Realitet og ingen Realitet), saa have vi dog det positivt gyldige i det,
der gjør, at vi saaledes maae baade bekræfte og benægte, og, naar vi ret see,
hvorledes vi just, idet vi bekræfte, maae benægte og omvendt, eller og hvorledes
vi kun see Bekræftelsen i sin Gyldighed, naar vi tillige see Benægtelsens
Gyldighed, saa have vi just i denne Indseen Ideen eller det umiddelbare Absolute
under denne dens Fremtræden. Og en ligefrem Følge bliver nu at søge Foreningen i
en Deduction af Begrebet, i hvilken Begrebet skues paa sin \emph{relative} Plads
med sin relative Gyldighed, og den dialectiske Tankebevægelse viser saaledes
ligefrem hen paa, og fører til, den videnskabelige Construction.

2.~I Dialectiken gaaer det \emph{logiske} og det \emph{metaphysiske} i eet, og det
er ofte bleven bemærket, at i de Gamles Inddeling i Dialectik, Physik og Ethik,
den første staaer istedetfor baade den senere Tids Logik og dens Metaphysik eller
i det mindste dens Ontologie. En speculativ Logik (ej en blot iagttagende,
reflecterende og beskrivende logisk Formlære) sees i det dialectiske i sin
practiske Fremtræden i Undersøgelser, i hvilke givne Begreber gjennemtænkes og
bestemmes med Hensyn til det højeste Sande, hvilket højeste Sande er
Metaphysikens Gjenstand; saa at baade i denne Henseende, saa og forsaavidt en
speculativ Logik selv har dette højeste Sandes Erkjendelse til Basis, det
metaphysiske gaaer i eet med det logiske. Gjøres det, der i det dialectiske er i
levende Bevægelse, til Gjenstand for en særegen Opmærksomhed, saa ville logiske
Erkjendelser deraf kunne udspringe.

3.~Det dialectiske, i den almindelige Betydning, nemlig i sin negative Retning,
gaaer ud paa at bestride og ophæve, eller dog vise tilbage inden de behørige
Relationer og Indskrænkninger. Først vender det sig mod saadanne
\emph{Forestillingsmaader}, hvilke, under Menneskehedens Bestræbelse efter at
udgrunde og erkjende, monne have dannet sig og sat sig fast med en Gyldighed, der
ej kan tilkomme dem. Dernæst gaaer det ud paa selve de for Anskuelsen givne
\emph{Gjenstande}, eller paa de Mennesket i den umiddelbare reelle Anskuelses
Sphære nødvendige Forestillinger. Det højeste er endelig, at det vender sig mod
\emph{Begreberne selv}, de nødvendige Begreber nemlig, især de a priori erholdte,
og ogsaa gjennemryster disse saaledes, at deres Gyldighed kun gjennem en Indsigt i
deres relative Ugyldighed kommer til at bekræfte sig. Hvilket sidste ej heller kan
andet end blive en dialectisk Stræbens Meed, men dette er det, Platon, (see hans
Parmenides pag.\ 129), idet han lægger særdeles Vægt derpaa, ogsaa især finder
vanskelig.

\begin{center}\large\S.~15.\end{center}

Ideens Udførelse i en fuldkommen Fremstilling, hvori alt skues paa sit Sted og i
sin Betydning, tilmed i sin Sandhed, ifølge Ideen og i Kraft af Ideens
Erkjendelse, fører til et \emph{System}, og den videnskabelige Erkjenden bliver nu
en \emph{systematisk}. I Systemets Udførelse maa Ideen staae sin \emph{Prøve}, om
den nu kan vise sig overalt som det reelle og i det Enkelte kan yde og opfylde,
hvad den som Idee skal være i Stand til at yde og udrette. At al Erkjenden bør
være systematisk, er klart nok, da kun i et System alt kan sees paa sin relative
Plads og i sine Relationer, men alt netop kun er, hvad det er, og har sin Betydning
i sine Relationer og paa sin Plads.

Et saadant System bliver nu nødvendigviis udført i og ved \emph{Begreber} og
bliver derved selv et System af Begreber; thi da deri alt skal sees i sin
Sammenhæng og bestemmes med Hensyn til denne Sammenhæng saaledes, som det ifølge
denne Sammenhæng viser sig, saa bliver Erkjendelsen heelt igjennem en Erkjendelse
af Sammenhæng og af Gjenstandenes derved bestemte Realbetydninger og
Fremtrædelsesmaader eller Grundformer, saa at overalt ethvert Punct bestemmes
saaledes, som det bestemmes, ifølge en Indseen og Begriben; ja endog
Grunderkjendelsens Gjenstand, Ideen, maa nu (som allerede forhen er bemærket)
gaae ind med som Led, nemlig ypperste Led, i det Hele, ret som om den, om end
højeste Gjenstand, dog kun var en vis særdeles Gjenstand med, ligesom de øvrige og
ved Siden af dem. Men om Ideen end saaledes selv maa træde op som et Begreb, saa
maa den dog tillige sees som det altreelle, der gaaer igjennem det hele System, og
fremstiller sig i, og straaler frem af, ethvert Punct, saaledes som vi forhen have
søgt at vise det. Sees da nu det Evige dog som det reale Væsen i og under de mange
Grundformer, skues og haves Ideen som det altbestemmende og realiter overalt
virkende, der gjennem hiint System af Begreber fremstiller sig selv, i sin Eenhed,
men dog under mange Skikkelser, saa indlyser det, hvorledes i Videnskaben (ligesom
i den reelle Virkelighed) Ideen, dens Eenhed uskadt, kan træde frem i en
\emph{Fleerhed af Grundformer}, i hvilke den selv hver Gang viser sig selv i en ny
Skikkelse, og at der nu kan tales om en \emph{Fleerhed af Ideer}. Hvorledes Ideen
kun er en eneste og ej selv er flere, og dog igjen saaledes træder ud i flere, at
flere Ideer kunne nævnes, det bliver klart, naar Ideen sees som den, der
yderliggjør sig selv i et System, hvilket, saavidt det afledes af den, er dens
Udtryk, og hvori forsaavidt altsaa i enhver derved bestemt Form den selv maa
gjenskues, men gjenskues som i mange sig for sig hypostaserende eller
constituerende Virkninger og Principer med ligesaa mange dertil svarende
ejendommelige Individualisationer, altsaa som i mange Ideer. Hvilke da, hver i sin
Sphære, vise sig som det altreelle og absolute, hvorpaa alt kommer an, hvorfra alt
gaaer ud og hvori alt gaaer tilbage, og nu ogsaa træde ud i en \emph{Fleerhed af
Videnskaber}. Men denne en vis Idees Alrealitet og Absoluthed er dog ikkun en
saadan, der finder Sted med Hensyn til alt, hvad der hører under denne Idees
Sphære, saa at den kun er en relativ Absoluthed.

Ideens Leven viser sig da, ligesaavel under Erkjendelsens og Videnskabens, som
under den reelle Virkeligheds Form, som en \emph{organiserende}, der, idet den
bliver det bestemmende Princip for et mangfoldigt Indbegreb af sondrede, men dog i
deres Sondring og relative Modsætning nøje forbundne enkelte Erkjendelser, nu
ogsaa, idet den saaledes giver dette Indbegreb dets hele Skikkelse, tillige deri
selv tager Skikkelse. Det stringente, skarpe og strenge, der er i al sig
organiserende og sig tilmed i en ejendommelig Form paa en bestemt Maade
fastsættende Leven, kommer da og til Syne her.

\begin{center}Anmærkning og Tillæg.\end{center}

1.~Kan end Ideen vistnok træde levende frem i det frie \emph{judiciøse}, saa maa
dog ogsaa allerede derunder et System, i det mindste visse Grundtræk dertil, haves
i \emph{Sinde og for Øje}; og skal paa en nøjagtig Maade det Enkelte bestemmes og
tilfulde sees i sin Betydning, saa vil Overvejelsen derunder nødvendigviis gaae
ind i en omfattende Sammenhængs Undersøgelse til alle Sider, og det vil snart
blive fornummet, at kun en total Udførelse af alle de mange i Bevægelse kommende
Betragtninger til et vel sammenhængende Hele kan bringe til de nøjagtige
Bestemmelser, som allene kunne tilfredsstille.

Til ægte videnskabelig Stræben hører derfor ej blot Interesse og Respect for Sagen
selv, som skal udtales, eller for Erkjendelsens Indhold, men og for den Form, hvori
den skal udtale sig, og en omhyggelig, ja nøjeregnende Eftertænkning over, paa
hvilken Plads og i hvilken Sammenhæng hvert enkelt Punct skal fremstilles, for at
der kan være Bestemthed i det enkelte, Orden og Følge eller Conseqvents i det
Hele. En vis \emph{logisk Enthusiasmus} maa gaae frem af Enthusiasmen for Ideen og
ledsage den. Hiin er af samme Kilde som denne. For at Ideen kan nyde sin fulde Ret
og vise sig i sin fulde Kraft og Totalitet i Fremstillingen, derfor bliver ogsaa
Formen saa magtpaaliggende.

2.~Men kun som Ideens væsentlige og nødvendige Udtryk og Realvordelse er Systemet
det, det er. Derfor maa det frie judiciøse, saavelsom det dialectiske, gaae foran
Systemet og bane det Vejen. Letteligen fører ellers den systematiske Tendents til,
at der, for dog at bringe det mangfoldige, som er kommet i Bevægelse, de mange
isolerede Erkjendelser, som ere opnaaede, til en vis Eenhed og Heelhed, dannes
\emph{Meninger}, som uagtet fra først af hypothetiske, dog efterhaanden betage og
hilde. Med et Systems Uddannelse maa derfor ikke \emph{hastes}.

Systemet maa dernæst ogsaa i alle dets Dele vedvarende holdes i Live eller
\emph{levende Bevægelse} derved, at disse bestandigen skues, som de skues, eller
sættes, som de sættes, i Kraft af Ideen, der da under idelig Reproduction
fremkalder dem paa ny. Herved blive de mange enkelte Erkjendelser først fastholdte
saaledes, at de, bevarende sig i bestandig Friskhed, nu ogsaa hver Gang sees i
deres Sandhed, hvilken de overhovedet ikke have eller vise i deres Isolerethed
eller i deres blot systematiske Sammenknytning med de øvrige Led, men kun i deres
Sammenhæng med Ideen, som de, der ere medhenhørende til dennes levende totale
Udtrykkelse.

3.~To Ting ere her at befrygte og forebygge. For det første, det
\emph{Nebulisterie} (for at bruge et af Göthe i Konstneres Bedømmelse optaget
Ord), der blot holder Formen eller Systemet i en vis Taage, hvoraf det kun i
ubestemte Træk skinter frem i en vis Almindelighed, uden at det enkelte bestemmes
nøjagtigen, saa at der ingen fast Holdning bliver i det Hele. Lyst til at skue
Ideen strax i dens store Totalitet fører vel stundom endog til en vis Sværmen,
hvori kun visse Erkjendelser antage en nogenlunde staaende Skikkelse, men den
større Deel kun haves ganske løseligen, og store Huller udfyldes af en dunkel
Phantasering, paa det at Øjet dog kan fornøjes ved Skinnet af et Hele. Dernæst er
at vaage over, at Systematiseringen ikke skal ende i en \emph{Chrystallisation},
hvorved tilsidst kun et uddødt Substrat bliver tilovers, der engang har baaret en
Idee og været Udtrykket af dens levende Bevægelser, men nu ikke mere opfyldes
deraf. Skal Systemet blive ved at være det, hvad det i sin Sandhed er og tilmed
ogsaa bør være, saa maa det bevares under bestandige nye Formationer, hvori det
bestandigen gjenfødes, idet det da ogsaa tillige bestandigen omdannes og træder
frem med nye Bestemmelser, indtil det endog paa visse Udviklingstrin modnes til en
Metamorphose, hvori det udaf en vis indtil et vist Punct uddannet og saa længe med
Grund fastholdt Form gaaer over i en anden.

Dog ogsaa uden at et System saaledes chrystalliserer sig, endog om det under en
bestandig levende Indvirkning af, og Vexelvirkning med, den Grundtanke, hvoraf det
gik frem, holder sig i den meest levende Bevægelse, kan det let føre til
\emph{Eensidighed}, og saaledes hilde og binde, at, idet alt tilsidst kun sees i
dette ene Systems Lys, det nu ikke engang gaaer ind i de nye Betragtninger, der
kunde føre det til at gjenføde sig i en ny Skikkelse, ved at optage noget nyt i
sig i Henseende til de første Grundtanker. Og dertil skal dog altid holdes Rum
eller Aandsfrihed og Villighed aaben. Thi da vi dog kun skue \emph{det Ene},
hvorpaa alt kommer an, \emph{saavidt som}, eller i det Maal, i hvilket vi see det
som eet og alt i det tilsvarende Hele, saa maa Erkjendelsen ikke allene
bestandigen kunne udvide sig og blive mere inderlig og kraftig i Henseende til den
Mangfoldighed, som mere og mere gjennemtrænges af Ideen, men ogsaa maa derunder
denne selv bestandigen kunne vise sig ej blot i et klarere, men og i et andet Lys;
og saaledes maa Erkjendelsen just ved at udvide sig i den ene Retning, drives til
bestandigen at søge et nyt, nemlig et mere primitivt og mere til alle Sider
henvisende, \emph{første Begyndelsespunct}; saa at altsaa med Mangfoldigheden ogsaa
Grunderkjendelsen og Grundprincipet selv, d.\ e.\ den Maade, hvorpaa
Grundprincipet sees og fattes, maa kunne undergaae en bestandig \emph{Reformation}.
Og derfor maa da hiint frie umiddelbart \emph{judiciøse}, ligesaavel som det maa
gaae forud for Systemet og berede til det, saaledes og \emph{følge paa det}, naar
der nu gaaes fra Systemet til det mangfoldige Detail. Saa vigtigt et bestemt
System er, saa nødvendigt er det, indtil en vis Grad igjen at kunne befrie sig fra
Systemets Baand og lade Ideen bevæge sig frit, uafhængigt af den Form, hvori den
har søgt at komme til en bestemt Skikkelse. Dertil hører, at vi, i levende
Vexelvirkning med Phænomenet eller Problemet selv, lade dette udøve sin Magt over
os til at lede os, hvorhen det vil, og lade det selv fremkalde den Tanke hos os,
hvorved der skal bringes Lys ind deri, hvilket iøvrigt vistnok forudsætter, at der
er et Fond af Tanke, hvoraf den nye, der nu skal svare til Tingen saaledes, som den
selv viser hen derpaa, kan fremspringe. Sand Interesse og Kjerlighed til Sagen, og
Lyst til at udgrunde den selv, som den er, vil føre til at lade det
Vexelvirkningsforhold imellem det, vi have og ere, og den nye Enkelthed, vi skulle
betragte, blive som det bør være, saa at et frit Blik lærer, hvad enten det nye
givne ligefrem maa gaae ind i, og ligesom assimileres af, Systemet, saa at dette
derved ej kan bringes til nogen Omdannelse af Betydenhed, eller om dets Indflydelse
maa blive den, at Systemet derved kommer til at bevæge sig tilbage i sig selv og
til en vis Grad at løsne sig ud af sin faste Tilstand, for, idet det assimilerer
sig det nye, at gaae fremad i sin egen Udviklingsproces, og ifølge denne Optagelse
selv at komme til en Omdannelse, være det og kun i een enkelt Henseende. Især kan
Bestræbelsen for at bringe det erkjendte til andre, naar det bliver os ret
magtpaaliggende at overbevise dem, og vi derfor nødes til at gaae ind i deres
Tankegange, virke imod den Eensidighed, et Systems Uddannelse let kan føre til, de
Baand, hvilke det kan lægge paa den frie Tænken. Overhovedet kan Vexelvirkning med
andre og det flere forskjellige, tænkende Væsener, være det under en gjensidig
levende Meddelelse i Omgang, eller under Studiet af forskjellige Forfatteres
Værker eller af Videnskabens Historie og historiske Udviklingsgang, ikke andet end
være af Vigtighed i denne Henseende.

4.~Et fuldkomment (nemlig et til en vis Grad af Fuldendelse kommet) System lægger
sig \emph{practiskt} for Dagen derved, at ethvert Begreb og enhver Gjenstand for
Undersøgelse nøjagtigen sees i sine Forbindelser med alle øvrige ved dens
Betragtning tillige i Bevægelse kommende Begreber og Gjenstande, og at Ideegangen
fra ethvert Punct kan føres lige op til det højeste, fra hvilket som det faste
Udgangspunct alt gaaer ud, og kan føres til alle Sider, saaledes, at det
betragtede enkelte Punct derved aldeles kan sættes i sit rette Lys og stilles fast
i sine Relationer og med sin Betydning. Forsaavidt nu alle Hovedpuncter i en
Videnskab sees i denne Klarhed og Sammenhæng, kan allerede et heelt System i sikre
og mangfoldige til een Grundform henvisende Træk lyse frem af enkelte forekommende
Spørsmaals Behandling og af Maaden at dømme og bestemme Begreberne paa i det Hele,
uden at det dog endnu haves udtalt i den systematiske Form. At levere Systemet som
System i en udvortes videnskabelig \emph{Fremstilling} kan selv, om det allerede
saaledes practisk besiddes, have sine store Vanskeligheder. Især vil den Orden og
Følge, hvori alt skal foredrages, kunne volde mange Betænkeligheder og Tvivl. Det
oplysende og begrundende skal dog staae foran, men da alt enkelt, som i et
organiskt Hele, hviler i Ideen, men denne igjen kun haves i det hele Systems
Udførelse, og kun indlyser i sin fulde Sandhed, naar den lyser frem af det hele
videnskabelige Indbegreb af Erkjendelser, som det sande i disse, saa hviler enhver
enkelt Erkjendelse egentligen i det hele System selv, hvilket dog, naar det skal
fremstilles, kun kan opbygges ved de efterhaanden til hinanden knyttede enkelte
Erkjendelser; og ligesom Tænkeren, der vil udtale Systemet, ved ethvert Punct
allerede har det hele System for Øje og seer det enkelte i det Hele, saa skulde
ogsaa i Fremstillingen ved de enkelte Puncter, tilmed allerede ved dem, hvormed de
begyndes, den hele Totalerkjendelse træde levende frem og stilles Betragteren for
Øje. Men dette kan dog ej skee i en Fremstilling, der kun kan gives Led for Led.
Noget maa tages først og stilles foran. Ja dette kan ej engang være
Grunderkjendelsen, hvilken dog først maa staae op for vor Øje og begrunde sig som
Grunderkjendelse, inden den kan opstilles som saadan, ej heller kan den, som den,
uden hvilken det Hele ikke sees i sit sande Lys, komme tilsidst. Der maa ledes op
til den, og da maa den træde ind som den, hvorfra nu hele Systemet skal gaae ud.
Ogsaa hvad de mange af Grunderkjendelsen udspringende enkelte Erkjendelser
angaaer, der blive Systemets mangfoldige Led, saa vil det ej sjeldent vise sig, at
den ene allerede maa haves og fremstilles for blot at kunne \emph{komme til den}
anden, men at denne dog igjen maa være \emph{fremstillet}, inden hiin kan sees i
sin \emph{fulde Sandhed og Betydning}. Ja ej blot enkelte Erkjendelser og
Sætninger, men \emph{hele Ideegange} (endog hele Videnskaber) kunne staae i et
saadant Forhold til hinanden, at, medens dog hver af dem fordrer at udføres og
fremstilles for sig i sin Retning og i sin Sphære, dog den ene ej kan, hverken
fattes tilfulde eller indlyse aldeles, uden at forudsætte den anden, og det
gjensidigen. I slige Tilfælde fører den systematiske Anordning store
Vanskeligheder med sig.

Det bliver da nødvendigt \emph{at tage fra en vis Side}, hvad der dog i sig selv ej
gaaer ud fra denne ene Side eller mere hæver denne Side frem hos sig end de andre.
Men ligesom deraf sees \emph{Mueligheden af flere Fremgangsmaader eller Methoder},
der kunne gjøre lige Fordring paa Grundethed og Værd, saa indlyser deraf ogsaa,
hvorledes i samme Videnskab \emph{flere videnskabelige Systemer} kunne opstilles
ved Siden af hinanden, uden at noget af dem kan gjøre Krav paa Enegyldighed eller
en afgjort Fortrinlighed i alle Hensender, og det, om der end i Henseende til den
videnskabelige Gehalt i det Hele hersker Overeensstemmelse.

5.~Ægte System er kun det, hvortil vi komme fra Ideen og ved Ideen, idet denne
betragtes, som den udtrykker sig i et vist Indbegreb af Detail, hvilket Details
Omfang og Art da bestemmes ved Ideen selv. Men ogsaa naar et Indbegreb af
Enkeltheder er os givet for sig (paa en empirisk Maade) og bliver Gjenstand for en
Iagttagelse, der gaaer ud paa at skue det, som det nu engang selv frembyder sig,
uden at det endnu erkjendes i Ideens Lys, vil Bestræbelsen altid blive at ordne
dette Indbegreb af Iagttagelser eller iagttagne Former og de derom dannede (blot
formale) Fælledsbegreber, eller Slægts- og Artsbegreber, i et System til en
omfattende Oversigt. Men ogsaa heri kan da og bør den philosophiske Stræben vise
sig \emph{deels} i den Maade, hvorpaa et vist Ideal om en saadan Classification med
Hensyn til det, den ifølge sin Hensigt skal yde, bestandigen haves for Øje,
\emph{deels} i den Maade, hvorpaa der ved dens Udførelse arbejdes en højere
Betragtning i Haanden. Jfr.\ §.~6. Anm.~2 i Slutningen.

6.~Et Systems Nødvendighed i Videnskaben følger saa ligefrem deraf, at intet
Enkelt kan sees i sit rette Lys, uden nøjagtigen at sees i sine Forhold og sin
Forbindelse med alt dermed sammenhængende, at derfra ingen Undtagelse kan finde
Sted. Een Videnskab synes imidlertid dog at gjøre Undtagelse, en Videnskab, som
dog, naar der er Tale om dens \emph{Fremstilling}, mere betragtes som en Konst:
\emph{Historien}. Men Begivenheder sees dog kun, ligesom alt andet, i deres rette
fulde Lys, naar de sees i deres, ofte mangesidige Sammenhæng med andre deels
samtidige, deels for- og efterfølgende, og Charakterer maae ogsaa sees i deres
Tider og paa deres Pladser. Saa at de samme Grunde som ellers fordre det
systematiske i Erkjendelsen ogsaa klarligen vise sig her, men hvad der ellers i
andre Videnskaber har en anden Form og kaldes System, er her den \emph{historiske
Composition}.

\begin{center}\large\S.~16.\end{center}

Et System kan ikkun finde Sted, en Idee ikkun udtale sig, i et mangfoldigt
Indbegreb af enkelte Erkjendelser, der, som det videnskabelige Heles Dele eller
Led, ogsaa kunne betragtes hver for sig, og forsaavidt tages som specielle
Erkjendelser. De udgjøre da det til al Videnskab væsentligen medhenhørende
\emph{Detail}, hvis nøjagtige Studium er af højeste Vigtighed. Kun naar et Begreb
sees gjennemført, en Idee sees som det reale Væsen, i Gjenstandenes hele
Mangfoldighed indtil dens mindste Dele eller fjærneste Arter, saa at alting i samme
begribes ifølge Begrebet og sees i Ideens Lys og som dens Udtryk, eller dog som
noget, hvori ogsaa Ideen er det reale Væsen, er selve Begrebet og Ideen fattet
tilfulde. Først i det mangfoldige Concrete viser sig Ideen i sin fulde Kraftighed,
som det i Sandhed altbestemmende, og Ideens Erkjendelse gjennemgaaer først i
Detaillets Gjennemgranskning sin fulde Prøve. I Detaillets nøjagtige Studium viser
sig den til al Videnskabs Uddannelse fornødne \emph{Iagttagelsesaand}.

Men skulde nu vel og et særligt Studium af Detaillet for sig kunne tænkes,
\emph{forud for}, eller dog \emph{adskilt fra}, Ideen og Begreberne, og nu kunne
træde disse imøde? Hvorledes skulde vel et saadant være mueligt? Forsaavidt de
specielle Erkjendelser skulle være ægte videnskabelige, hvorved det Enkelte
erkjendes i sin Sandhed, tilmed være intelligente Erkjendelser, kunne de kun haves
som Følger af, og i Kraft af, Idee og Begreb; de ere da netop disses
Fremtrædelser i, eller Bestemmelser med Hensyn til, det Enkelte, og Detaillets
Studium kan da ej være andet end den fuldkomne Udførelse af selve Ideens og
Begrebets Erkjendelse; det bestaaer da netop i Grunderkjendelsens Forfølgelse
indtil alle de mange forskjellige Enkeltheder, eller særdeles Tilfælde, som ere at
henføre under Grunderkjendelsen og skulle oplyses og bestemmes ved samme.
Iagttagelse kan altsaa her ikkun paa en \emph{indirect} Maade være frugtbringende,
idet den deels bemærker og fremdrager de specielle \emph{Spørsmaal og Problemer},
der nu vente deres Opløsning af den intelligente Erkjenden i Kraft af Begrebet,
deels henvender Opmærksomheden paa de mange Tilfælde, hvori det erkjendte kan
finde \emph{Anvendelse}, en Anvendelse, der dog nu maa skee ifølge den
andenstedsfra kommende Indsigt. Forsaavidt derimod, saaledes som i Empiriens
Sphære, forud for al klar Indsigt og intelligent Erkjendelse, Resultaterne af Idee
og Begreb kunne være givne, som udtalte Resultater af en Intelligents, der dog
dermed endnu ej strax er given som saadan, fordi den endnu kun har fremstillet sig
i sine Følger, kan et særligt Iagttagende Studium af et givet Detail paa en mere
\emph{direct} Maade føre til ægte videnskabelig Erkjendelse, idet Tingenes
Grundformer, deres Sammenhænge og Forhold, og de i dem udtrykte Begreber og Love
nu eftergranskes saaledes, som de i Tingene selv give sig tilkjende. Og her bliver
da, uagtet det ogsaa her er en intelligent Erkjenden, hvorpaa det kommer an, da en
i de givne Ting sig fremstillende Intelligents skal udgrundes, \emph{Detaillets
særlige Studium} af saa stor Vigtighed og Væsentlighed.

\begin{center}Anmærkning og Tillæg.\end{center}

1.~Hvorledes \emph{Iagttagelsesaand}, ja Empirie, kan være af Vigtighed for
saadanne Videnskaber, hvori intet skal eller kan abstraheres af Erfaringen,
f.~Ex.\ i \emph{Mathematik}, saa og i \emph{Philosophie}, endog den meest
speculative, vil af det fremsatte blive klart. Vel kunne en mathematisk Sætnings
Følger og dens Anvendelser i forskjellige specielle Tilfælde ikkun \emph{sees} ved
at \emph{indsees} ifølge Sætningen selv, og kunne tilmed ikke simpelhen opfattes og
iagttages som noget umiddelbart givet eller forefundet. Men i \emph{Opmærksomheden}
paa alt det, der vil kunne udledes af en Sætning og paa de forskjellige muelige
Combinationer, i hvilke det erkjendte vil kunne sees under ejendommelige
Modificationer eller i særlige Former, kan en vis mathematisk Iagttagelsesaand vise
sig og fremme Videnskabens Uddannelse meget. Empirien dernæst opgiver
Mathematikeren blot Problemer; men en, allene under Forudsætning af et mangfoldigt
Studium af empiriske Gjenstande muelig, Agtsomhed paa deslige Problemer med en
nøjagtig Bestemmelse af de Spørsmaal, hvis Besvarelse Naturforskeren venter af
Mathematikeren, har dog ej sjeldent fremkaldet ellers neppe opstaaede mathematiske
Undersøgelser, ja betydelige Udvidelser af de mathematiske Discipliner, og
saaledes kan da endog en i Empiriens Sphære sig bevægende Iagttagelsesaand blive
frugtbringende for en aldeles ikke empirisk Videnskab. Vi ville blot minde om et
af de nyeste Exempler: \emph{Chrystallographien}.

Ogsaa i Philosophien gjelder det samme og lignende Betragtninger lade sig anstille
om samme. De Spørsmaal, hvilke Erfaringen kan gjøre f.~Ex.\ til den philosophiske
\emph{Ethik}, en Videnskab, hvis Sætninger, lige indtil de allermeest specielle,
eftersom ogsaa disse skulle udtale et Bør, ikke kunne hentes fra Erfaringen,
fremkalde ofte Overvejelser, der blive frugtbringende for Ethiken. Ja den
tilbagevirkende Indflydelse af deslige, ved Agtsomhed paa empirisk givne Forhold og
deres mangelunde Modificationer opstaaede, Betragtninger og Undersøgelser,
strækker sig endog tilbage til de første Principer eller til hele
Grunderkjendelsen, der vistnok ikke beroer paa nogen Empirie, men til hvis dybere,
bestemtere og mere omfattende Fattelse vi dog ofte anlediges, ja vejledes ved
Opmærksomhed paa, og Eftertænkning over, de mange empiriske Gjenstande, i hvilke
den skal bringe Lys og Indsigt, eller med Hensyn til hvilke den skal udøve sit
Herredømme.

Allerede i denne Henseende kan altsaa en vis \emph{philosophisk Iagttagelsesaand}
vise sig. Men ogsaa endnu i en ganske anden Henseende. Den philosophiske
Erkjendelses Gehalt gives os endnu paa en anden Maade og ad en anden Vej, end i
den philosophiske Erkjendelses (videnskabelige) Form og ad den videnskabelige Vej:
nemlig igjennem en vis umiddelbar Fornemmelse eller Følelse, \emph{en
Sandhedsfølelse} (hvortil vi ogsaa henregne den æsthetiske og den moralske
Følelse). Heri viser sig det sig i os forkyndende Uendelige og Evige, Ideen, i
sine mere \emph{reelle} Virkninger, refererende sig til vor Subjectivitet, vort
Velbefindende eller vor personlige Stilling: heraf har den uafhængigt af
videnskabelig Betragtning sig forkyndende umiddelbare Tro, ligeledes en vis
umiddelbar Tact eller Skjønsomhed (det foromtalte umiddelbare judiciøse) sin
Oprindelse, og kun derved er samme muelig. Men ej i sin Almindelighed, kun i sin
specielle Yttring med Hensyn til noget individuelt og concret viser sig den almene
og almeengyldige Idee i denne umiddelbare, sig selv endnu ej aldeles klare,
Sandhedsfølelse. Deri er da noget, hvormed Ideens videnskabelige Udvikling, naar
den udføres, ofte kan og vil mødes. Den mødes heri med sig selv, med sine egne
specielle Sigforkyndelser, men i en ganske anden Form, end i den videnskabelige
Afledning og Slutning. Og saa er da heri et \emph{Detail}, der kan blive Gjenstand
for en særegen Agtpaagivenhed, og med Hensyn til hvilket Ideen (den videnskabelig
udtalte Idee) kan undergaae en Prøve og finde en Stadfæstelse.

2.~Al ægte og levende Kjerlighed til Ideen og Glæde ved samme føder af sig en
\emph{Interesse og Glæde ved Detaillet}, en Lyst til at forfølge det Enkelte
indtil dets mindste Dele eller særligste Modificationer, for ret tilfulde at skue
og tilegne sig det intelligente igjennem det empiriske, det almene igjennem det
specielle; ja denne Lyst til at forfølge det Enkelte hører væsentligen med til at
vidne om et sandt Kald til en vis Videnskab; hvorimod det at ringeagte Detaillet
vidner om Mangel af ægte Sands og Følelse for Ideen, til hvis Væsen det netop
hører at træde frem som det altbestemmende i sin Mangfoldighed, og hvis Kraft kun
lader sig tilsyne i dens Virkninger, saa at vi kun kunne knytte os til den og skue
den, ved at knytte os til det, som bærer dens Præg og Udtryk, og skue den i dette.
Naar vi spørge os, hvorledes dog en genial Maler af den nederlandske Skole har
kunnet finde Glæde ved saa troligen at iagttage og studere Lysets Reflectioner paa
en Stol eller Kobberkjædel under visse Belysninger, for troligen at kunne gjengive
den i sit Malerie, saa ligger Svaret nær: netop Lyset, det almene alt oplivende
Lys, glædede det ham at skue og forfølge i de utallige Nuancer af dets
Reflectioner paa Gjenstandene; det almene Lys var det, der oplivede ham i dette
Studium, og var det, hvilket han vilde fremstille i nogle af dets utallige
Virkninger. Og naar saaledes et tænkende Hoved med Sands for det højeste ej
sjeldent har indskrænket sig til en speciel Gjenstands (et Insects eller en
Plantearts) nøjagtige og trolige Studium, saa var det, fordi det Hele var ham
magtpaaliggende og fordi han saae, at det for dette Heles Skyld var nødvendigt
ogsaa at studere denne enkelte Gjenstand. Og naar han studerede den enkelte
Gjenstand med den Aand og Tendents deri at levere et Bidrag til det højeste og
almeneste, kom han derved ogsaa selv til at studere det almeneste i den enkelte
Gjenstand og derunder at føle sig i levende Berøring med den hele almene,
altomfattende, alt i sig bærende Natur, og kunde nu ogsaa i den enkelte Gjenstands
grundige Fremstilling levere noget, der, ligesom det ej kunde komme i Stand uden
et mangfoldigt Henblik til den hele øvrige Natur, nu ogsaa kunde kaste et Lys hen
over en vidtomfattende Kreds. Ogsaa det meest sig begrændsende Studium, som vi i
Anm.\ til §.~1. allerede have sagt, har sin Muse.

\begin{center}\large\S.~17.\end{center}

I den fuldendte Videnskab træder Ideen i og formedelst et System frem som den, der
oplyser og bestemmer et Detail, være dette sidste nu givet som et Indbegreb af
Gjenstande, der skulle udgrundes, eller som Problemer, der skulle løses. Alle
trende høre væsentligen sammen, og maae under Bestræbelsen efter fuldkommen
Erkjendelse og under Videnskabens Uddannelse alle haves for Øje, og det alle tre
paa eengang og med hinanden. Vel er Ideen, den altoplysende og altbestemmende
Grunderkjendelse, i sig selv det første, men den haves dog først, ej blot i dens
fulde Lys, men og paa den rette Maade, naar den viser sig i det Mangfoldige, hvori
den er det altbestemmende og altoplysende; og saaledes maa, for at Ideens
Erkjendelse kan antage en bestemt Skikkelse, endog kun saavidt, at, den kan fattes
i et saadant Udtryk, at en videnskabelig Afledning kan begynde at gaae en fast
Gang, allerede Detaillet indtil en vis Grad være studeret, Systemet udkastet. Men
dog vilde atter, hvis ej Ideen allerede var sprunget frem for Blikket og fattet
med en vis Grad af Bestemthed, denne Systemets foreløbige Udkastelse og det
Studium af Detaillet, der kan blive frugtbringende for Ideens Erkjendelse, ej være
mueligt. Og hvad det indbyrdes Forhold imellem System og Detail angaaer, da staaer
begges Studium i den samme gjensidige Afhængighed; thi det Enkelte kan ej studeres
saaledes, at derved kan arbejdes hen til et System, uden at en vis Totalsammenhæng
derunder allerede svæver Betragteren for Øje, men Systemets Udkastelse forudsætter
igjen et Overblik over alle de enkelte Led og Dele, der i Systemet skulle sees i
deres Sammenhæng, stilles paa deres Pladser, og sees i deres relative Betydning.

Der kunne derfor tænkes flere Veje, ad hvilke, eller Maader, paa hvilke, man kan
komme til Maalet i en Videnskab, og flere kunne og maae under Granskningen og
Henarbejdelsen forbindes med hinanden. Saaledes bliver da, foruden Idee, System og
Detail, endnu i enhver Videnskab, med Hensyn til det tænkende Væsens Stræben, at
lægge Mærke til \emph{Methoden}. Deels maa der være en vis Methode i det Studium,
hvorved vi selv arbejde os frem, deels viser en saadan sig og i at udtale det
erkjendte i dets Sammenhæng og meddele det til andre. I sidste Henseende staaer
Methode og System i nær Forbindelse med hinanden (see §.~15. Anm.~4.); imidlertid
kan man dog i at meddele det samme System, den samme Erkjendelse eller Antagelse i
Henseende til alle Delenes Betydenhed, Sammenhæng og deres Forhold i et vist
videnskabeligt Hele, gaae forskjellige Veje; man kan fra forskjellige Sider og paa
forskjellige Maader føre andre ind i den videnskabelige Bygning, uden at denne
selv tænkes anderledes. Og saaledes bliver der altid Forskjel imellem System og
videnskabelig Foredragsmethode.

Gjøres i en Videnskab Maaderne, hvorpaa man kan gaae tilværks i at naae sit Maal,
og de Puncter, man herved har at iagttage, til Gjenstande for en særegen
Reflection, saa bliver Betragtningen en \emph{logisk}; enhver Videnskab kan altsaa
give Anledning til \emph{særegne logiske Overvejelser}, i hvilke det almindelige
logiske drages saaledes frem for Opmærksomheden, som det viser sig med Hensyn til
de særdeles Granskninger og Fremstillinger, hvilke denne Videnskab fordrer.

\begin{center}Anmærkning og Tillæg.\end{center}

1.~I Almindelighed er Forskjellen imellem den \emph{analytiske Methode} og den
\emph{synthetiske} at bemærke. Til Systemet, der ligesom udgjør den videnskabelige
Erkjendelses Midte, og giver den Fasthed og Soliditet, kunne vi bevæge os hen fra
to Sider. Ægte videnskabelig Fremgangsmaade maa, som den, der gaaer ud fra Ideen
og afleder et systematisk Hele deraf, blive \emph{synthetisk}; de samme
Betragtninger, der lede til at indsee Nødvendigheden af et System, vise ogsaa den
synthetiske Methodes Væsentlighed. Men til denne gjør først en indtrængende og
omfattende Betragtning af Detaillet, altsaa en \emph{Analyse}, tilstrækkelig
moden; og det ej blot i empiriske Gjenstandes videnskabelige Studium og
Fremstilling, men og i de rene philosophiske Videnskaber, i hvilke en Analyses
Muelighed indsees af forrige Paragraph. Dog maa det vel bemærkes, at Analysen kun
da har ægte Værd, naar den bestandigen sigter hen til et System, tilmed en muelig
synthetisk Fremstilling af et Hele, ligesom al grundig Analyse ogsaa i sig selv
allerede bærer Spiren til en saadan i sig. At skue og besidde Ideen som det alt
gjennemtrængende Lys, derpaa beroer dog den højeste Erkjenden, men saa maa og en
Deduction nødvendigviis gaae frem af den. Imidlertid maa dog derunder det
Mangfoldige, hvilket skal oplyses, bestandigen haves for Øje, for at den
videnskabelige Afledning nu ogsaa virkelig kan føres ned dertil og komme an
derved: det Territorium, der skal betragtes og overskues fra det højere
Standpunct, og hvori vi fra dette skulle orientere os, ja hvilket vi først ret
skulle lære at kjende ved denne Betragtning fra oven, maa dog være gjennemløbet
først nede i det selv og i dets enkelte Partier. Til at forebygge Eensidighed er
en foregaaende mangfoldig og nøjagtig Analyse saa vigtig. Ogsaa kan, naar en
Ideegang fra oven har ledet ned til et bestemt Punct og ført til et bestemt
Begreb, nu tilsidst kun et Blik paa Erfaringsverdenen selv lære, i denne at finde
det, der svarer til dette Begreb; men da maa denne givne Virkelighed i Erfaringen
først være nøjagtigen iagttaget, tilmed gjort til Gjenstand for Analyse. I blot
empiriske Videnskaber er Analysen Studiet af Bogstaven i Naturens Bog, gjennem
hvilken man dog først kan komme til \emph{dens Aand}, dens Tanker. Men ligesom man
dog kun ret har studeret og fattet en anden Bog, naar man saaledes er trængt ind
til dens Aand og intelligente Indhold, at man paa en fri Maade kan udtale dette,
gaaende ud fra dens Grundtanker, og kan i Kraft af disse selv udføre det Hele igjen
i en ligesom reproduceret Fremstilling, saaledes er da og Naturen først studeret
til Grunde, naar Mueligheden af en synthetisk Fremstilling er opnaaet. I rene
philosophiske Granskninger er Analysen egentligen selve \emph{Ideens Studium},
(see §.~16. Anm.~1.) saaledes nemlig som dens Lys reflecterer sig paa en enkelt
speciel Gjenstand, hvorpaa dens Kraft, Betydenhed og Fylde ej sjeldent meget
bestemt kan vise sig og iagttages, førend den har viist sig klarligen i og for
sig, eller i det videnskabelige Hele. ---

Paa en hældig \emph{Forbindelse af analytisk og synthetisk Fremgangsmaade} kommer
derfor i Studiet saavel som i Fremstillingen saa meget an.

2.~Bedst uddanner sig Methoden i en Videnskab af denne selv og dennes egen
historiske Udviklingsgang. Enhver Gransker finder Videnskaben uddannet til en vis
Grad og i en vis Form, og ingen kan med Kraft og Virkning føre den videre, uden at
gaae ud fra det Punct, paa hvilket den staaer i hans Tid. Kan Geniet end see langt
ud derover og udkaste en Sæd til nye Udviklinger for Aarhundreder, saa maa dog og
Geniet udtale det nye i en bestemt Relation til det bestaaende. Ere vi alle
bestemte Led i Menneskehedens hele Udviklingsrække, opstaae vi paa bestemte
Puncter i det Menneskehedens Liv, der udtaler sig og lever hen gjennem
Aarhundrederne, saa maae vi og gribe ind i dette Liv og bidrage til det store
Værks Fremme fra det Punct af, paa hvilket vi ere satte. Ellers vil Sæden, vi
kunne kaste ud, ej svare til den Jordbund, hvori den skal falde, Stødet, vi give,
ej vække den Reaction, eller fortsætte sig i den Forplantelse, uden hvilken det ej
kommer til at føre til det, det skulde. Vistnok skal Aanden, der skal tale i os og
gjennem os, være af Sandheden allene, og denne skulle vi ej lade bindes ved nogen
Tidsalders bestaaende videnskabelige eller overhovedet aandelige Tilstand. Men
idet den just skal føre ud af den bestaaende Tilstand for at føre videre fremad
under bestandige til sine Tider indtrædende Gjenfødelser af det Ydre, hvori den
fremtræder, maa den Virksomhed, hvorved den skal formaae dette gjennem Individernes
Bestræbelser, netop for at kunne føre til Gjenfødelser og Omdannelser, referere
sig paa de bestemteste Maader til det bestaaende.

Men i det bestaaende kan kun den Tænkers Virken gribe fuldkommen ind, der kjender
det tilfulde, men kjende det tilfulde kan kun den, der har betragtet det i dets
Opstaaen. Og saaledes bliver i enhver Videnskab et \emph{Studium af Videnskabens
Historie} af saa høj Vigtighed.

\begin{center}\large\S.~18.\end{center}

I Videnskaben træder Ideen, paa en ideel og universel Maade, ud i en ydre Form og
kommer derved til ydre Realitet eller Tilværelse (Existents). Men i og gjennem
Videnskaben kommer den tillige til at leve i og for en Bevidsthed og i personlige
Væsener; og derved kommer den først til den \emph{fuldkomne ydre Tilværelse}, og
træder nu først frem som en Magt og Vælde i den reelle Virkelighed. Her, i Sjælene
selv, see vi først det Stof eller Materiale, der ved Ideen skal beaandes og
levendegjøres, saa at frie personlige Individualiteter blive de Organer for Ideens
levende Fremtræden og Virken, hvori den skal tage Skikkelse og komme til at boe.

Dertil stunder og stræber da det i Ideen sig forkyndende Uendelige og Evige selv,
paa det at det kan komme til et Rige i Menneskeheden. Og derfor forkynder det sig
selv og \emph{arbejder sig frem}; ikke blot, at det fremstiller sig af og ved sig
selv, det bevæger ogsaa, griber og hentager. Heri bestaaer det \emph{geniale},
hvorfra al højere Erkjenden nødvendigviis gaaer ud.

\emph{Anm.} Hvad der herved endvidere kunde være at bemærke, er deels allerede
forekommet i det foregaaende, deels kan det læses i Forfatterens Psychologie i
§.~22. og Anmærkningerne dertil, hvor Genialiteten udførligen er omhandlet.

\begin{center}\large\S.~19.\end{center}

Skal Ideen komme til at leve som saadan i Bevidstheden og lyse i og ved sin egen
Sandhed, saa kan den ej komme frem og virke paa en instinctmæssig Maade, men maa
træde saaledes ind, at det sig selv bevidste Individ med Klarhed siger sig selv,
hvad det er, der taler i det, og at det nu, med Agtsomhed derpaa og Ærefrygt
derfor, lader det faae Rum og Magt hos sig, fordi det selv erkjender det for noget,
det maa aabne sig for; saa at altsaa Individet, om det end vistnok kun har det,
forsaavidt det er grebet deraf, dog ej er grebet blindthen, men med selvbevidst
Klarhed, med Indsigt og Frihed, lader sig hentage og bestemme. Er end
Overbeviisningen væsentligen en uimodstaaelig Grebethed, saa er den dog ogsaa en
villig og den er kun en ægte Overbeviisning, naar den er begge Dele i eet. Vi maae
selv see Sandheden i dens Sandhed, naar en sand Erkjendelse skal indtræde.

Ideen refererer sig altsaa nødvendigviis til en \emph{Selvbevidsthed og Frihed} og
sætter, idet den opfordrer samme, en \emph{Selvtænkning i Bevægelse}, for igjennem
dennem at komme til den Leven i Individet, hvortil den skal komme. Vi have allerede
forhen (§.~6. i Anm.) havt Anledning til at bemærke dette.

Ideen forkynder sig derfor ej blot i et Noget, som fremstiller sig af sig selv for
Bevidstheden, men og i en \emph{Drivt og Tilskyndelse}, nemlig en Tilskyndelse til
ved egen Opmærksomhed og Eftertænkning at drage hiint sig fremstillende Noget frem
for Bevidstheden, og nu at opnaae, hvad Tilskyndelsen viser hen paa og kræver:
overhovedet en Tilskyndelse til noget, som nu \emph{hos Individet selv}, eller i
den sig selv bevidste personlige Leven, \emph{skal paafølge}. See §.~4. Anm.~3.

Ideen fremkalder altsaa, for at komme til at træde frem og leve som Idee,
\emph{nødvendigviis sin Modsætning}, og bringer selv det Andet, hvori den skal
aabenbare sig i sin Sandhed, den \emph{selvbevidste personlige Leven}, til at
yttre sig som saadan og nu at træde Ideens egen sig fremarbejdende Indflydelse
imøde.

\emph{Anm.} Hvad der her allerede kunde paatrænge sig os til nærmere Overvejelse
og videre Udførelse, og ikke allerede forhen er kommet frem i Anmærkningerne,
bringe vi beqvemmere i Forbindelse med de Betragtninger, hvilke de følgende §§.
ville fremkalde.

\begin{center}\large\S.~20.\end{center}

Selvtænkningen sættes i Bevægelse ved en \emph{Spørgen}, og det --- eftersom
Friheden skal blive sig bevidst som saadan, paa det at en villig Hengivenhed kan
følge, --- en Spørgen, hvis Tilfredsstillelse ikke strax gives. Men et Spørsmaal,
der lader sig fornemme som saadant, et Spørsmaal, hvis Svar først skal
\emph{søges}, sætter ind imellem \emph{Tvende}, af hvilke \emph{Eet skal
foretrækkes}, medfører altsaa \emph{Mueligheden af Tvivl}, og, idet Individet
finder sig som det, der kun ifølge sin egen Bestemmen og Gjøren kan komme til den
Forestilling, i hvilken det skal bringe den attraaede Erkjendelse til Bevidsthed
hos sig og skal have og besidde den, kommer Tvivlen ogsaa til Virkelighed.

En \emph{Antagelse og Ikke-Antagelse}, eller, saafremt det virkelig er det Sande,
som forkyndte sig, (hvilket Individet dog først skal komme til Forvisning om), en
\emph{Erkjendelse og Ikke-Erkjendelse} træde nu imod hinanden. Men kommer det da
nu til Erkjendelse, saa faaer denne en ganske anden Charakteer, Kraft og
Betydenhed, end om Tvivlen og den deraf følgende Kamp og Prøve ej var
igjennemgaaet; det til Erkjendelse bestemmende, og det, hvortil det bestemmer,
ere nu traadte ud fra hinanden, og Erkjendelsen bliver nu først fuldkomment en
Erkjendelse af Sandheden i Kraft af dens Sandhed, en Erkjendelse med Indsigt og
fuld Forvisning. En Bekræftelse, der gaaer frem af en Negations eller Oppositions
Benægtelse, altsaa af en dobbelt Negation, er vistnok en ganske anden, end en
ligefrem Antagelse, i hvilken ingen Muelighed af Negation endnu har rørt sig. Uden
denne Benægtelse af Nægtelsen eller Oppositionen kommer (som Franz Baader med
Henviisning til Hegel siger: über die Ekstase, zweites Stük, S.~6) Livet ej til
\emph{at begrunde sig}, tilmed ej heller til at oplade eller at opløfte sig. Men
uden at hiin Nægtelse opvækkes, er det (som han føjer til) vistnok ikke mueligt
igien at benægte den.

\begin{center}Anmærkning og Tillæg.\end{center}

1.~\emph{Spørsmaalet} (Grundspørsmaalene nemlig, hvorfra den hele Stræben efter at
erkjende gaaer ud) opstaaer af \emph{sig selv}, fordi det højere Liv, der skal
frem, nu ogsaa vil frem, og derfor i sin Tids Fylde, ogsaa virkelig forkynder sig
og stræber frem i det Individ, der jo overhovedet kun er til, for at vorde dette
højere Livs Organ eller Kar. Derfor komme ogsaa \emph{a priori}, i det mindste i
dunkle Yttringer, de Grundbegreber frem, der ligge til Grund for disse Spørsmaal
og gjøre dem muelige. Men til \emph{Svarene} komme vi, om end ogsaa disse tilsidst
dog maae findes, tilmed gives, kun formedelst en Stræben og Gransken, der er vor
egen Stræben og Gransken. Hvorfor, er klart nok af det foregaaende (§.~4. Anm.~3.
og flere Steder). Drivten viser sig, som al Drivt, baade som det \emph{Universelles
Drivt} (Driven og Fremskynden) i os, saa og som vor egen Drivt, vor egen Attraa og
Begjær. Tendentsen er ogsaa her, paa eengang og i eet, en Tendents til at eniges
fuldkomnere med det almene og algyldige, altreelle, saa og en Tendents til
fuldkomnere Selvbestaaen; og naar vi ej lade os nøje med det første givne, saa er
det en dobbelt Tilskyndelse, som driver til at søge mere. Til Harmonie kommer den
dobbelte Stræben, naar den ene netop vil frem, for at den anden aldeles kan finde
sin Fyldest, naar vi kun stræbe efter Selvstændigheden, for at vorde Ideens og
Erkjendelsens fuldkomne, ved at vorde dens selvstændige, Organer.

2.~Det samme, som gjør Tvivl muelige, fører ogsaa \emph{Mueligheden af Vildfarelse}
med sig. Fordi vi under egen Søgen og Stræben skulle med Selvbevidsthed selv
bringe Erkjendelsen i Stand, derfor kunne vi gjøre andre Bestemmelser og danne os
andre Forestillinger, end de sande, ledede dertil ved at bestemmes af noget andet
end det, der skulde herske, eller gaaende videre, end dettes Tilsagn gaae.

\begin{center}\large\S.~21.\end{center}

Under Granskningen stræber Individet efter en Erkjenden, til hvilken det ej kan
komme uden selv at tragte derefter og anstrenge sig derfor. Men alt kommer nu an
paa, om ogsaa hos Individet det paafølger, der skal paafølge: for det første, om
Individet \emph{virkelig kommer til at erkjende} d.\ e.\ ej blot at erfare og
erholde en Underretning og Kundskab, men og at kjende det opfattede for gyldigt
og at optage og beholde det med Forvisning om dets Sandhed; dernæst om denne
Erkjendelse bliver saa levende og indtrængende, at en frugtbringende Selvtænkning
vil gaae frem deraf.

Forsaavidt er Ideens Kommen og Sigforkynden en \emph{Forsøgen}. Ja der er nu ogsaa
noget \emph{kritiskt} deri; det kommer an paa, hvad Vending Virkningen tager, idet
nu den frie Selvtænkning og Prøvelse, overhovedet den egne frie Gjøren træder den
imøde. Det er dette \emph{periculum vitæ}, som Franz Baader nævner (Sätze aus der
Bildungs- oder Begründungslehre des Lebens, Anmærkning til 6). Fordi Ideen skal
lyse ved sin egen Sandhed, i Intelligents og Frihed, derfor kunne dens reelle
Virkninger hos det Individ, hos hvilket den vil gaae op, kun have denne Charakteer,
men derfor kommer det da ogsaa an paa, om nu hos Individet Erkjendelse og
Selvtænkning komme eller ikke.

Men om vi end vel kunne tragte efter at erkjende, søge og stræbe, saa kunne vi dog
lige saa lidet give os Overbeviisningen selv, som vi selv kunne bringe eller give
os Sandheden: endog den fulde, for andre Tænkende aldeles indlysende og visse,
Sandhed kan staae lige overfor et Individ, uden at det endnu seer Sandheden deri
og bliver grebet deraf, altsaa uden at det endnu formaaer at erkjende, hvad det
dog allerede kjender. Just fordi Erkjendelsen skal være Sandhedens Indlysen ved
dens egen Sandhed allene, altsaa ved dens egen Kraft, formaaer ingen at give sig
eller andre den. Kun arbejde paa at \emph{vække} eller bringe til at yttre sig,
hvad der, constituerende Fornuftvæsenets hele Tilvær, er i ethvert, kunne vi; men
nøde det til at komme kunne vi ej.

Saaledes beroer da alt tilsidst paa noget, der kun kan tages, som ligefrem
\emph{Gave og Naade fra Herren}. Den, der, efter Paulus (Phil. 2, 13), udretter i
os baade at ville og at udrette, er og allene den, som virker i os at erkjende.

Er her da paa den ene Side noget, som staaer til os eller er sat i \emph{vor Magt}
(ἐστιν ἐφ' ἡμῖν, Epict.), med Hensyn til hvilket vi maae vise \emph{moralsk Kraft}
i Flid og Udholdenhed, redelig Sandhedskjerlighed og omhyggelig Naarvaagenhed, saa
er her paa den anden Side noget, som \emph{ikke staaer i vor Magt} (τὰ οὐκ ἐφ'
ἡμῖν),%
% FLAG p.167: Greek from Epictetus/Stoic (ἐστιν ἐφ' ἡμῖν / τὰ οὐκ ἐφ' ἡμῖν); verify accents/breathings against scan.
 hvilket vi maae møde med et \emph{religiøst Sindelag}, en religiøs Hengivenhed i
Guds Villie og hans Beskikkelse over os, et Sindelag, som baade skal vise sig i
rolig Oppebien, Taalmod, Ydmyghed og Resignation, saa ogsaa i Tillid og Mod.

\begin{center}Anmærkning og Tillæg.\end{center}

1.~Ved os selv skulle vi søge at erkjende, hvad der er at erkjende. Det nok saa
rigtigt og bestemt overleverede bliver ej vort, kommer ej til ægte Gjenfødelse, ny
Bevægelse og Liv i os, at det, fordi det i os hersker i sin Sandhed, nu og kan
tale og virke igjen med sin Kraft udaf os, og vi blive dets sande Organer, med
mindre det saaledes griber os og indsees og erkjendes af os, at en levende
Overbeviisning derom hos os kommer frem. Det maa efterhaanden blive os saa
ejendommeligt, som om vi vare komne til det ved os selv allene og havde øst det af
os selv, det er: umiddelbart havde faaet det af den i og for os talende Sandhed
selv. Heraf bliver en Følge, at vi ligesom oversætte det for os i vort eget Sprog,
udtalende det for os paa en ejendommelig Maade.

Hvorledes vi herunder unddrage os for enhver anden Autoritet for allene og
umiddelbart at underordne os den højeste Autoritet, Sandhedens egen, have vi
allerede forhen (§.~6. Anm.~1.) havt Anledning til at bemærke. Har Antagelsen paa
den ene Side ingen sand Værd, naar den ej gaaer frem af egen Indsigt og fri
Selvtænkning, saa har den dog ogsaa paa den anden Side ingen Værd, naar det ej er
Sandheden selv, som deri erkjendes.

Paa den ene Side maa Selvtænkningen nødvendigviis gaae sin egen frie Gang, og man
skal ikke ville ankomme ved noget andet Maal, end det, hvortil klar Overvejelse og
egen Indsigt af sig selv ville føre. Og dog er paa den anden Side Maalet, hvortil
vi skulle komme, nødvendigviis foreskrevet; thi naar det ej er Sandheden, hvortil
vi komme, saa er hele Bestræbelsen intet; men Sandheden er kun en eneste. Det
forholder sig her ligesom ved Dyden, som viser sig i, at vi med villig Hu og
levende Tilbøjelighed (Kjerlighed) deels bestemme os til, deels finde os i, det
nødvendige. Vi skulle arbejde fremad, som om det var os ligegyldigt, til hvad Maal
vi kom, naar det kun var et Maal, hvortil klar Indsigt og virkelig Overbeviisning
havde ført os; men dog kan det ej være os ligegyldigt, om nu og vor Overbeviisning
og formeentlige klare Indseen stemmer sammen med det i sig selv Sande, om det er
af Sandheden selv, vor Forvisning udstammer. Vel maae vi tilsidst berolige os ved
at sige, at vi redeligen have stræbt og gransket og nu ikke kunne see eller
erkjende andet, end hvad vi antage; men ingen vilde, i videnskabelig Henseende,
kunne berolige sig derved, dersom dertil ej den Tanke knyttede sig: og det, jeg
saaledes ikke kan andet end see og erkjende, er sikkert ogsaa i sig selv rigtigt. I
religiøs Henseende gives der vistnok endnu anden Beroligelse.

De Modsætninger, der her skulle komme til Harmonie, idet de tvende Modsatte ganske
af sig selv skulle træffe sammen i eet, disse Modsætninger sees især i det
theologiske Studium. Ordets Prædikant og Lærer skal tale Ordet til Hjerterne og
deels fremkalde deels nære levende Tro og Overbeviisning. Hos ham er det derfor, da
ingen kan tale til Hjertet, uden at tale af Hjertet, en væsentlig Fornødenhed at
være personligen levende overbeviist om, og gjennemtrængt af, Sandheden af det
Ord, han har at forkynde, den Lære, han har at lære. Han kan ingenlunde, som
Juristen, lade sig nøje med at paaberaabe sig, at saaledes er det ham foreskrevet,
og nu, skjelnende imellem det foreskrevne og den private Mening eller Dom, holde
denne sidste tilbage. Er ikke Christendommens, ja Kirkens Lære, i det mindste i
alle vigtige og væsentlige Puncter, tillige hans egen Overbeviisning, saa kan han
ej uden Uredelighed paatage sig den Gejstliges Kald; han vilde ogsaa ved at
paatage sig det, hilde sig ind i en for ham selv farlig Existents, idet han vilde
komme til at bære eet paa Læberne og et andet i Hjertet, thi i sit Kalds Udøvelse
er han nødt til bestandigen at træde frem som personlig overbeviist. Ja saa megen
Frihed, der og kan herske i Foredraget, saa meget end, efter Gavernes
Forskjellighed, Lære- og Virkemaade kunne være af forskjellig Art, en Frihed og
Ejendommelighed, just Christendommen i saa høj Grad ej blot tilstæder, men og selv
fremkalder, saa udfordres dog til at være en sand christelig Præst, at ikke allene
den hellige Skrifts Lærdomme i eet og alt udgjøre hans Overbeviisning, men at han
ogsaa saaledes gaaer ind i dens hele Udtryksmaade og ejendommelige Sprog, at ogsaa
dette bliver ham ejendommeligt, at just Bibelens Udtryk og Ord ere ham de meest
træffende og betegnende, og at han føler, at just et saadant Guds Ord fuldkomnest
forkyndes i saadant et Menneskeord. Han vilde ellers ikke være Christi Ords
Forkynder. Saa aldeles positivt og foreskrevet er det Maal hvortil den theologiske
Studerende skal komme.

Og dog skal han ingenlunde saaledes lægge an paa at komme dertil, at han med
Undertrykkelse af Selvtænkningen skulde beqvemme eller accommodere sin hele
Forestillen efter den overleverede hellige Lære; hvorved ogsaa netop det, hvorpaa
alt kommer an, den levende Overbeviisning, vilde hindres fra at komme kraftigen
frem og grundfæste sig.

Det har kun kunnet være den sælsomste Tankeforvirring, som nogensinde har kunnet
bringe til at statuere, at Troens Lærdomme skulle antages og troes uden alt Hensyn
til den egne Sandhedsfølelse og til den, vistnok kun under Forudsætning af egen
Overvejelse og tilmed egen Prøvelse muelige, indre personlige Samstemming dermed;
en saa slavisk Fordring udelukkede endog Mueligheden af den af Overvejelse og
Prøvelse udspringende levende Overbeviisning. Den studerende Theolog maa som
Studerende netop sætte den frie Selvtænkning i Bevægelse og intet antage enten
tvertimod eller med Fordølgelse af egen Sandhedsfølelse%
% FLAG p.171: "Fordølgelse" (OCR "Fordsigelse"); reading uncertain -- verify against scan.
 og Overbeviisning, intet uden egen Indseen.

Men imedens han skal følge den egne Indseen og Selvtænkning og den virkelige egne
Overbeviisning, skal dog hans Overbeviisning vorde en vis bestemt, idet den paa den
nys beskrevne Maade skal stemme overeens med den overleverede hellige Lære, og
hvad han ej skal lægge saaledes an paa, at han skulde lempe sig derefter og vænne
sig dertil ved at fæste Øjet derpaa allene, dette skal af sig selv blive hans
Stræbens Resultat. Kun ved at gjennemgaae en saadan Ildprøve viser det sig, om
hans Stræben ogsaa gik frem af et sandt indvortes Kald; kun af en saadan Ildprøve
gaaer en ægte Ordets Forkynder frem. Den Tro, han som Christen allerede maaskee
kan have vundet, skal han endnu engang som granskende Selvtænker vinde, om han
skal finde sig kaldet til at være Ordets ægte, det er selvstændige, frit talende og
virkende Organ.

En Modsætning, ja Modstrid eller Conflict af denne Art vilde vist ikke kunne
løses, og saa utallige Gange være bleven løst, dersom ikke Christi Ord og Lære og
Apostlernes Prædiken var af ganske andet Udspring, eller Herkomst, end
nogensomhelst Philosophie eller anden Lære, dersom der ikke ved dens
Guddommelighed og ægte fuldkomne Sandhed og al menneskelig Dybsindighed
tilfredsstillende Dybde ligesom var forud sørget for, at jo enhver til Grunden
gaaende Tænker, som vil grundigen indlade sig med den, nok maa komme til at finde
den.

Saa nøje og skarp bestemt træder ikke i nogen anden Erkjendelsens Sphære
Fordringen paa det, der skal blive Resultatet, op ved Siden af Opfordringen til
fri Selvtænkning. Men i sig selv er dog i enhver videnskabelig Stræben en saadan
Prøve at gjennemgaae. Thi i enhver Videnskab kommer det an paa, om vi komme til
Sandheden, og om den griber os saaledes, at den Selvtænkning og Selvvirksomhed,
der til Kaldets Udøvelse er fornøden, gaaer op i os. Vel er i de andre
Erkjendelsens Sphærer Sandhedens aldeles og for alle gyldige Hovedsum ikke
saaledes forud udtalt og nedlagt i en bestemt ydre Fremstilling, der kunde tjene
som ydre Prøvesteen, som i Theologien. Men ligesom der i alle Erkjendelsens Kredse
antages at gives visse Erkjendelser, der kunne være saa stadfæstede og afgjorte,
at den, der ej kunde naae Overbeviisning om dem eller Indsigt i dem, maatte ansees
udelukket fra denne Videnskab, saa gives der ogsaa i Særdeleshed, hvad ethvert
Kalds Udøvelse angaaer, visse Betingelser i Henseende til personlig Overbeviisning,
Indsigt og Aandsgave, eller Anlæg og Talent, hvilke ere fornødne for baade fra det
Offentliges Side og ved den egne Samvittigheds Dom at kunne tilstædes dertil.

Heri er da noget, som ikke videre beroer paa os, men hvori Livet under sin egen
Fremgang allene kan lære, hvad Livet vil give. Men saa er da ogsaa heri noget,
hvori det religiøse Sind, der føler sig i, og giver sig hen til, en højere
Styrelses Haand, især maa vise sig. Ingen kan forud vide om Tingenes Væsen, det,
hvorpaa alt kommer an, saaledes vil oplade sig for ham, at han vil naae den
fuldkomne, sig selv visse, i sig selv med Tilforladelighed hvilende Erkjendelse,
med det sikkre Omdømme og den fornødne Færdighed, saa at han kan sige sig at have
fundet, hvad der var Formaalet for hans Stræben. Aanden giver sig, til hvem den
vil, og i det Maal, hvori den vil. Og vi? ville vi vinde Livet, saa maae vi vove
Livet. Et Pant paa, at vi skulle naae, hvad vi attraae, gives os ikke forud. Om
den Studerende gjelde aldeles Schillers Ord: „Du mußt hoffen, Du mußt wagen, Denn
die Götter leihn kein Pfand, Nur ein Wunder kann Dich tragen In das schöne
Wunderland.“ Den Studerende maa, som Fichte i sine Forelæsninger „über das Wesen
des Gelehrten“ lægger os paa Hjertet, med Flid og Sandhedskjerlighed, overhovedet
med den Retskaffenhed, som han ogsaa i sit Studium skal vise, og med det
tillidsfulde Mod, som ogsaa hører herhen, arbejde og stræbe, som om han havde
Anlægget og nok engang vilde naae Maalet, og saa stille det i Guds Haand, hvad
enten han skal naae det, eller engang skal erfare, at han ikke var bestemt for
denne, men for en anden Plads, ikke for en højere, men for en lavere. Ingen kan
lægge en Alen til sin Væxt, enddog han bekymrer sig derfor, som Skrivten siger. Men
vi maae dog række os efter det, som er foran, og jage efter Maalet (Phil. 3, 14).
Og naae vi det ikke, saa kan den Tid dog ej være spildt, i hvilken vi lærte at
kjende os selv, lærte at finde vor Plads og at kjende Guds Beskikkelse over os.

2.~Men endnu i ganske andre Henseender er her en Prøve at gjennemgaae, ofte en
\emph{vovelig Prøve}, og det viser sig, at Studiet ogsaa kan have sin Mislighed, og
sin Fare. Forsaavidt nemlig deels Egoismen kan „tage Anledning“ af den højere
Erkjendelse, til hvilken Vejen skal aabne sig for os, og faae en desto
fordærveligere Indflydelse, jo højere Formaalene ere, deels selve Bestræbelsen
efter at komme til Indsigt og Klarhed just kan føre ind i Forvirring og
Vildfarelse.

Vi opfordres til at prøve vore egne Kræfter, selv at bestemme og afgjøre, og at
bringe det dertil, at vi tilsidst i vort Fag intet antage eller erkjende, uden at
det i vor egen frie Overbeviisning og Indsigt har fundet Samstemming og ligesom
faaet Sanction. Hvor naturligt da, at vi, naar Tiden til at stige op paa dette
højere Trin er kommet, rivende os løs fra andres Vejledning, begynde selv at danne
os vore egne Forestillingsmaader, og naar vi selv have dannet os saadanne,
fastholde og pleje dem med en særdeles Forkjerlighed, fordi vi i dem have prøvet
vor egen Selvtænkning. Men heri er da noget, som let kan opholde eller hilde vor
Fremadstræben til Sandhed, og lede paa Afveje, ja aabne det, der aldrig bør betyde
noget ved Siden af Sandheden, Indpas og Indflydelse. At dette eller hiint er
udfundet eller udtænkt af os selv, kommer nu let til umærkeligen at bidrage med
til at give det Betydning for os. Større bliver Faren, naar vi nu endog lade os
bringe til, før \emph{Modenhedens Tid, at træde offentligen frem med vore
Betragtningsmaader}, noget, der allerede, fordi det begrændser Granskningen til et
vist Enkelt og standser det alsidige, hensynsfrie, blot efter at naae Erkjendelse
stræbende Studium, ikke er godt i de egentlige akademiske Aar, og, naar det kommer
for tidligt, bærer Vidne om en ikke ret dyb Ærefrygt for Videnskaben. Det eengang
paastaaede fastholdes let ivrigen, fordi det eengang er paastaaet, især naar det
er skeet paa en offentlig Maade. Især hvis Polemik, der overhovedet saa let, for
at bruge et Ord af Göthe, trænger rolig Uddannelse tilbage, naar den ej gaaer frem
af levende Varme for en vigtig Sag og er forbundet med høj Aandskraft, bringer til
at tage det Partie, man tager, som \emph{Partie}, saa at en \emph{Partiaand}
sætter sig fast, og Tingen efterhaanden ej mere forsvares eller angribes, fordi en
endnu livfuld Overbeviisning og Indseen fører dertil og saavidt den fører, men
fordi det nu eengang hører vort Partie til at dømme saaledes. Ja langt mere
hæslige, endog lave og slette Motiver have vi ej sjeldent seet, at faae Overhaand
i den videnskabelige Stræben. Hvor megen moralsk Slethed der ogsaa hos Skribenten,
som Skribent, kan lægge sig for Dagen, derom afgiver vore Tiders Literatur
Vidnesbyrd, som vække Forfærdelse. Nu da Afvigelse har vist sig som muelig, maa der
derfor udtrykkeligen vaages over, at reen Sandhedskjerlighed allene raader og
hersker, og at kun Sandhed eftertragtes og intet andet, enten istedetfor den, uden
hvis Opnaaelse den hele Stræben jo vilde være intet, eller og ved Siden af den,
hvorved Bestræbelsen og Studiet vilde forurenes og forvanskes. Fordærveligt er det
især, naar Mangel af sand Ærefrygt for Sandhed endog gaaer saavidt, at den blotte
\emph{Vilkaarlighed} (det ganske momentane og tilfældige) kommer til at spille
istedetfor trolig og samvittighedsfuld Agtsomhed paa Sandheden selv, saa at snart
eet snart et andet synes og antages efter øjenblikkeligt Godtbefindende, og blotte
Indfald bestemme Antagelse og Dom. Være det nu, at en Higen efter det nye, det
usædvanlige, det glimrende og stærkt i Øine faldende eller anden forfængelig Lyst
fører til et, maaskee endog tøjlesløst og frivolt \emph{Spil med Erkjendelsens
Gjenstande}, som om al Agtelse for Sandhedens Højhed, al Følelse for dens Hellighed
og Omsorg for at lade den selv allene tale, var forsvunden, og kun frappante
Indfald og Udtryk vare det, hvorpaa det kom an. Eller ogsaa, at mat Ligegyldighed,
en sig selv bevidst Mangel af sand Overbeviisning om nogetsomhelst, saaledes
opkommer, at alt kun antages og beholdes efter øjenblikkeligt Tilsyneladende som
en blot Mening, der kan slippes igjen, og hvori man ikke føler sig at eje noget,
man veed at kunne, eller troer at kunne, forlade sig paa med inderlig Tillid og
Tilforladelighed: en \emph{Indifferentismus}, der bærer den rene Intetheds Præg.
Noget maa der være, hvorom Mennesket er saa overbeviist, som om sin egen Tilvær, at
han vil \emph{leve og døe paa dets Sandhed}. Kun af en saadan Overbeviisning kan et
kraftigt og dygtigt Liv gaae frem. [Hiin Indifferentisme maa iøvrigt ikke
forverles med det, man ved Ordet \emph{Tolerance} har villet betegne.
\emph{Intolerance}, der gjærne vilde paatvinge andre sin Erkjendelse eller
overhovedet for at gjøre den gjeldende bruge andre Vaaben og Midler, end Sandhedens
egen Magt og Kraft og det Sandheden forkyndende, til Overbeviisning henrivende og
aandeligen nødende Ord, hvori Sandheden, vidnende om sig selv, ogsaa virker ved
sig selv, bærer lige saa fuldt som Indifferentisme, skjønt paa en ganske anden
Maade, Præget af Mangel paa ægte Følelse og Agtelse for Sandhedens indre Højhed og
Hellighed; thi Intolerancen gaaer ikke ud paa at fremme Sandheden som Sandhed, men
blot dens ydre Antagelse og Bekjendelse uden ægte Overbeviisning, som om Sandhedens
saakaldte Antagelse uden denne levende Overbeviisning virkelig var en Sandhedens
Antagelse, en Udbredelse af Sandhedens Herredømme. At sætte andre Midler i
Bevægelse for at fremme en Erkjendelse, end netop dem, der kunne virke sand levende
Overbeviisning, er ikke blot at vise en synderlig Mangel af Tillid til Sandhedens
Magt, som om dens egen Magt i og for sig ikke var stærk nok; men fører endog til
at forfejle det, hvorpaa det egentligen kommer an, og at lade en anden Myndighed,
end Sandhedens, indtage den Plads, hvilken dens egen Magt skulde have. Ethvert
Foretagende af denne Art gjør (ligesom det Enigmord,%
% FLAG p.179: "Enigmord, Saul begik" (OCR); word uncertain (Saul's sacrilege, 1 Sam) -- verify against scan.
 Saul begik) et Rov paa det, der hører Sandheden til, hvoraf da indlyser, at
Indifferentismus og Intolerance begge fattes den ægte dybe Kjerlighed til, og
Agtelse for, Sandhed, der kraftigen vil Sandheden, og derfor kun vil den rene
Sandhed selv i dens Reenhed. Tolerance er tilmed (for dog ogsaa at bemærke dette)
ikke noget, som ligger midt imellem begge, men noget, som i Henseende til Motiv og
Princip er ganske forskjelligt derfra].

3.~Dog ogsaa uden at nægte og lave Bevæggrunde spille ind med, eller en
forkastelig Tankemaade ligger til Grund, kan Eftergranskning føre til Vildfarelse;
endog hos den ivrigt stræbende. Den i Bevægelse værende Stræben driver fremad, men
Antagelse og Forestilling bliver endnu ikke, saalænge det Sande, som skulde
bestemme og beherske den, endnu ej har viist sig eller ej har viist sig
klarligen, bunden og bestemt til en vis eneste Maade at forestille paa, men finder
flere Mueligheder for sig, imellem hvilke der nu vakles. Derved standses
Bestræbelsen efter at komme til Erkjendelse i sin Fremgang. Da er det, at
letteligen, uagtet al redelig Sandhedskjerlighed, blot Drivt til at komme fremad
bringer til at gribe en Udvej og at forfølge den, uden at bestemt Vished haves for
dens Rigtighed; hvoraf da Følgen let bliver, at en urigtig Forestillingsmaade
sætter sig fast, der, om den end og i Begyndelsen erkjendtes for blot
\emph{hypothetisk}, dog ej sjeldent siden faaer Sandhedens og Vishedens
Herredømme, idet den, naar meget derved sees bragt til Eenhed eller saaledes
drages hen dertil, at det bringes i Forbindelse og Sammenhæng dermed, (noget, som
med den falske Forestilling er ganske mueligt), tilsidst bliver tilvant og som
hjemmehørende i vor hele Forestillingskreds. Ja for at komme videre og ej at
standses ganske, kan det stundom være fornødent eller dog naturligt og
befordrende, forsøgsviis at danne saadanne hypothetiske Forestillingsmaader og
forfølge dem, indtil de under Undersøgelsens Fremgang kunne prøves og vidne om sig
selv. Dog maa det ikke glemmes, at vi, som ej kunne beherske den højere Virken,
hvorved Erkjendelsen gives os, ikke bør haste fremad paa egen Haand ud over det,
der gives os, men maae oppebie, indtil det Sande selv vil vise sig for os. Men
heraf bliver da let Følgen, hvad man siger: \emph{at Vejen til Sandhed gaaer
igjennem Vildfarelse}. Den systematiske Tendents, der stræber efter at bringe det
mangfoldige, stykkeviis erkjendte og begrebne, til Eenhed og deri vil skue et
Hele, og til denne Ende snart ønsker en Mængde Puncter bestemte, fører især let
dertil. See §.~15. Anm.~2. og følgende. Overhovedet gjelder om den enkelte
Gransker, hvad vi see at gjelde om Menneskeheden i dens Fremgang til Erkjendelse,
og hvad der synes at være en Grundlov for alt Livs organiske Udviklingsgang i
Almindelighed: at den organiske Stræben, igjennem hvilken Tingenes indre Væsen skal
komme til ydre Fremtræden, Aanden og Sandheden til Herredømme, maa gjennemgaae en
Fleerhed, ja maaskee en talløs Fleerhed, af Standpuncter og Tilstande, paa hvilke
og i hvilke Livet efterhaanden maa sætte sig fast en Tid og komme til en vis Grad
af Klarhed, Eenhed og fast Bestaaen og til bestemte Resultater, inden det kan naae
til det Punct, da det fuldkomment kan modnes til at vige derfra igjen og stræbe op
til højere og derfor mere overskuende Standpuncter og nye Organisationer. Men saa
tilsigtes da den \emph{fuldkomme Alsidighed}, hvori alt kommer til i alle
Henseender at vise sig retteligen i sit Lys og med sin relative Gyldighed, igjennem
en \emph{Række af Eensidigheder}. Deels tage vi efterhaanden det, der skal
erkjendes i sin Sammenhæng og Heelhed, fra forskjellige Sider, (vi have omtalt
dette forhen §.~15. Anm.~4.), deels trænge vi under Betragtningen og Undersøgelsen
af det sig frembydende Mangfoldige altid først ind i en \emph{vis bestemt
Mangfoldighed}, søgende dens Eenhed og indre Væsen, inden vi søge at omfatte den
tilligemed en vis anden, der maaskee ogsaa for sig allerede er undersøgt, og nu at
finde det højere Eenhedspunct for begge, hvorved da ogsaa det deri erkjendte indre
Væsen, det er det deri fremtraadte Udtryk af Ideen, vil komme til at vise sig i et
nyt Lys, naar det nu skues som det lige og det primitive og constitutive i begge
disse forskjellige Indbegreb af Mangfoldighed. Idet vi fra Mangfoldighedens Side
nærme os Eenheden og Ideen, saa og idet vi søge at gjenskue en erkjendt Idee i det
Mangfoldige eller at gjøre den gjeldende deri, og nu derfor gjøre dette Mangfoldige
til Gjenstand for en særegen Randsagelse (see §.~16.), maae vi begynde med at tage
visse enkelte af det større hele omfattede mindre Heelheder for os, hver for sig,
og først trænge ind til det relativ første og til Eenheden i enhver af dem, inden
vi kunne komme til at faae den dybere Eenhed i dem alle; thi just det inderste i
enhver af dem er det Punct, hvori de væsentligst ere sammenknyttede. Men naar
Tiden kommer dertil, saa vil nu og det, hvad der i enhver af disse Sphærer syntes
at staae fast, ej mere kunne det; det højere Standpunct, hvortil vi ankomme, vil
vise alt i et nyt Lys, og det til en vis Grad af fast Bestaaen komne maa da løsne
sig igjen for paa en levende og organisk Maade at gaae ind i den større Heelhed.

Men saaledes vil der da letteligen, under Studiets Fremgang, paa en fast, rolig og
sig selv tilfredsstillende Erkjendelsestilstand kunne følge en Tilstand, hvori
Usikkerhed, Uvished, Uro indtræder. Hvad vi havde, kunne vi ej beholde, som vi
havde det, og det nye, der skal komme, har ikke vundet Skikkelse i os endnu, men
bevæger sig i os i en gjærende Masse, hvori tusinde enkelte Elementer til en ny
Dannelse bevæge sig ved Siden af hinanden og imod hinanden, hver for sig gjørende
sin Betydning gjeldende, uden at de endnu have fundet deres relative Plads. I
Videnskabernes Udviklingsgang i Menneskehedens Historie finde vi, ligesom
Geognosten finder det i Jordklodens, ikke blot forskjellige Formationer, som have
afløst hinanden, hver for sig uddannet paa sin ejendommelig Maade og med sine egne
Retninger, men ogsaa finde vi midt imellem tvende saadanne ej sjeldent
Overgangstider, hvori Brudstykker af en tidligere Formation findes ved et aldeles
heterogent Bindemiddel forbundne og ligesom sammenæltede til en sammenhængende, men
ingenlunde heelt igjennem eensartet, end sige chrystalliniskt uddannet Masse.
Ogsaa den videnskabelige Udviklingsgang er en undulatorisk; Stigen og Synken
afvexle med hinanden. Hvad der imidlertid kan og skal, som det constante, gaae
igjennem denne Vexel af rolig Bestaaen og gjærende Vorden, og selv midt i denne
holde en højere Ro vedlige, at Uroen ej kan trænge ind til Livets og
Bestræbelsens Kjærne, er en fast Overbeviisning om Grunderkjendelsens, og tilmed
selve den derved fremkaldte Stræbens, Realitet, Overbeviisningen om, at der er
noget at søge og derfor ogsaa at finde, en Overbeviisning, der forbunden med det
religiøse Sindelag vil gjøre, at der under al til sine Tider kommende Vaklen,
Utilfredshed og Formørkning, vil blive en urokkelig Grund, og med den en tryg
Grundfred og uformørkelig Grundklarhed i vort Allerinderste, en Stemning, der ej
lidet kan vedligeholdes ved hiint foromtalte Henblik til et Ideal (§.~13. Anm.~3.),
hvilket Schiller i sit Digt: „die Ideale und das Leben“ saa herligt har skildret.

Hvad i Særdeleshed den til ægte Videnskabelighed fornødne philosophiske Grunden
angaaer, saa har enhver, længe førend han kan være moden til philosophiskt
Studium, allerede sat sig fast i en Maade at tage og skue Livet paa, hvilken
Philosophien ej kan lade blive staaende, han har bortgivet Realitetens Navn til
det, der ingen sand Realitet har. En anden Forestillingsmaade end den
philosophiske holder sig ogsaa hele Livet igjennem fast ved Siden af denne,
omtrent som vi blive ved at tale, som om Solen gik omkring Jorden, saa vel vi end
vide, at dette ej er saa, og saa let den øvede kan i ethvert Øjeblik sætte sig over
paa det rette Standpunct. Allerede derfor er det intet ringe, der skal skee, naar
Philosophien skal komme til et Individ og blive levende i dette; thi dets hele
Forestillingskreds skal løsne sig og ligesom smelte hen, for at undergaae en ny
Dannelsesproces og at gjenfødes. Dog har det, naar Tiden dertil kommer, ikke saa
stor Vanskelighed i denne Henseende hos den lærelystne Studerende, thi netop en
Længsel efter Ideen vaagner i den Tid, da han er moden til det akademiske Studium,
lige med hans Stræben efter Selvstændighed, og han er let beredt til at aabne sig
for det, der først skal bringe det rette Lys og det rette Liv ind i alt, hvad han
hidtil lærte. En Tid lang giver han sig villig hen til Ideens Betragtning. Under
den Form, den bydes ham, tager han aabent imod den, og den højere Selvtænkning
gaaer i eet med en højere Læren. Det ham saaledes paa en livfuld Maade, igjennem
en fri Selvtænkning til hans egen Selvtænkning, overleverede tænker han sig let
ind i, tilegner sig det, og glæder sig i det og ved det; fører det ogsaa i mangt
et enkelt Punct videre paa en selvstændig Maade, og knytter dertil mangen egen
Besvarelse af selvopkastede Spørsmaal; det synes nu ganske at være assimileret.
Men lever en dybere Drivt til Philosophie og ægte Videnskab i ham, saa vil det med
Tiden vise sig, at Tilegnelsen dermed ej var til Ende. Det saaledes optagne vil
efterhaanden komme op igjen, gjennemarbejdes endnu engang for nu først under en
egentlig og fuldkommen Gjenfødelse at assimileres tilfulde og, (for at bruge endnu
et physiologiskt Ord) optages i, eller gaae over i, Constitution; alle eller dog en
Deel af de Hovedspørsmaal, der syntes vel besvarede og bragte til Ro, især det til
Grunden for dem alle liggende Grundspørsmaal, komme nu først saaledes op, at
Braaden deri eller Veerne derved (ὠδίς, for at tage et Udtryk af det andet af de
platoniske Breve) blive fornumne; og nu spørges der først egentligen ret af
Hjertens Grund. Hvad vi forhen fattede og overskuede i dets hele Sammenhæng og
derfor erkjendte i denne dets Sammenhæng og Conseqvents, og hvori vi saaledes
fandt det, vi søgte, Ideen og Grunderkjendelsen, at vi villig optoge det, dette see
vi nu i et andet Lys. Hvad den allerede modnede Tænker maaskee strax havde fundet,
naar en saadan Lære var bleven bragt ham, at den, sin maaskee gode, velgrundede og
meget omfattende Sammenhæng uagtet, dog enten aldeles ikke eller ej tilstrækkeligen
havde iagttaget dette eller hiint, der ellers var tænkt i denne Videnskab eller dog
kunde være bleven tænkt paa, dette vil nu og under sit Studiums Fremgang den
Studerende kunne bemærke om det, han tidligere beredvillig lod finde Indgang hos
sig. Fuldkomment er jo intet, der som videnskabeligt System overleveres, lige saa
lidet hvad et alsidigt Hensyn og Omblik angaaer, som ellers. Endog en vis
umiddelbar, vistnok ej altid ubedragelig, Sandhedsfølelse, som, uden at
ugjendrevne Indvendinger endnu have dannet sig, kan stræbe det modtagne eller dog
den Maade, hvorpaa det er udviklet, imod, kommer nu i Betragtning med. Hos den, hos
hvem en saadan Følelse eller Fornemmelse modstræber, er den noget reelt, der, lige
saa vel som bestemte Indvendinger, skal klares, og ikke tør tilsidesættes. Idet
man skal føres ind i et ellers dybtgaaende og vel sammenhængende Tankehele, maa
man vistnok ikke lade deslige modstræbende Fornemmelser, ja ej engang mange
bestemte, fra den hidtil havte Forestillingsmaade hentede eller øjenblikkeligen
opstaaende, Indvendinger faae en betydelig eller endog forstyrrende Indflydelse.
Der er mangen ny Lære, mangen ny Tanke eller ny Fremstillingsmaade, hvorom der med
Føje kan siges: Prøver den practiskt, optager og sætter Eder ind i den, og I
skulle vel forfare, at den er af Sandheden. Men har nogen just optaget den saaledes
og levet med den, og er han netop derved bleven modnet til at stille sig imod den
(som man da overhovedet kun ved at erkjende, bliver ret i Stand til paa den rette
Maade at bestride, det, der ellers har nogen sand Gehalt), saa kan han ogsaa give
en modstræbende Følelse Magt; maaskee deri noget anmelder sig, der igjennem ham
vil bryde en ny Bane, aabne Indgangen til en ny Erkjendelse.

Studiet fremkalder overhovedet en Vexelkamp imellem en i en bestemt Form udtalt og
sig i denne med Fasthed haandhævende Videnskab paa den ene, og en Individualitet,
der ogsaa gjør sin Ejendommelighed gjeldende, paa den anden Side. Kun paa vor
Maade kunne vi fatte, i vort Sprog skulle vi ogsaa oversætte for os, den
almeengyldige Sandhed, ja som den almeengyldige skal den vise sig ved ogsaa at
kunne gaae herskende ind i os, som vi nu eengang ere, naar vi kun ei fattes
fornødent Anlæg og stræbe ivrigen og redeligen. Men i en conseqvent udført, i sine
Dele nøjagtigen bestemt, Videnskab kommer den os dog imøde og synes ej at kunne
træde ud af denne eller give noget efter deraf. Imidlertid maa den dog løsne sig,
og fra den faste Form gaae over i en vis Fluiditet, for at kunne gaae ind i os; thi
den kommer jo kun til at blive vor derved, at den ved Kraften af dens indre første
bevægende Livsprincip gjenføder eller bygger sig op igjen i os; og saa kan da
ogsaa, under denne dens Gjenoplevelse eller Gjenopbyggelse i den nye
Individualitet, det nye Bidrag, hvad denne Individualitet er i Stand til at levere,
komme til, og nu altsaa ogsaa Individualiteten under den indtrædende Vexelkamp nyde
sin Fremme.

Men heraf opkommer da let en Gjæring, hvilken den kraftigt stræbende kan have at
gjennemgaae og det vel endog flere Gange i forskjellige Epocher. Det almene skal
blive hans, men for at blive hans, bygge sig op igjen af sine Elementer, ja
undergaae en vis ny Modification i Henseende til sin Udtryksmaade, maaskee endog en
Metamorphose og Reformation (hvilken man dog maa vogte sig for af Lyst til
Originalitet at lægge an paa), maa det i ham løse sig op i sine Elementer. En
Mængde Puncter lyse frem, en Mængde adskilte, ofte modsatte Antagelser paatrænge
sig, men Eenhed og Harmonie kan endnu ej bringes ind i den i Bevægelse komne, men
endnu ej til Klarhed komne Masse. Da fører vel ej ganske sjeldent Tvivl og Uro
saavidt, at der endog tvivles om Realiteten af den hele Bestræbelse, om Realiteten
af den til Grund for denne Bestræbelse liggende Grunderkjendelse, der nu ej synes
at bestaae sig og at holde ud under den hele Fremgang, der dog --- for at gjentage
et forhen forekommet Ord --- ved at lykkes skulde vidne om den. Om vort eget Kald
til at stræbe frem paa denne Vej opstaaer endnu lettere foruroligende Tvivl.
Herunder maae vi da lægge paa Hjertet, at, imedens vi endnu først stræbe fremad og
arbejde paa at finde Sandheden og bringe den til Klarhed og Indlysenhed hos os ved
at uddanne et Indbegreb af enkelte Erkjendelser og Fremstillinger, som skal blive
dens adæqvate Udtryk, maae vi, eftersom dette dens Udtryk, dens klare Fremtræden
dog kun kan fremkomme som en Virkning og Følge af dens egen Indflydelse og Magt,
nødvendigviis have givet os hen til den og have ladet os bestemme og lede af den,
førend Evidentsen og Forvisningen kan komme. Productet, den klare og bestemte
Fremstilling, tilmed og den Sandhedens Fylde eller fuldkomne Sigaabenbarelse, som
først dermed kommer, kan dog ikke indtræde strax tilligemed det, der sætter den
productive Stræben i Bevægelse.

Ogsaa i det indre Liv i Bevidstheden, hvori vi skulle voxe i Indsigt og
Erkjendelse, er der en af sig selv opererende Virken og Vorden (φύσις), hvortil vor
egen Arbejden refererer sig, og hvorved der nu af sig selv maa gaae et Lys op for
os. Men dette af sig selv arbejdende Lys i os kræver sin Tid og har sine Stadier,
ligesom sine Kriser.

Og her er det da, at, baade i religiøs og i videnskabelig Henseende, en Tro (see
§.~5. Anm.~2), som tillidsfuld Tilknytning til, og Overbeviisning om det, hvis
Kraft og Betydenhed allerede har vist sig og vidnet om sig selv, skjønt dets Fylde
ej er kommet endnu, maa finde Rum. Ja ligesom der nu kan behøves Taalmod og Tillid
til ej at overvældes af Uroen, saaledes er og en vis Respect for et αὐτὸς ἔφα
fornøden.%
% FLAG p.191: Greek "αὐτὸς ἔφα" (ipse dixit) and "φύσις" p.190, "ὠδίς" p.186 -- verify accents against scan.
 Skulle vi end ikke antage paa Autoritet, saa kunne vi dog ikke undslaae os for at
have Agtelse for Autoritet, for grundige Tænkeres, ja overhovedet for tænkende
Menneskers, Dom, og ligesom det ej lidet styrker vor Overbeviisning, at vi i samme
mødes med en anden Tænker, eftersom det, der viser sig eens for Tvende, som ad
forskjellige Veje ere komne dertil, derved synes at vise sig som gyldigt over
begge, saa er det og naturligt, at vi, naar vi støde an imod grundige Tænkere,
derved bringes til at standse og til ny Overtænkning. Men især bør den, der endnu
ikke er kommet til den Gjennemgranskning af en stor Mangfoldighed og det Overblik
derover, der i enhver Videnskab er fornødent til en solid Dom om nogetsomhelst
enkelt Punct, af Agtelse for deres Ord, der ere komne saavidt, holde sine
modstridende Antagelser og Domme in suspenso. Som jo overhovedet den fremad
stræbende Tænker altid, indtil han er naaet til en vis Oversigt, og indtil en vis
Grad har forfulgt en Grundtanke i dens Følger, altid befinder sig i en vis Svæven
imellem Forvisning og Ikkeforvisning, medens ogsaa det, hvorom han begynder at
forvisses, endnu kun har en vis ubestemt og svævende Skikkelse.

4.~Men naar vi da nu saaledes see den Fare eller \emph{vovelige Prøve}, man i det
philosophiske Studium har at gjennemgaae, saa tør vi ej fordølge for os, at det
just ikke kan være raadeligt for enhver at føres ind i Philosophiens Dyb, og at vi
ingenlunde tør undre os over, om selv Mænd af fortrinlig Aandsgave ej sjeldent have
Utilbøjelighed til at indlade sig med den i den abstracte Form, der er den
Philosophien egentligen tilkommende Form. Et udmærket Talent til Konst eller til en
vis Art af practisk Virksomhed, ja til visse Arter af videnskabelige Granskninger,
kan, naar Reflection og indtrængende Grunden ikke naae Dybden af det, man i sit
Talent besidder, letteligen forvirres og forstyrres eller standses ved
philosophiske Speculationer. Philosophering i videnskabelig Form er ej heller
hverken den eneste eller den højeste Maade, paa hvilken vi kunne stræbe at tilegne
os og besidde Ideen og den aandelige Gehalt, Ideen giver Livet. Fromhed,
Kjerlighed og Dyd ere højere end Philosophie, og (Poesien kan stille sig ved Siden
af den). Eet er det, under en levende Receptivitet og Productivitet at besjæles af
Idee, et andet, at stræbe at udgrunde Beskaffenheden af det, der saaledes ligger
til Grund for den kraftige Virksomhed, til hvis Udøvelse vi føle Kald og Gave. Med
Talent til Virksomheden selv er just Talent til at udgrunde dens indre Natur og
Væsen ej altid, ja ei letteligen forbunden. Men saa kan den umiddelbare kraftige
Udøvelse let gaae tabt over Bestræbelsen efter, philosopherende at trænge ind i
det, der constituerer den. Bekjendte ere Bacos Ord, at Philosophien, naar man blot
kommer halvt ind i den, fører fra Gud, og kun, naar den grundig studeres, fører
tilbage igjen til Gud. Og saaledes vil man og ved et ikke noksom til Grunden
gaaende Studium af sin egen Levens ideelle Gehalt komme bort fra denne. Ideen, man
formedelst abstract Tænkning vil stræbe at skue klarere, som den i og for sig er,
kan ikke vedblive at leve i os, som den forhen levede i os, som den ikke særligen
til Syne kommende Sjæl i vor practiske Virksomhed, og det gjelder nu om den
Tænkning, der har gjort den til Gjenstand for indtrængende Reflectioner, kan
igjenvinde den under den philosophiske Form, og paa ny bringe den til den
umiddelbare kraftige Leven, eller om denne midt under hiin Tænkning kan bevare sig
ved Siden deraf.

\begin{center}\large\S.~22.\end{center}

Ikke blot leve i Individet skal og vil Ideen, men ogsaa udtale sig igjennem
Individets udadgaaende Fremstillen og Virken og derved komme til en fuldkommen
Fremtræden, en fuldkommen ydre Hersken og Tilværelse i Endeligheden. Og ligesom
Ideen kun kommer til at udgjøre, og at vise sig som, al Fyldes Indbegreb i
Individet, naar den behersker denne fra Individet udgaaende Virksomhed, saa kommer
ogsaa Individet kun til, fuldkomment at have Ideen, idet den søger at realisere
dens Fordringer i en saadan udadgaaende Virksomhed.

I blot videnskabelig Henseende gaaer denne Realisationstendents ud paa at
fremstille Ideen i et paa en ydre Maade udtalt, og derved til udvortes Realitet
kommende, videnskabeligt Hele, paa hvis Istandbringelse Individet altsaa arbejder.
Men desuden kan og maa endnu ogsaa Ideen objectivere sig gjennem Individet i en i
Menneskelivet mere umiddelbart indgribende Praxis eller practiske Udøvelse.

Denne Stræben knytter nu Individet nødvendigviis til et \emph{Menneskesamfund}, og
i et saadant Menneskesamfunds eller Menneskeheles fuldkomne Organisation og
Uddannelse til at vorde besjælet af Ideen, søger Ideen Fuldendelsen af sin Kommen
til ydre Bestaaen og Tilvær eller Existents. Tendentsen til Fremstilling bliver en
Tendents til \emph{Meddelelse} og en Stræben efter at finde Deeltagelse. Vi have
kun Ideen, forsaavidt den gjennem os udtaler sig i sin Kraft, men dens fulde Kraft
i vor Udtalelse fornemme vi ved intet saa meget, som derved, at den gjennem os
kommer til at gribe andre med samme Magt, som os, og nu at gjenføde sig i dem.
Allerede til fuldkomment at udtale sig hos den Enkelte kommer den først under hans
levende Tankemeddelelse og Vexelvirkning med andre.

Det højere Liv, der vil leve i den Enkelte og komme til Skikkelse i ham, vil
ligefuldt leve i dem alle, og kommer kun fuldkomment frem, kommer kun til den
Skikkelse (Gestaltung), hvortil det vil komme, naar det træder frem som det ene Liv
i det hele harmoniskt forbundne Samfund. Saa er der da intet, som mere er
Incitament for dets Rørelser hos den Enkelte, for dets geniale Fremkommen i den
Enkelte, end dennes Sammentræf eller Sammenstød med andre tænkende Individer.
Derfor er ogsaa Studiet af Videnskabens Historie saa vigtigt, da dette fører os ind
i Samqvem med alle Aarhundreders fortrinlige Granskere.

At bringe en Videnskab i Stand og til Fuldendelse bliver nu en fælleds menneskelig
Sag, et Menneskehedens Anliggende. Interessen bliver en Almeeninteresse. Men skal
den Enkelte fuldkomment kunne arbejde herfor, saa maa han begrændse sig, det er
begrændse sin Productivitet, til noget vist bestemt, paa hvis Udførelse han vil
arbejde, medens han dog, idet han derved som Led griber ind i det store Hele, maa
opfyldes af det Heles Aand, og gjøre Studiet af det valgte Enkelte til Middelpunct
for en heel individuel Uddannelse.

\begin{center}Anmærkning og Tillæg.\end{center}

1.~Allerede fra først af, ved Naturen, komme vi frem og constitueres som Led af et
Hele, og vise os derved som dem, der ere bestemte til at voxe ind i det sig
formedelst den menneskelige Stræben og Arbejden dannende Menneskehele. Paa den
meest afgjorte Maade viser Naturen selv hen paa et saadant ved bestemte Talenter,
hvorved der anvises Individers Bestræbelser bestemte Retninger og Formaal, hvortil
da ogsaa den med Talentet forbundne Drivt ansporer, men hvorunder naturligviis saa
mange andre vigtige Formaal, om ogsaa Sands og Interesse for dem ikke udelukkes,
dog hvad Productiviteten angaaer, blive uforsøgte. En høj Gave er et saadant
afgjort Talent at agte for, saavel i saa mangen anden Henseende, som og fordi det
hæver op over Nødvendigheden af et ved Reflection og Overvejelse ledet Valg, og
befrier for dettes mangelunde Misligheder; thi hvor Naturen saa bestemt har talt,
der er, som Steffens med Rette ofte erindrede, en højere Beskikkelse at respectere,
uden alt Hensyn til borgerlig Stilling, ja til de naturlige Menneskeforskjelle, og
altsaa er da intet Valg fornødent. Imidlertid maa dog her, ligesaavel som ellers,
Individet modtage sin særdeles, om end af Naturen allerede anviiste, Gjenstand af
Ideen selv, og derfor tage den saaledes, som denne anviser den; Granskeren maa,
som Fichte etsteds siger „lade sig sin særlige Videnskab give af Videnskaben;“
denne er det, han i hiin maa opsøge og søge at besidde. Den Almeenidee, der fordrer
Begrændsningen, maa ogsaa lede os til den og lede os ved den.

2.~Men ikke enhver lærer Naturen gjennem et særdeles bestemt fortrinligt Anlæg og
en bestemt Drivt, hvad han bør ansee for sit Livs Opgave. Og da kunne Forhold og
Omstændigheder, Livets hidtil tagne Retning og Gang, der for den, som eengang er
deri, er at tage for en Tilskikkelse eller en højere Bestemmelse over ham,
ligeledes ogsaa andres Oversigt og modnere Dom, lige saa snart blive os en
Lederraad, som egen Lyst eller Forkjerlighed. Med afgjort Ulyst til et Fag bør
vistnok ingen vælge det; thi at gaae ind i en Virkekreds med bestemt Mangel af det,
der skal drive Værket, er at gjøre et voveligt Forsøg med det, der ikke er vort
eget. Men Lyst er endnu ej strax Trang eller Drivt, og hvor Lysten ej har en
bestemt Drivts Fynd og fremadstræbende Magt, der er Lyst endnu ej strax at tage
for indvortes Kald. Især er Lystens indre første Udspring og Motiv at iagttage. Om
den Lyst, for hvilken Hensyn til Erhverv, forfængelig Ære eller deslige
uvedkommende Hensyn, der aldeles bør holdes borte, ligge til Grund, ville vi ikke
engang tale. Men vel bør der gjøres opmærksom paa, at Lyst til en vis Virkekreds,
Lyst til en vis Art af Liv og Existents, der indtager ved sin egen Skjønhed og ved
den velgjørende Indflydelse paa mange Mennesker, hvortil den efter sin Natur maa
føre, ikke forverles med Lyst til selve den Videnskab og Granskning, der skal
berede og gjøre duelig til en saadan Virkekreds. Man kan for Exempel udmale sig en
Landsbypræstes eller en practisk Læges Stilling og Virkekreds saaledes, at man
indtages af Lyst til den, for at kunne komme til at gribe saa velgjørende ind i
Menneskelivet og at glæde og lyksaliggjøre saa mange. Men det er en ganske anden
Lyst, hvorpaa det hos den Studerende, der skal afgjøre sit Valg, kommer an. Kun
naar han tilligemed hiin Lyst ogsaa føler Lyst til at randsage Skrivten og
efterspore ethvert Hjelpemiddel til ret fuldkomment at trænge ind i dens Aand og i
ethvert af dens Ord, saavelsom Lyst til at randsage Menneskehjertet og
Menneskesindet for at udgranske, hvad det i de forskjelligste Stillinger og Forhold
til sin Frelse og Fremme kan behøve; eller naar han (for ogsaa at tage det andet
Exempel) føler Lyst til at studere Legemets hele Bygning, Functionernes
mangfoldige Arter, Ejendommeligheder og deres Forhold, naar det interesserer ham at
efterspore Sygdommenes Natur og saa mange Lægemidlers Egenskaber og Virkninger, og
saa fremdeles: kun da bevæges han af den Lyst til selve Videnskaben og
Granskningen, der først lærer, at hiin Lyst til Virkesphæren var et sandt Kald til
den.

3.~Hensynet til en Praxis, hvortil man vil gjøre sig duelig, er vistnok allerede
under Studiet forsaavidt af megen Vigtighed, som det kan gjøre opmærksom paa alt,
hvad der engang i Tiden vil kunne udfordres dertil og nu allerede maa omfattes med
af Studiet; som da i alle Hovedfag egentligen Ideen om et bestemt Kald, en bestemt
i Livet indgribende Virken, leder til at samle en Mængde ellers adskilte
Granskninger omkring et Middelpunct. Men ikke maa dette Hensyn hilde saaledes, at
noget, som Videnskabens Idee fordrer med til sin totale Udførelse til et
systematisk Hele, skulde tilsidesættes, fordi det ej synes vigtigt for det senere
practiske Liv, eller at overhovedet det Spørsmaal, hvorvidt det, Videnskabens Gang
fører til, er fornødent for det practiske Liv, skulde komme frem med Hensyn til de
enkelte Puncter. Den bedste Beredelse til Praxis er et grundigt Studium af
Videnskaben selv saaledes, som den er; thi til en solid Praxis udfordres at være
aldeles hjemme i Videnskaben og at have tilegnet sig dens Aand, dens Principer og
Methoder, men dertil kommer kun den, som studerer Videnskaben i dens Heelhed,
tilmed dens Principer og Methoder i deres hele Udførelse og forfølger dem i alle
deres Hovedgrene eller Retninger. Men ligesom Hensynet til Praxis ej bør betinge
eller hilde, saa bør det ej heller bekymre eller ængste, saa at vi ved ethvert
enkelt Punct skulde sørge for, hvorledes vi skulde komme udaf det dermed, naar
engang Regnskabens Dag kommer. Det første er at trænge ind til Aanden og de
Grunderkjendelser, hvorpaa alt kommer an; ere vi grebne og opfyldte heraf og have
ganske tilegnet os disse, saa vil ogsaa Detaillets mangfoldige Erkjendelse komme
efterhaanden. Ogsaa i Studiet maae vi have for Øjne, hvad der staaer skrevet: Søger
først Guds Rige og hans Retfærdighed, saa skulle og alle disse Ting tillægges
Eder. Endog hvad den practiske Udøvelse angaaer, saa er vistnok nøjagtig Kundskab
om et mangfoldigt Detail fornøden, og det væsentligen fornøden dertil, men det
rette practiske er dog Besiddelsen af det, hvorfra dette Details sande Erkjendelse
gaaer ud. Hvor lidet gives der vel af den egentlige, dybere gaaende Philosophie,
som kan anvendes ligefrem i en ydre Virkekreds? Og dog er, næst Religionen, intet
mere practisk end Philosophien, der, idet den bringer det, hvorpaa hele Livet
hviler, til fuldkommen Erkjendelse og Klarhed, danner til overalt at lægge de rette
Principer, den rette Anskuelsesmaade, til Grund for sine Domme, og bestandigen at
have Hensyn til det, hvorpaa alt kommer an.

4.~Den practiske Udøvelse er kun ægte, naar Videnskab udgjør det bevægende Princip
i den. Vi maae overhovedet ikke glemme at \emph{tage Videnskaben med ind i Livet},
og fortsætte dens Studium under den hele practiske Virksomhed, der ogsaa
bestandigen, frembydende nyt Detail for Iagttagelse, frembyder nyt Stof for en
Tænkning, som ej blot umiddelbart kan blive frugtbringende for Praxis, men og, ved
at føre os tilbage til og i den hele Videnskab, kan bringe nyt Lys ind i dennes
Hele og derigjennem atter bringe nyt Lys ind i det practiske Liv. Kun saaledes
bevare vi Ideens Liv i os midt i Livet i den ydre Virkelighed, og kun saaledes
kommer den igjennem os til det Herredømme, hvortil den skal komme, til at beherske
Menneskelivet. Men saa er det og af Vigtighed, at vort ydre Kald, vor Virkekreds
virkelig staaer i nær Forbindelse med den Idee og Videnskab, vi især have studeret,
og at vi søge at bringe begge i den nærmeste muelige Berøring med hinanden. Det
sees da nu og let, hvor misligt det ogsaa i Videnskaben er at ville \emph{tjene
tvende Herrer}, ved at dyrke en anden Videnskab end den, hvis Dyrkelse vort ydre
Kald fordrer af os. Vel formaaer den menneskelige Aand ej sjeldent at omfatte
tvende ueensartede og usammenhørende Videnskabers Studium, men til vort Hovedfag
fordrer dog Samvittigheden, at vi gjøre den, til hvis Udøvelse vi ere satte, at vi
ikke for Exempel over Mineralogiens Studium forglemme at grunde over alt det, der
er vigtigt for en Sjælesørger, hvis Embede især fordrer et uafbrudt Studium hele
Livet igjennem, eller at vi ikke over Botanikens Studium lade vor medicinske Praxis
synke ned til det haandværksmæssige, eller at vi ikke som Skolemænd over Studiet af
en bestemt Green af Philologien forglemme Fornødenheden af den bestandige
Opmærksomhed paa vore Disciples Tarv i alle Henseender, noget, hvortil Iagttagelse
og grundende Eftertænkning uafbrudt er fornøden. Den, der ikke midt i sit Kalds
Arbejder føler sig nødt i sin Videnskab, ligner den, der ikke føler sig ret hjemme
i sit eget Huus, men har en Cirkel udenfor sit Huus, hvori han helst befinder sig,
en Stilling, der er unaturlig og mislig.

\begin{center}\large Slutning.\end{center}

Saa stort det er, den Erkjendende i sin Indsigt og Videnskab besidder, saa maae vi
dog ikke glemme, at det ej er dette, hvorpaa alt kommer an, det ene, som er
fornødent. Som da Paulus siger: at, om vi talte med Menneskenes og Englenes
Tungemaal og havde ikke Kjerlighed, saa vilde vi være et skingrende Malm og en
klingende Bjelde; og om vi havde Prophetie og vidste alle Hemmeligheder og al
Kundskab og havde al Tro, saa vi kunde flytte Bjerge, og havde ikke Kjerlighed, saa
vare vi intet. Kunde vel end Platon tale, som om den sande philosopherende, for sin
højere Erkjendelses Skyld, havde en højere Plads, et skjønnere Lod at vente efter
Døden, saa fører dog christeligt Sindelag en ganske anden Bedømmelse med sig. Kunne
end naturligt Anlæg og Talent og begunstigende ydre Omstændigheder give mangt et
Fortrin i Henseende til Livets højere Goder, saa staaer dog det højeste Gode alle
lige nær, og den Dannede har deri intet forud for den Udannede, saa meget han i
andre Henseender kan have forud for denne. Ikke, hvad vi vide, ikke, hvad vi
formaae med det dybere Blik at see og udtale, er det, hvorpaa alt kommer an, saa
lidet som det er det, vi ere blevne satte i Stand til at fremvirke og udrette.
Hvad vi ere, er det, hvorpaa alt kommer an, og det ej, hvad vi ere i Henseende til
den udadtrædende Bestaaen, hvori vi træde frem, men hvad vi ere i Henseende til det
Udspring, hvoraf vort Liv gaaer frem, om vi her have fuldkomment hengivet os i den
højestes Magt og have Gud overalt for Øje og ville, hvad Gud vil, hvad enten det
saa er, at vi skulle erkjende og skue meget eller lidet, at det, vi skulle udrette,
skal være stort eller ringe. Den, der formaaer at prædike Ordet, er derfor endnu ej
Frelsen og Saligheden nærmere. Det kommer an paa, om han nu og selv personligen
tager sig det Ord, det er givet ham at kunne forkynde, til Hjerte, og det netop
saaledes til Hjerte, som de skulle tage det til Hjerte, for hvilke det prædikes. Og
saaledes maa der gives et Punct i vort Indre, hvor vi føle, at alt, hvad vi i vor
Viden og ifølge vor Grunden besidde, endnu er intet, og i denne Mening ende vi med
intet hellere end med Christi Ord: Salige ere de, som ere fattige i Aanden.

% [BODY COMPLETE. Rettelser (errata) leaf — printed p.205 = PDF 227 — reproduced verbatim below.]

\bigskip
\begin{center}\large Rettelser.\end{center}

\noindent
Side 10 Linie 17 efter Sindsbevægelse sættes (,)\\
-- 13 -- 13 læs: i os.\\
-- 18 -- 14 efter viser sig udslettes (,)\\
-- 21 -- 11 istedetfor (,) sættes (.)\\
-- 21 -- 13 efter angaaer sættes (,)\\
-- 29 -- 3 fra neden udslettes begge ( )\\
-- 48 -- 2. I det læs: Idet.\\
-- 64 -- 5 læs: i det mindste.\\
-- 72 -- 16 den intelligente l. den fuldkomne intelligente.\\
-- 142 -- 6 fra neden: de l. der.\\
-- 145 -- 5 fra neden foran som sættes (,)\\
\hspace*{3em}-- 4 \hspace{2em}-- \hspace{2em}foran ogsaa (,)\\
-- 153 -- 7 fra neden udslettes (,) efter at.\\
-- 155 øverste Linie vis l. vist.\\
-- 161 -- 8 dennem l. denne.\\
-- 172 -- 14 efter Udspring udslettes (,) og sættes efter Herkomst,\\
-- 173 -- 5 fra neden sættes (,) efter vide.

\medskip
Endnu bemærkes, at der især i de første Ark ofte staaer maae istedetfor
maa, og et eller andet Sted igjen maa (i Singularis) istedetfor maae
(i Pluralis).

\begin{center}\rule{0.28\linewidth}{0.4pt}\end{center}

% [TRANSCRIPTION COMPLETE — front matter + §§1–22 + Slutning + Rettelser.]

\end{document}
